\chapter{Vectors and vector calculus}
\label{vol02:ch08}

\epigraph{Vector analysis is the language in which the local and the
global are joined: a quantity that flows through a point and a quantity
that accumulates inside a region are the same statement, written from
two ends of an integral.}{The working engineer's gloss on the
divergence theorem, after \textcite{text:strang2016}.}

\chapmeta{Half-life: foundational. Archetypes: balance (the divergence
theorem as the formal accounting equation between an interior source
and a surface flux); transport (the gradient as the local driver of
flux, the divergence as the local rate of net outflow per unit volume).
Exercise target: 35. Project track: household + analysis (hybrid).
Hazard class: none. Reader path default: standard, with sections
explicitly tagged. Named cases: none required; the failure section
treats orientation and sign-of-the-normal errors generically.}

\noindent
A vector is a quantity that carries both magnitude and direction;
vector calculus is the differential and integral apparatus by which
vector-valued functions of position are described, manipulated, and
balanced. The reader has met directed quantities since Volume I:
displacement, velocity, force, electric current, the gradient of a
temperature field. The mathematics of this chapter formalises a
practice the reader already has, and turns three common physical
operations, gradient, divergence, and curl, into a single bookkeeping
discipline. The chapter closes with the three integral theorems of
Green, Gauss, and Stokes, which bind the local rate of change of a
field to the integrated boundary behaviour of that field, and with
the project that earned the chapter its place in the volume: the
divergence theorem applied as a volume calculation against a sampled
surface, reconciled with the empirical water-displacement result the
reader produced in Volume~I, Chapter~7.

\section{Vectors as algebraic and geometric objects}\pathtag{core}

A \engterm{vector} in $\R^{n}$ is an ordered tuple
$\vect{v} = (v_{1}, v_{2}, \ldots, v_{n})$ together with the
operations of componentwise addition and scalar multiplication. We
write the same object as a column,
\[
\vect{v} = \begin{pmatrix} v_{1} \\ v_{2} \\ \vdots \\ v_{n} \end{pmatrix},
\]
and use the boldface convention $\vect{v}$ throughout this volume; the
older $\vec{v}$ notation appears in handwritten work, in some primary
sources, and in figure annotations. Both denote the same object. A
scalar, by contrast, is set in upright italic, $v$, and represents a
single real number.

The geometric content is older than the notation. A vector represents
a directed segment with a tail and a head; the segment's length is the
\engterm{magnitude} or \engterm{norm} $\norm{\vect{v}}$, and its
direction is the unit vector $\unitvec{v} = \vect{v} / \norm{\vect{v}}$.
The two pictures, algebraic tuple and geometric arrow, are connected
by a choice of \engterm{basis}: a set of $n$ linearly independent
vectors $\vect{e}_{1}, \ldots, \vect{e}_{n}$ such that any $\vect{v}$
is uniquely written as
$\vect{v} = v_{1}\vect{e}_{1} + \cdots + v_{n}\vect{e}_{n}$.
The Cartesian basis $\unitvec{x}, \unitvec{y}, \unitvec{z}$ in $\R^{3}$
is the working default, and the components $v_{i}$ in that basis are
the numbers an engineer reads off an instrument or writes on a drawing.

\subsection{Operations and identities}

Vector addition is componentwise, $(\vect{u} + \vect{v})_{i} = u_{i} +
v_{i}$, and obeys the parallelogram law geometrically: place
$\vect{u}$ and $\vect{v}$ tail to tail; their sum is the diagonal of
the parallelogram they span. Scalar multiplication
$\alpha \vect{v} = (\alpha v_{1}, \ldots, \alpha v_{n})$ stretches
the arrow by $\abs{\alpha}$ and reverses it when $\alpha < 0$.
Subtraction is addition of the negation: $\vect{u} - \vect{v} =
\vect{u} + (-1)\vect{v}$ points from the head of $\vect{v}$ to the
head of $\vect{u}$ when both share a tail.

The norm $\norm{\vect{v}} = \sqrt{v_{1}^{2} + \cdots + v_{n}^{2}}$
satisfies the three axioms a length must satisfy: it is non-negative,
it is zero only for the zero vector, and it obeys the triangle
inequality $\norm{\vect{u} + \vect{v}} \leq \norm{\vect{u}} +
\norm{\vect{v}}$ \cite{text:strang2016}. The triangle inequality is
the algebraic statement that the shortest path from a tail to a head
is a straight line; it becomes the engineer's tool for bounding
errors in vector-valued sums.

\begin{example}[Resultant of three coplanar forces]
Three cables pull on a fixed ring. Cable one applies
$\vect{F}_{1} = (200, 0, 0)\,\si{\newton}$, cable two applies
$\vect{F}_{2} = (-120, 160, 0)\,\si{\newton}$, and cable three applies
$\vect{F}_{3} = (0, -90, 0)\,\si{\newton}$. The resultant is the
componentwise sum,
\[
\vect{R} = \vect{F}_{1} + \vect{F}_{2} + \vect{F}_{3} =
(80, 70, 0)\,\si{\newton},
\]
with magnitude $\norm{\vect{R}} = \sqrt{80^{2} + 70^{2}} =
\sqrt{11300} \approx 106\,\si{\newton}$ and direction
$\arctan(70/80) \approx 41^{\circ}$ above the $+x$ axis. The ring is
not in equilibrium: a fourth cable supplying $-\vect{R} =
(-80, -70, 0)\,\si{\newton}$ would close the force polygon and bring
the net force to zero. The triangle inequality sets an upper bound on
the resultant before any arithmetic: $\norm{\vect{R}} \leq
\norm{\vect{F}_{1}} + \norm{\vect{F}_{2}} + \norm{\vect{F}_{3}} =
200 + 200 + 90 = 490\,\si{\newton}$. The actual resultant of
$106\,\si{\newton}$ sits well below the bound because the three
forces partly oppose one another; the gap between the bound and the
result is the geometric statement that the forces do not all pull in
the same direction.
\end{example}

\begin{example}[Closing error in a survey traverse]
A surveyor walks a closed loop, recording four displacement legs
$\vect{d}_{1} = (45.0, 0)\,\si{\meter}$, $\vect{d}_{2} = (0, 30.0)\,
\si{\meter}$, $\vect{d}_{3} = (-44.8, 0.2)\,\si{\meter}$, and
$\vect{d}_{4} = (-0.1, -30.3)\,\si{\meter}$. A perfectly closed loop
returns to its start, so the sum of the legs should be the zero
vector. The actual sum is the \engterm{closing error},
\[
\vect{e} = \sum_{i} \vect{d}_{i} = (0.1, -0.1)\,\si{\meter},
\]
with magnitude $\norm{\vect{e}} = \sqrt{0.1^{2} + 0.1^{2}} \approx
0.14\,\si{\meter}$. The total traverse length is roughly
$45 + 30 + 45 + 30 = 150\,\si{\meter}$, so the relative misclosure is
$0.14/150 \approx 1{:}1100$, within the $1{:}1000$ tolerance a
low-order traverse must meet. Vector addition turns four independent
field measurements into a single accountable error, and the norm of
that error vector is the surveyor's quality gate. A misclosure that
exceeds the tolerance signals a blunder (a transposed digit, a
mismeasured leg) before the parcel area is ever computed.
\end{example}

\subsection{Why vectors, not scalars}

A force on a beam, a velocity in a pipe, and an electric field at a
point all carry direction. A scalar description loses information the
engineer needs. A force of $100 \,\si{\newton}$ is not the same load
on a column as $100 \,\si{\newton}$ tangent to the column's axis; the
first is a compression, the second is a shear. A velocity of
$10 \,\si{\meter\per\second}$ in a duct is not the same air supply as
$10 \,\si{\meter\per\second}$ at $80^{\circ}$ to the duct axis; the
first delivers airflow, the second delivers swirl. The engineering
record is built from vectors precisely because the algebra of scalars
hides the directional information that determines whether a system
carries its load, delivers its mass flow, or aligns its sensor.

\begin{archetype}[Transport, vector form]
A transport quantity is one that flows from one region to another:
mass, momentum, charge, energy, signal. Its local description is a
\engterm{flux density}, a vector $\vect{J}$ whose magnitude is the
amount transported per unit area per unit time and whose direction is
the direction of net transport. The product $\vect{J} \cdot \unitvec{n}
\,dA$ is the rate of transport across an oriented surface element.
The chapter's later sections build this transport accounting into the
divergence theorem, but the vector $\vect{J}$ is the irreducible
object: a directed rate per unit area at a point.
\end{archetype}

\begin{estimation}
\textbf{Question.} A garden hose delivers water to a household lawn at
a measured volumetric flow rate of $8 \,\si{\liter\per\minute}$
through a nozzle of inner diameter $12 \,\si{\milli\meter}$. Estimate
the water's exit velocity vector at the nozzle, and the magnitude of
the velocity at a point on the lawn where the jet has spread to a
diameter of $40 \,\si{\milli\meter}$. State the dominant uncertainty
source and a factor-of-three confidence range.

\smallskip
\textit{Estimate, before reading on.}

\smallskip
The volumetric flow rate is $Q = 8 \,\si{\liter\per\minute} = 1.33
\times 10^{-4} \,\si{\meter\cubed\per\second}$. The nozzle area is
$A_{0} = \pi (0.006)^{2} = 1.13 \times 10^{-4} \,\si{\meter\squared}$.
The exit speed is $\norm{\vect{v}_{0}} = Q / A_{0} = 1.18 \,\si
{\meter\per\second}$, directed along the nozzle's axis. At the
$40 \,\si{\milli\meter}$ jet patch, the area has grown by a factor
$(40/12)^{2} \approx 11$; conservation of mass for an incompressible
jet gives a spread velocity $\norm{\vect{v}_{1}} \approx 0.11 \,\si
{\meter\per\second}$, again directed along the local jet axis.

The dominant uncertainty is the spread profile: the jet is not a
plug-flow column, and the boundary between jet and ambient air mixes
at a rate that grows with axial distance. A factor-of-three confidence
range on the local velocity at the patch, $0.04$ to $0.33 \,\si{\meter
\per\second}$, brackets the kind of estimate a domestic measurement
can support; a working engineer would either measure the velocity
directly with a small impeller anemometer or compute the spread from
a published self-similar jet profile. The estimate's value is the
algebraic structure: the same volumetric flow rate, divided by two
different cross-sectional areas, gives two velocities that differ by
the area ratio. The vector points along the jet axis at both points,
not at the same direction in space.
\end{estimation}

\subsection{The zero vector and linear independence}

Two vectors are \engterm{linearly dependent} if one is a scalar
multiple of the other; otherwise they are linearly independent. A set
of $k$ vectors in $\R^{n}$ is linearly independent if no non-trivial
linear combination $\alpha_{1}\vect{v}_{1} + \cdots + \alpha_{k}
\vect{v}_{k}$ equals the zero vector. A set of three vectors in
$\R^{3}$ is linearly independent if they do not all lie in a common
plane; geometrically, the parallelepiped they span has non-zero
volume. The reader meets this condition formally as the determinant
test in \cref{vol02:ch10}; for now we use the geometric criterion.

Linear independence governs whether a coordinate system is well-posed.
A pilot navigating with three compass bearings whose three directions
all lie in a horizontal plane has a well-posed two-dimensional
position fix, but no altitude information; the third bearing was not
linearly independent of the first two in $\R^{3}$. The error is not
in the magnitudes of the bearings but in their directions. Linear
independence is the discipline of checking, before measurement, that
the proposed reference frame can resolve the quantity that needs
resolving.

\begin{example}[Measurement directions versus a sufficient set]
A strain-gauge rosette measures normal strain along three in-plane
directions $\vect{a}_{1} = (1, 0)$, $\vect{a}_{2} = (\cos 45^{\circ},
\sin 45^{\circ})$, $\vect{a}_{3} = (0, 1)$. The plane strain has three
independent components, and each gauge returns one scalar projection
$\vect{a}_{i}\cdot\tens{\varepsilon}\,\vect{a}_{i}$, so three readings
are needed to invert for the three components. As vectors, only two of
the three directions are linearly independent, since any two
non-parallel plane vectors already span the plane and the third is a
combination of them. The rosette nonetheless needs all three gauges,
because the unknown is a tensor with three numbers and each gauge
supplies one equation. The distinction is between independence of the
measurement \emph{directions}, a geometric property of the vectors,
and sufficiency of the measurement \emph{set}, a count of how many
scalars the unknown carries. A rosette with all three gauges parallel
fails both tests and resolves nothing; a rosette with three distinct
in-plane directions passes the sufficiency test even though its third
direction is dependent on the first two. The engineer who conflates
the two tests either over-instruments (adding redundant directions
that add no information) or under-instruments (assuming two gauges
suffice for a three-component unknown).
\end{example}

\section{Dot and cross products}\pathtag{core}

Two products of vectors recur in working engineering: the
\engterm{dot product}, which produces a scalar, and the
\engterm{cross product}, which produces a vector. They are the
algebraic embodiments of two distinct geometric questions.

\subsection{The dot product}

The dot product of $\vect{u}$ and $\vect{v}$ in $\R^{n}$ is
\[
\vect{u} \cdot \vect{v} = u_{1} v_{1} + u_{2} v_{2} + \cdots +
u_{n} v_{n}.
\]
The geometric form is $\vect{u} \cdot \vect{v} = \norm{\vect{u}}
\norm{\vect{v}} \cos\theta$, where $\theta$ is the angle between
$\vect{u}$ and $\vect{v}$ when their tails are joined. The two forms
agree because the law of cosines applied to the triangle with sides
$\vect{u}$, $\vect{v}$, and $\vect{u} - \vect{v}$ rearranges to the
componentwise expression \cite{text:strang2016}.

The dot product has three working engineering uses:

\begin{itemize}[leftmargin=*]
  \item \engterm{Projection}. The component of $\vect{v}$ along the
    direction $\unitvec{u}$ is $\vect{v} \cdot \unitvec{u}$. This is
    how an engineer reads off the load along a beam axis when the
    applied force is at an angle, or how a controller extracts the
    along-track and cross-track components of a vehicle's velocity
    relative to a desired course.
  \item \engterm{Orthogonality test}. Two non-zero vectors are
    perpendicular if and only if $\vect{u} \cdot \vect{v} = 0$. The
    test is exact: any deviation from zero is the cosine of the
    actual angle, scaled by the magnitudes.
  \item \engterm{Work and rate of work}. The work done by a constant
    force $\vect{F}$ over a displacement $\vect{d}$ is
    $W = \vect{F} \cdot \vect{d}$; the instantaneous power
    $P = \vect{F} \cdot \vect{v}$. Volumes III and V take up these
    two uses in the contexts of mechanics and energy systems.
\end{itemize}

The dot product is commutative ($\vect{u} \cdot \vect{v} = \vect{v}
\cdot \vect{u}$), distributes over addition, and is bilinear in its
two arguments. It is positive on a non-zero vector dotted with itself:
$\vect{v} \cdot \vect{v} = \norm{\vect{v}}^{2} \geq 0$, with
equality only for $\vect{v} = \vect{0}$. The Cauchy-Schwarz
inequality, $\abs{\vect{u} \cdot \vect{v}} \leq \norm{\vect{u}}
\norm{\vect{v}}$, follows immediately and is the algebraic content of
$\abs{\cos\theta} \leq 1$. The Cauchy-Schwarz inequality is the
discipline behind every bound that says one vector quantity controls
another.

\begin{example}[Work done by an angled force]
A worker drags a crate $6.0\,\si{\meter}$ across a floor by a rope
inclined $30^{\circ}$ above the horizontal, pulling with a force of
$150\,\si{\newton}$. The displacement is $\vect{d} = (6.0, 0)\,
\si{\meter}$ and the force is $\vect{F} = (150\cos 30^{\circ},
150\sin 30^{\circ}) = (130, 75)\,\si{\newton}$. The work done is
\[
W = \vect{F}\cdot\vect{d} = 130 \times 6.0 + 75 \times 0 = 780\,
\si{\joule}.
\]
Only the horizontal component of the force does work, because the
displacement is horizontal; the vertical component
$75\,\si{\newton}$ contributes nothing because $\vect{d}$ has no
vertical part. A worker who reports $150 \times 6.0 = 900\,\si{\joule}$
has used the magnitude of the force instead of its component along the
displacement, and overstates the work by the factor
$1/\cos 30^{\circ} \approx 1.15$. The dot product enforces the
projection automatically; the scalar arithmetic does not.
\end{example}

\begin{example}[Decomposing a velocity into along-track and cross-track]
A vehicle moves with velocity $\vect{v} = (12, 5)\,\si{\meter\per
\second}$ relative to ground. Its desired course runs along the unit
vector $\unitvec{c} = (1, 0)$. The \engterm{along-track} speed is the
projection $\vect{v}\cdot\unitvec{c} = 12\,\si{\meter\per\second}$,
and the along-track velocity vector is $(\vect{v}\cdot\unitvec{c})
\unitvec{c} = (12, 0)\,\si{\meter\per\second}$. The
\engterm{cross-track} velocity is the remainder,
\[
\vect{v}_{\perp} = \vect{v} - (\vect{v}\cdot\unitvec{c})\unitvec{c}
= (12, 5) - (12, 0) = (0, 5)\,\si{\meter\per\second},
\]
a cross-track drift of $5\,\si{\meter\per\second}$ that the
controller must null. The decomposition $\vect{v} = \vect{v}_{\parallel}
+ \vect{v}_{\perp}$ into a part along $\unitvec{c}$ and a part
perpendicular to it is the building block of every track-following
controller; the dot product extracts the first part and subtraction
yields the second. Their magnitudes satisfy
$\norm{\vect{v}}^{2} = 12^{2} + 5^{2} = 169$, so $\norm{\vect{v}} =
13\,\si{\meter\per\second}$, and the Pythagorean split $12^{2} +
5^{2} = 13^{2}$ confirms the two parts are orthogonal.
\end{example}

\subsection{The cross product}

The cross product of $\vect{u}$ and $\vect{v}$ in $\R^{3}$ is the
unique vector $\vect{u} \times \vect{v}$ that is perpendicular to
both $\vect{u}$ and $\vect{v}$, has magnitude $\norm{\vect{u}}
\norm{\vect{v}} \sin\theta$, and points in the direction given by
the right-hand rule when $\vect{u}$ and $\vect{v}$ are not parallel.
In Cartesian components,
\[
\vect{u} \times \vect{v} =
\begin{pmatrix}
u_{2} v_{3} - u_{3} v_{2} \\
u_{3} v_{1} - u_{1} v_{3} \\
u_{1} v_{2} - u_{2} v_{1}
\end{pmatrix}.
\]
The componentwise expression is the determinant
\[
\vect{u} \times \vect{v} = \det
\begin{pmatrix}
\unitvec{x} & \unitvec{y} & \unitvec{z} \\
u_{1} & u_{2} & u_{3} \\
v_{1} & v_{2} & v_{3}
\end{pmatrix},
\]
when the determinant is expanded along the first row.

The cross product has three working engineering uses:

\begin{itemize}[leftmargin=*]
  \item \engterm{Area of a parallelogram}. The magnitude
    $\norm{\vect{u} \times \vect{v}} = \norm{\vect{u}}\norm{\vect{v}}
    \sin\theta$ is the area of the parallelogram with sides $\vect{u}$
    and $\vect{v}$. Triangle area is half of that. The cross product
    is the algebraic engine behind every triangulated surface area
    calculation in the second half of this chapter.
  \item \engterm{Torque and angular velocity}. The torque about a
    point produced by a force $\vect{F}$ applied at position
    $\vect{r}$ is $\vect{\tau} = \vect{r} \times \vect{F}$. The
    velocity of a point at position $\vect{r}$ on a rigid body
    rotating with angular velocity $\vect{\omega}$ is $\vect{v} =
    \vect{\omega} \times \vect{r}$. Volume III takes up rigid-body
    dynamics with these two formulas as the working tools.
  \item \engterm{Surface normal}. Given two non-parallel tangent
    vectors $\vect{t}_{u}$ and $\vect{t}_{v}$ to a surface at a
    point, the unit normal is $\unitvec{n} = (\vect{t}_{u} \times
    \vect{t}_{v}) / \norm{\vect{t}_{u} \times \vect{t}_{v}}$. The
    direction of the normal is fixed by the order of the cross
    product, which is the chapter's failure section.
\end{itemize}

\begin{figure}[h]
  \centering
  % Dot product (projection) and cross product (parallelogram area + normal).
% Two side-by-side panels.
\begin{tikzpicture}[>=Stealth, scale=1.0]

  % --- Panel A: dot product as projection ---
  \begin{scope}[xshift=0cm]
    \draw[->, thick] (0,0) -- (3.5,0) node[below right]{$\unitvec{u}$};
    \draw[->, thick, engblue] (0,0) -- (2.2,2.0) node[above right]{$\vect{v}$};
    \draw[dashed, enggray] (2.2,2.0) -- (2.2,0);
    \draw[very thick, engred] (0,0) -- (2.2,0);
    \node[engred, below, font=\footnotesize] at (1.1,-0.05) {$\vect{v}\cdot\unitvec{u}$};
    \draw (0.6,0) arc[start angle=0, end angle=42, radius=0.6];
    \node[font=\footnotesize] at (0.95,0.25) {$\theta$};
    \node[font=\footnotesize, align=center] at (1.75,-0.95) {Projection of $\vect{v}$\\onto $\unitvec{u}$};
  \end{scope}

  % --- Panel B: cross product as parallelogram + normal ---
  \begin{scope}[xshift=7cm, yshift=0cm]
    \coordinate (O) at (0,0);
    \coordinate (U) at (2.4,0.4);
    \coordinate (V) at (0.8,2.0);
    \coordinate (W) at ($(U)+(V)$);
    \fill[englime!30] (O) -- (U) -- (W) -- (V) -- cycle;
    \draw[->, thick] (O) -- (U) node[below right]{$\vect{u}$};
    \draw[->, thick] (O) -- (V) node[above left]{$\vect{v}$};
    \draw[engblue, dashed] (U) -- (W);
    \draw[engblue, dashed] (V) -- (W);
    \node[font=\footnotesize, align=center] at (1.6,1.0) {$\norm{\vect{u}\times\vect{v}}$\\$=$ area};
    % Normal vector out of plane indicated symbolically
    \draw[->, very thick, engred] (1.2,1.0) -- ++(0.0,-1.3) node[below, font=\footnotesize]{$\vect{u}\times\vect{v}$};
    \node[font=\footnotesize] at (0.2,-0.6) {right-hand rule};
  \end{scope}

\end{tikzpicture}

  \caption[Dot product as projection; cross product as oriented
  parallelogram.]{The dot product reads off the component of one vector
  along the direction of another (left). The cross product produces a
  vector whose magnitude is the area of the parallelogram spanned by
  $\vect{u}$ and $\vect{v}$ and whose direction is fixed by the
  right-hand rule (right). Source: authors' diagram.}
  \label{fig:vol02:ch08:dot-cross}
\end{figure}

\begin{example}[Torque from a force on a wrench]
A mechanic applies a force $\vect{F} = (0, -200, 0)\,\si{\newton}$
(downward) at the end of a wrench whose handle runs from the bolt
along $\vect{r} = (0.30, 0, 0)\,\si{\meter}$. The torque about the
bolt is
\[
\vect{\tau} = \vect{r}\times\vect{F} =
\begin{pmatrix}
(0)(0) - (0)(-200) \\
(0)(0) - (0.30)(0) \\
(0.30)(-200) - (0)(0)
\end{pmatrix}
= (0, 0, -60)\,\si{\newton\meter}.
\]
The torque points along $-\unitvec{z}$, a clockwise rotation seen from
above, which is what a downward pull on a rightward handle produces.
The magnitude $60\,\si{\newton\meter}$ equals $\norm{\vect{r}}
\norm{\vect{F}}\sin 90^{\circ} = 0.30\times 200$, because the force is
perpendicular to the handle. If the same force were applied along the
handle, $\vect{F} = (200, 0, 0)$, the cross product would vanish: a
pull along the handle's axis exerts no torque on the bolt. The
sine factor in $\norm{\vect{u}\times\vect{v}}$ is the geometric
reason a mechanic pulls perpendicular to the wrench, not along it.
\end{example}

\begin{example}[Unit normal to a triangle from two edges]
A triangular roof panel has vertices $\vect{p}_{1} = (0, 0, 0)$,
$\vect{p}_{2} = (4, 0, 1)$, $\vect{p}_{3} = (0, 3, 1)$, in metres.
Two edge vectors are $\vect{a} = \vect{p}_{2} - \vect{p}_{1} =
(4, 0, 1)$ and $\vect{b} = \vect{p}_{3} - \vect{p}_{1} = (0, 3, 1)$.
The cross product is
\[
\vect{a}\times\vect{b} =
\big((0)(1) - (1)(3),\ (1)(0) - (4)(1),\ (4)(3) - (0)(0)\big)
= (-3, -4, 12).
\]
Its magnitude is $\sqrt{9 + 16 + 144} = \sqrt{169} = 13$, so the
panel area is $\tfrac{1}{2}\times 13 = 6.5\,\si{\meter\squared}$, and
the unit normal is $\unitvec{n} = (-3, -4, 12)/13 \approx
(-0.23, -0.31, 0.92)$. The normal tilts mostly upward (large
$+z$ component), as a roof panel that rises one metre over its
extent should. Reversing the edge order to $\vect{b}\times\vect{a}$
would flip the normal to $(3, 4, -12)/13$, pointing into the roof
rather than out of it; the order of the cross product is the
chapter's recurring orientation decision.
\end{example}

\begin{figure}[h]
  \centering
  % Right-hand rule for the cross product u x v.
% Two coplanar vectors and the out-of-plane product, with a curved
% sweep arrow indicating the finger direction from u to v.
\begin{tikzpicture}[>=Stealth, scale=1.0]

  \coordinate (O) at (0,0);
  \coordinate (U) at (2.6,0);
  \coordinate (V) at (1.0,2.0);

  % the spanning vectors
  \draw[->, thick, engblue] (O) -- (U) node[below right]{$\vect{u}$};
  \draw[->, thick, engblue] (O) -- (V) node[above left]{$\vect{v}$};

  % angle theta between them
  \draw (0.7,0) arc[start angle=0, end angle=63, radius=0.7];
  \node[font=\footnotesize] at (1.05,0.45) {$\theta$};

  % sweep arrow from u toward v (finger curl direction)
  \draw[->, engamber, thick]
    (1.7,0.55) arc[start angle=18, end angle=58, radius=1.5];
  \node[engamber, font=\footnotesize] at (2.2,1.25) {sweep};

  % the product, out of plane (drawn vertically as a proxy for "up")
  \draw[->, very thick, engred] (O) -- (0,2.6)
    node[above, font=\footnotesize]{$\vect{u}\times\vect{v}$};
  \node[engred, font=\footnotesize, align=left] at (-0.95,1.6)
    {thumb};

  \node[font=\footnotesize, align=center] at (1.3,-0.85)
    {Right-hand rule: fingers sweep $\vect{u}\!\to\!\vect{v}$,\\
     thumb gives $\vect{u}\times\vect{v}$};

\end{tikzpicture}

  \caption[The right-hand rule for the cross product.]{The right-hand
  rule fixes the direction of $\vect{u}\times\vect{v}$: with the
  fingers of the right hand sweeping from $\vect{u}$ toward $\vect{v}$
  through the angle $\theta$, the thumb points along
  $\vect{u}\times\vect{v}$. Reversing the order to
  $\vect{v}\times\vect{u}$ reverses the thumb, hence the
  anti-commutativity $\vect{u}\times\vect{v} = -\vect{v}\times\vect{u}$.
  Source: authors' diagram.}
  \label{fig:vol02:ch08:right-hand-rule}
\end{figure}

The cross product is anti-commutative, $\vect{u} \times \vect{v} =
-\vect{v} \times \vect{u}$, distributes over addition, and is bilinear.
It is not associative: $\vect{u} \times (\vect{v} \times \vect{w})
\neq (\vect{u} \times \vect{v}) \times \vect{w}$ in general. The
\engterm{BAC-CAB rule}, $\vect{u} \times (\vect{v} \times \vect{w}) =
\vect{v}(\vect{u} \cdot \vect{w}) - \vect{w}(\vect{u} \cdot \vect{v})$,
expands the triple cross product into a sum of vectors weighted by
dot products, and is the working identity for unwinding a triple
cross product in a derivation.

\subsection{The scalar triple product and parallelepiped volume}

The \engterm{scalar triple product} $\vect{u} \cdot (\vect{v} \times
\vect{w})$ is a scalar whose absolute value equals the volume of the
parallelepiped spanned by $\vect{u}$, $\vect{v}$, $\vect{w}$. The
sign reports the handedness of the triple: positive when
$(\vect{u}, \vect{v}, \vect{w})$ is right-handed (in the same sense
as $\unitvec{x}, \unitvec{y}, \unitvec{z}$), negative when
left-handed. The scalar triple product equals the determinant of the
$3 \times 3$ matrix whose rows are $\vect{u}, \vect{v}, \vect{w}$:
\[
\vect{u} \cdot (\vect{v} \times \vect{w}) = \det
\begin{pmatrix}
u_{1} & u_{2} & u_{3} \\
v_{1} & v_{2} & v_{3} \\
w_{1} & w_{2} & w_{3}
\end{pmatrix}.
\]
The chapter's project rests on the scalar triple product: the
volume of a tetrahedron with one vertex at the origin and the
opposite face spanned by three vertices is one-sixth of the absolute
value of the scalar triple product of those three position vectors.

\begin{example}[Coplanarity test by scalar triple product]
An engineer must confirm that four survey points lie in a common
plane before fitting a flat slab to them. Take the points
$\vect{p}_{0} = (0, 0, 0)$, $\vect{p}_{1} = (2, 0, 1)$,
$\vect{p}_{2} = (0, 2, 1)$, $\vect{p}_{3} = (2, 2, 2)$, in metres.
Form three edge vectors from $\vect{p}_{0}$: $\vect{a} = (2, 0, 1)$,
$\vect{b} = (0, 2, 1)$, $\vect{c} = (2, 2, 2)$. The four points are
coplanar exactly when the parallelepiped spanned by the edges has
zero volume, that is, when $\vect{a}\cdot(\vect{b}\times\vect{c}) = 0$.
Compute $\vect{b}\times\vect{c} = ((2)(2) - (1)(2),\ (1)(2) - (0)(2),\
(0)(2) - (2)(2)) = (2, 2, -4)$, then
\[
\vect{a}\cdot(\vect{b}\times\vect{c}) = (2)(2) + (0)(2) + (1)(-4)
= 0.
\]
The triple product vanishes, so the four points are coplanar and the
flat slab fits exactly. Had the result been non-zero, its absolute
value would have measured the parallelepiped volume, and that volume
divided by the base parallelogram area would have given the
out-of-plane deviation the slab would have to absorb. The coplanarity
test is the scalar triple product set to zero; the same number, when
non-zero, reports the magnitude of the failure to be coplanar.
\end{example}

\subsection{Identities the working engineer keeps to hand}

A short list of vector identities recurs across Volumes III through
XII. Each can be derived from the componentwise definitions; the
reader is expected to know that they exist and to look them up when
needed, rather than to memorise them.

\begin{align*}
  \vect{u} \cdot (\vect{v} \times \vect{w}) &=
    \vect{v} \cdot (\vect{w} \times \vect{u}) =
    \vect{w} \cdot (\vect{u} \times \vect{v}), \\
  \vect{u} \times (\vect{v} \times \vect{w}) &=
    \vect{v}(\vect{u} \cdot \vect{w}) - \vect{w}(\vect{u} \cdot \vect{v}), \\
  (\vect{u} \times \vect{v}) \cdot (\vect{w} \times \vect{x}) &=
    (\vect{u} \cdot \vect{w})(\vect{v} \cdot \vect{x}) -
    (\vect{u} \cdot \vect{x})(\vect{v} \cdot \vect{w}), \\
  \norm{\vect{u} \times \vect{v}}^{2} &=
    \norm{\vect{u}}^{2}\norm{\vect{v}}^{2} -
    (\vect{u} \cdot \vect{v})^{2}.
\end{align*}

The last identity is sometimes called Lagrange's identity and is the
algebraic statement of $\sin^{2}\theta + \cos^{2}\theta = 1$.

\begin{mastery}
\textit{Index notation and the Levi-Civita symbol.} The vector
identities above are tedious to prove componentwise and routine to
prove in index notation, the working language of continuum mechanics
and tensor analysis. Write $\vect{u}\cdot\vect{v} = u_{i}v_{i}$ with
the Einstein convention that a repeated index is summed over
$1, 2, 3$. The cross product is $(\vect{u}\times\vect{v})_{i} =
\varepsilon_{ijk}u_{j}v_{k}$, where the \engterm{Levi-Civita symbol}
$\varepsilon_{ijk}$ is $+1$ for an even permutation of $(1, 2, 3)$,
$-1$ for an odd permutation, and $0$ when any two indices repeat. The
single identity
\[
\varepsilon_{ijk}\varepsilon_{ilm} = \delta_{jl}\delta_{km} -
\delta_{jm}\delta_{kl},
\]
with $\delta$ the Kronecker delta, collapses every double cross
product to a sum of dot products. The BAC-CAB rule, for instance,
follows in three lines: $(\vect{u}\times(\vect{v}\times\vect{w}))_{i} =
\varepsilon_{ijk}u_{j}\varepsilon_{klm}v_{l}w_{m} =
(\delta_{il}\delta_{jm} - \delta_{im}\delta_{jl})u_{j}v_{l}w_{m} =
v_{i}(u_{j}w_{j}) - w_{i}(u_{j}v_{j})$, which is $\vect{v}(\vect{u}
\cdot\vect{w}) - \vect{w}(\vect{u}\cdot\vect{v})$. A reader who learns
the $\varepsilon$-$\delta$ identity once derives the entire identity
list above without memorising any of it, and reads the index-notation
constitutive laws of Volumes III and V without translation. The cost
is one new symbol and the summation convention; the return is every
vector identity on demand.
\end{mastery}

\section{Coordinate systems and transformations}\pathtag{standard}

A coordinate system is a rule that assigns three numbers to each
point in $\R^{3}$. Three systems recur in working engineering, and the
discipline is to use each one where the geometry of the problem fits
the symmetry of the system.

\subsection{Cartesian, cylindrical, spherical}

\engterm{Cartesian coordinates} $(x, y, z)$ assign a triple of signed
distances along three mutually perpendicular axes; the basis
$\unitvec{x}, \unitvec{y}, \unitvec{z}$ is constant in space. A point
at position $\vect{r}$ has $\vect{r} = x\unitvec{x} + y\unitvec{y} +
z\unitvec{z}$.

\engterm{Cylindrical coordinates} $(\rho, \phi, z)$ assign a radial
distance from the $z$-axis, an azimuthal angle in the $xy$-plane
measured counterclockwise from the $x$-axis, and a height along the
$z$-axis. The conversion is $x = \rho\cos\phi$, $y = \rho\sin\phi$,
$z = z$. The basis vectors $\unitvec{\rho}, \unitvec{\phi},
\unitvec{z}$ vary with $\phi$: at azimuth $\phi$, $\unitvec{\rho} =
\cos\phi\,\unitvec{x} + \sin\phi\,\unitvec{y}$, $\unitvec{\phi} =
-\sin\phi\,\unitvec{x} + \cos\phi\,\unitvec{y}$, and $\unitvec{z}$
remains itself.

\engterm{Spherical coordinates} $(r, \theta, \phi)$ assign a radial
distance from the origin, a polar angle from the $+z$-axis, and an
azimuthal angle in the $xy$-plane. The conversion is $x = r\sin\theta
\cos\phi$, $y = r\sin\theta\sin\phi$, $z = r\cos\theta$. The basis
$\unitvec{r}, \unitvec{\theta}, \unitvec{\phi}$ varies with both
$\theta$ and $\phi$. The convention adopted here is the
physics-and-engineering convention: $\theta$ is the polar angle (zero
at the north pole, $\pi$ at the south pole), $\phi$ is the azimuthal
angle. The mathematics convention swaps the two; a reader who consults
a primary source in another field should check the convention before
trusting a formula.

The choice of system is governed by symmetry. A flow in a circular
pipe has cylindrical symmetry; cylindrical coordinates make
$\partial / \partial \phi$ vanish for many of its fields. A radiation
problem from a point source has spherical symmetry; spherical
coordinates make $\partial / \partial \theta$ and $\partial /
\partial \phi$ vanish for the radial field. A working engineer
wastes time using a coordinate system that does not match the
problem's symmetry; the coordinate-choice step is part of the
problem setup.

\begin{figure}[h]
  \centering
  % Curvilinear basis vectors: cylindrical (rho-hat, phi-hat) on the left,
% spherical (r-hat, theta-hat) on the right, showing that the basis turns
% with position, which is the origin of the scale factors.
\begin{tikzpicture}[>=Stealth, scale=1.0]

  % cylindrical panel
  \begin{scope}[xshift=-3.2cm]
    \draw[->, enggray] (-1.8,0) -- (1.8,0) node[right, font=\footnotesize]{$x$};
    \draw[->, enggray] (0,-1.8) -- (0,1.8) node[above, font=\footnotesize]{$y$};
    \draw[engblack, dashed] (0,0) circle (1.2);
    % point at azimuth 50 deg
    \coordinate (P) at (50:1.2);
    \fill[engblack] (P) circle (0.04);
    \draw[engblack, thin] (0,0) -- (P);
    % rho-hat (radial) and phi-hat (tangential) at P
    \draw[->, engred, thick] (P) -- ($(P)+(50:0.8)$)
      node[right, font=\footnotesize]{$\unitvec{\rho}$};
    \draw[->, engblue, thick] (P) -- ($(P)+(140:0.8)$)
      node[above, font=\footnotesize]{$\unitvec{\phi}$};
    \node[font=\footnotesize] at (0,-2.4) {cylindrical $(\rho,\phi)$};
  \end{scope}

  % spherical panel
  \begin{scope}[xshift=3.2cm]
    \draw[->, enggray] (0,-1.8) -- (0,1.8) node[above, font=\footnotesize]{$z$};
    \draw[->, enggray] (-1.8,0) -- (1.8,0) node[right, font=\footnotesize]{$x$};
    \draw[engblack, dashed] (0,0) circle (1.2);
    % point at polar angle 40 deg from z-axis
    \coordinate (Q) at (50:1.2); % 50 deg from x-axis = 40 deg from z
    \fill[engblack] (Q) circle (0.04);
    \draw[engblack, thin] (0,0) -- (Q);
    % r-hat (radial) and theta-hat (increasing polar angle, tangent)
    \draw[->, engred, thick] (Q) -- ($(Q)+(50:0.8)$)
      node[right, font=\footnotesize]{$\unitvec{r}$};
    \draw[->, engblue, thick] (Q) -- ($(Q)+(-40:0.8)$)
      node[below right, font=\footnotesize]{$\unitvec{\theta}$};
    \node[font=\footnotesize] at (0,-2.4) {spherical $(r,\theta)$};
  \end{scope}

\end{tikzpicture}

  \caption[Curvilinear basis vectors turn with position.]{The
  cylindrical basis $(\unitvec{\rho}, \unitvec{\phi})$ and the
  spherical basis $(\unitvec{r}, \unitvec{\theta})$ point in different
  directions at different points, unlike the fixed Cartesian basis. The
  radial unit vector aligns with the position direction; the angular
  unit vectors are tangent to the coordinate circles. The rotation of
  the basis with position is what introduces the $1/\rho$, $1/r$, and
  $1/(r\sin\theta)$ scale factors into the gradient, divergence, and
  curl in these systems. Source: authors' diagram.}
  \label{fig:vol02:ch08:curvilinear-basis}
\end{figure}

\subsection{Volume and area elements}

Volume integrals and surface integrals require the appropriate
infinitesimal element in the chosen system. The forms below cover
the three cases:

\begin{itemize}[leftmargin=*]
  \item Cartesian: $\dd V = \dd x \,\dd y \,\dd z$. A surface
    perpendicular to $\unitvec{z}$ has area element $\dd A = \dd x
    \,\dd y$.
  \item Cylindrical: $\dd V = \rho \,\dd\rho \,\dd\phi \,\dd z$. A
    cylindrical surface at fixed $\rho$ has $\dd A = \rho \,\dd\phi
    \,\dd z$; a flat top or bottom at fixed $z$ has $\dd A = \rho
    \,\dd\rho \,\dd\phi$.
  \item Spherical: $\dd V = r^{2} \sin\theta \,\dd r \,\dd\theta \,
    \dd\phi$. A spherical surface at fixed $r$ has $\dd A = r^{2}
    \sin\theta \,\dd\theta \,\dd\phi$.
\end{itemize}

The factors of $\rho$ and $r^{2}\sin\theta$ are the
\engterm{Jacobian determinants} of the coordinate change from
Cartesian. Volume~II Chapter~10 takes up the Jacobian as the formal
volume-distortion factor of any change of variables; here it is
enough to know the three forms above by inspection.

\begin{example}[Surface area of a spherical cap by the area element]
A pressure-vessel dome is a spherical cap of a sphere of radius
$R = 1.2\,\si{\meter}$, subtending a polar half-angle $\theta_{0} =
40^{\circ}$ from the vertical axis. The spherical area element at
fixed $r = R$ is $\dd A = R^{2}\sin\theta\,\dd\theta\,\dd\phi$. The
cap area is
\[
A = \int_{0}^{2\pi}\int_{0}^{\theta_{0}} R^{2}\sin\theta\,\dd\theta\,
\dd\phi = 2\pi R^{2}\big(1 - \cos\theta_{0}\big).
\]
With $R = 1.2\,\si{\meter}$ and $\cos 40^{\circ} = 0.766$,
$A = 2\pi(1.44)(1 - 0.766) = 2\pi(1.44)(0.234) \approx 2.12\,
\si{\meter\squared}$. The $\sin\theta$ factor in the element is what
makes the integral elementary: it is the circumference of the
latitude circle at polar angle $\theta$, scaled by $R$. An engineer
who forgets the $\sin\theta$ Jacobian and integrates $R^{2}\,\dd\theta
\,\dd\phi$ overestimates the cap area by treating every latitude band
as if it were at the equator. The same area enters the wall-material
estimate and the heat-loss calculation for the dome.
\end{example}

\subsection{Rotations as orthogonal transformations}

A \engterm{rotation} of $\R^{3}$ is a linear map that preserves
lengths and orientations. Its matrix representation $\mat{R}$ is a
$3 \times 3$ matrix satisfying $\mat{R}\transpose\mat{R} = \mat{I}$
and $\det\mat{R} = +1$. The first condition says columns are
orthonormal; the second excludes reflections, which would have
$\det\mat{R} = -1$. Volume~I Chapter~3 introduced rotations as
length-preserving transformations of position; here we generalise:
any vector quantity transforms under a rotation by left-multiplication
of its column representation by $\mat{R}$.

The composition of two rotations is again a rotation: if $\mat{R}_{1}$
and $\mat{R}_{2}$ are rotation matrices, $\mat{R}_{2} \mat{R}_{1}$ is
also a rotation matrix. The order matters, because matrix
multiplication does not commute in general; rotating first about
$\unitvec{x}$ and then about $\unitvec{y}$ produces a different
orientation from rotating first about $\unitvec{y}$ and then about
$\unitvec{x}$. The reader who has not seen this should rotate a
small block on a desk in the two orders and confirm by inspection.

A rotation about the unit axis $\unitvec{n}$ by angle $\alpha$ has
matrix
\[
\mat{R}(\unitvec{n}, \alpha) = \cos\alpha \,\mat{I} + (1 - \cos\alpha)
\,\unitvec{n}\unitvec{n}\transpose + \sin\alpha \,[\unitvec{n}]_{\times},
\]
where $[\unitvec{n}]_{\times}$ is the $3 \times 3$ matrix that, when
multiplied on a vector $\vect{v}$, returns $\unitvec{n} \times
\vect{v}$. This is Rodrigues's formula, and it is the working tool
behind every quaternion-based attitude system in aerospace, robotics,
and computer graphics.

\begin{example}[A vector under a $90^{\circ}$ rotation about $z$]
A robot's gripper points along $\vect{v} = (1, 0, 0)$ in the base
frame. The base rotates the gripper $90^{\circ}$ counterclockwise
about $\unitvec{z}$. The rotation matrix is
\[
\mat{R}_{z}(90^{\circ}) =
\begin{pmatrix}
\cos 90^{\circ} & -\sin 90^{\circ} & 0 \\
\sin 90^{\circ} & \cos 90^{\circ} & 0 \\
0 & 0 & 1
\end{pmatrix}
=
\begin{pmatrix}
0 & -1 & 0 \\
1 & 0 & 0 \\
0 & 0 & 1
\end{pmatrix}.
\]
Then $\mat{R}_{z}\vect{v} = (0, 1, 0)$: the gripper now points along
$+\unitvec{y}$, the expected quarter turn. The columns of
$\mat{R}_{z}$ are orthonormal, $\mat{R}_{z}\transpose\mat{R}_{z} =
\mat{I}$, and $\det\mat{R}_{z} = +1$, confirming a proper rotation
rather than a reflection. The inverse rotation is the transpose,
$\mat{R}_{z}^{-1} = \mat{R}_{z}\transpose = \mat{R}_{z}(-90^{\circ})$,
which the orthogonality condition guarantees for every rotation
matrix. An engineer who needs to undo a rotation transposes rather
than inverts, which is exact and free of the rounding that a general
matrix inverse introduces.
\end{example}

\begin{example}[Order of rotations does not commute]
Start with a book lying face up, spine to the left. Rotate it
$90^{\circ}$ about $\unitvec{x}$ (the spine axis pointing right),
then $90^{\circ}$ about $\unitvec{z}$ (the vertical axis). The
composite is $\mat{R}_{z}\mat{R}_{x}$. Reversing the order gives
$\mat{R}_{x}\mat{R}_{z}$, a different matrix. Numerically,
\[
\mat{R}_{z}\mat{R}_{x} =
\begin{pmatrix} 0 & 0 & 1 \\ 1 & 0 & 0 \\ 0 & 1 & 0 \end{pmatrix},
\qquad
\mat{R}_{x}\mat{R}_{z} =
\begin{pmatrix} 0 & -1 & 0 \\ 0 & 0 & -1 \\ 1 & 0 & 0 \end{pmatrix},
\]
which map the same starting vector to different orientations. The
non-commutativity is not a notational subtlety; it is the reason an
attitude-control system must record rotations in order and apply them
in that order, and the reason a sequence of Euler-angle commands sent
out of order leaves a spacecraft pointing the wrong way.
\end{example}

\subsection{Active and passive transformations}

A subtle distinction recurs across the engineering record. An
\engterm{active transformation} rotates the vector while the
coordinate frame stays still; the same physical body has new
components. A \engterm{passive transformation} rotates the coordinate
frame while the vector stays still; the same physical vector has new
components. The two operations are inverses of each other: an
active rotation by $\alpha$ corresponds to a passive rotation by
$-\alpha$. A working engineer who records that a frame was rotated
when in fact a body was rotated has signed the wrong calculation.
The discipline is to state, with each rotation matrix, whether the
matrix transforms vectors or transforms frames. Volumes III and VIII
take up the active-passive distinction in the contexts of rigid-body
dynamics and inertial navigation.

\section{Curves: parametrisation, arc length, curvature}\pathtag{standard}

A \engterm{curve} in $\R^{3}$ is the image of a continuous map
$\vect{r}: [a, b] \to \R^{3}$, $t \mapsto \vect{r}(t) =
(x(t), y(t), z(t))$. The variable $t$ is the \engterm{parameter}; the
choice of parametrisation is part of the engineer's modelling
decision. A curve has many parametrisations, all describing the same
geometric set of points.

\subsection{The tangent vector}

The \engterm{tangent vector} at parameter $t$ is the derivative
\[
\vect{r}\,'(t) = \dv{\vect{r}}{t} =
\left(\dv{x}{t}, \dv{y}{t}, \dv{z}{t}\right).
\]
Its direction is the direction of motion along the curve, and its
magnitude is the speed of motion under the chosen parametrisation. A
curve and its tangent picture together form the formal record of a
trajectory: the position $\vect{r}(t)$ at each $t$ and the velocity
$\vect{r}\,'(t)$ at each $t$.

\subsection{Arc length}

The \engterm{arc length} of the curve from $t = a$ to $t = b$ is
\[
s = \int_{a}^{b} \norm{\vect{r}\,'(t)} \,\dd t.
\]
The integrand is the speed under the parametrisation; the integral
accumulates infinitesimal arc-length contributions. The arc length
is a property of the curve, not the parametrisation: any
reparametrisation that traces the same set of points in the same
direction yields the same $s$.

A particularly clean parametrisation uses the arc length itself as
the parameter. Define $s(t) = \int_{a}^{t} \norm{\vect{r}\,'(\tau)}
\,\dd\tau$; if $\vect{r}\,'$ never vanishes, $s$ is a strictly
increasing function of $t$ and can be inverted to give $t(s)$. The
\engterm{arc-length parametrisation} $\vect{r}(s) = \vect{r}(t(s))$
has the property that $\norm{\vect{r}\,'(s)} = 1$ everywhere: the
tangent is a unit vector $\unitvec{T} = \vect{r}\,'(s)$. The
arc-length parametrisation is rarely available in closed form, but
its existence simplifies many derivations.

\subsection{Curvature}

The \engterm{curvature} $\kappa$ of a curve at a point is the
magnitude of the rate of change of the unit tangent with respect to
arc length:
\[
\kappa = \norm{\dv{\unitvec{T}}{s}}.
\]
For a circle of radius $R$, $\kappa = 1/R$ at every point. For a
straight line, $\kappa = 0$. For a curve given by an arbitrary
parametrisation, the working formula is
\[
\kappa = \frac{\norm{\vect{r}\,'(t) \times \vect{r}\,''(t)}}
              {\norm{\vect{r}\,'(t)}^{3}},
\]
which can be derived by the chain rule from the arc-length form.

The \engterm{principal normal} $\unitvec{N}$ at a point of non-zero
curvature is the unit vector in the direction of $\dd\unitvec{T}/
\dd s$; the \engterm{binormal} $\unitvec{B} = \unitvec{T} \times
\unitvec{N}$. The triple $(\unitvec{T}, \unitvec{N}, \unitvec{B})$
is the \engterm{Frenet-Serret frame}, an orthonormal basis attached
to the curve at each point. The frame's derivatives along the curve
satisfy the Frenet-Serret formulas, which the reader will meet in
the Volume~III treatment of rigid-body kinematics.

\begin{mastery}
The Frenet-Serret formulas, in arc-length parametrisation, are
\[
\dv{\unitvec{T}}{s} = \kappa \unitvec{N}, \qquad
\dv{\unitvec{N}}{s} = -\kappa\unitvec{T} + \tau\unitvec{B}, \qquad
\dv{\unitvec{B}}{s} = -\tau\unitvec{N},
\]
where $\kappa$ is the curvature and $\tau$ is the
\engterm{torsion}, the rate at which the curve twists out of the
osculating plane spanned by $\unitvec{T}$ and $\unitvec{N}$. A planar
curve has $\tau = 0$ everywhere; a helix has constant $\kappa$ and
constant non-zero $\tau$. The three equations are the working tools
of differential geometry of curves \cite{text:v2ch08-do-carmo-1976},
and they recur in the analysis of aircraft trajectories, of charged
particles in magnetic fields, and of helical springs and screw
threads. The reader who has not yet met them in a working context can
defer the derivation to the Volume~III treatment and treat the three
identities above as a reference triple.
\end{mastery}

\begin{example}[Arc length of a parabolic cable]
A cable hangs in a shallow parabola $y = \tfrac{1}{200}x^{2}$ for
$x \in [-50, 50]\,\si{\meter}$, a span of $100\,\si{\meter}$ with a
sag of $50^{2}/200 = 12.5\,\si{\meter}$ at the ends relative to the
midpoint. Parametrise by $x$ itself: $\vect{r}(x) = (x,
x^{2}/200)$, so $\vect{r}\,'(x) = (1, x/100)$ and
$\norm{\vect{r}\,'(x)} = \sqrt{1 + x^{2}/10^{4}}$. The cable length is
\[
L = \int_{-50}^{50}\sqrt{1 + \frac{x^{2}}{10^{4}}}\,\dd x.
\]
The integral has the closed form $L = \big[\tfrac{x}{2}\sqrt{1 +
x^{2}/10^{4}} + 50\,\operatorname{arcsinh}(x/100)\big]_{-50}^{50}$.
Evaluating, $\sqrt{1 + 0.25} = 1.118$ and
$\operatorname{arcsinh}(0.5) = 0.4812$, giving
$L = 2\big(25\times 1.118 + 50\times 0.4812\big) = 2(27.95 + 24.06)
\approx 104.0\,\si{\meter}$. The cable is $4.0\,\si{\meter}$ longer
than its $100\,\si{\meter}$ horizontal span, a $4\,\%$ excess that
the procurement order must include. A flat-cable estimate that
ignores the sag underbuys the cable, a recurring error in spans where
the deflection looks small on a drawing.
\end{example}

Curvature controls the lateral acceleration of any object travelling
along a curved path: an object moving at speed $v$ along a curve of
radius of curvature $R = 1/\kappa$ experiences a centripetal
acceleration $a_{c} = v^{2}/R = \kappa v^{2}$. A working highway
designer uses this relation to set the maximum design speed for a
banked curve given a target lateral acceleration; a working
roller-coaster designer uses the same relation in inverse to choose
the curvature for a target sensation. The relation links the
geometric curvature of the curve to the dynamic experience of
following it.

\begin{example}[Minimum curve radius for a highway design speed]
A highway curve is to be driven at $v = 30\,\si{\meter\per\second}$
(about $108\,\si{\kilo\meter\per\hour}$) with a lateral acceleration
no greater than $a_{c} = 0.3g \approx 2.9\,\si{\meter\per\second
\squared}$, the comfort limit for unbanked travel. The required
radius of curvature follows from $a_{c} = v^{2}/R$:
\[
R = \frac{v^{2}}{a_{c}} = \frac{30^{2}}{2.9} = \frac{900}{2.9}
\approx 310\,\si{\meter}.
\]
The curvature must satisfy $\kappa = 1/R \leq 3.2\times 10^{-3}\,
\si{\per\meter}$ everywhere along the curve. A curve laid out with a
tighter radius forces either a lower design speed or super-elevation
(banking) to supply part of the centripetal force. The geometry sets
the kinematics: the curvature the surveyor stakes out is the
acceleration the driver feels, scaled by $v^{2}$. The factor $v^{2}$
is why a curve safe at city speed becomes uncomfortable at highway
speed without any change to the road.
\end{example}

\subsection{Line integrals}

Two line integrals recur. The \engterm{scalar line integral} of a
function $f$ along a curve is
\[
\int_{C} f \,\dd s = \int_{a}^{b} f(\vect{r}(t)) \norm{\vect{r}\,'(t)}
\,\dd t,
\]
and reports a quantity-times-arc-length, such as a mass per unit
length integrated along a wire to give total mass.

\begin{example}[Mass of a wire with variable density]
A wire bent into the helix $\vect{r}(t) = (\cos t, \sin t, t)$ for
$t \in [0, 2\pi]$ has linear density $\lambda(\vect{r}) = 1 + z =
1 + t$ kilograms per metre. The speed is $\norm{\vect{r}\,'(t)} =
\sqrt{2}$, constant along the helix. The total mass is the scalar
line integral
\[
m = \int_{C}\lambda\,\dd s = \int_{0}^{2\pi}(1 + t)\sqrt{2}\,\dd t =
\sqrt{2}\Big[t + \tfrac{1}{2}t^{2}\Big]_{0}^{2\pi} =
\sqrt{2}\big(2\pi + 2\pi^{2}\big) \approx 36.8\,\si{\kilogram}.
\]
The $\sqrt{2}$ arc-length factor is the bridge between the parameter
$t$ and physical length: integrating $(1 + t)\,\dd t$ alone would
return the mass of a wire of unit speed, undercounting by the helix's
$\sqrt{2}$ stretch per unit parameter. The scalar line integral
carries the geometry through the density automatically; an engineer
who integrates the density against $\dd t$ instead of $\dd s$
mismeasures the mass by the speed factor.
\end{example} The \engterm{vector
line integral} of a vector field $\vect{F}$ along a curve is
\[
\int_{C} \vect{F} \cdot \dd\vect{r} = \int_{a}^{b} \vect{F}(\vect{r}(t))
\cdot \vect{r}\,'(t) \,\dd t,
\]
and reports a quantity-times-displacement-component-along-curve, such
as the work done by a force field on a particle moving along $C$.
The vector line integral is the prototype of the circulation
integral that appears in Stokes's theorem.

\begin{example}[Work along a curved path]
A force field $\vect{F}(x, y) = (y, -x)$ acts on a bead that moves
along the quarter circle $\vect{r}(t) = (\cos t, \sin t)$ for
$t \in [0, \pi/2]$, from $(1, 0)$ to $(0, 1)$. The tangent is
$\vect{r}\,'(t) = (-\sin t, \cos t)$, and along the path
$\vect{F}(\vect{r}(t)) = (\sin t, -\cos t)$. The integrand is
\[
\vect{F}\cdot\vect{r}\,' = (\sin t)(-\sin t) + (-\cos t)(\cos t)
= -\sin^{2} t - \cos^{2} t = -1,
\]
so the work is $\int_{0}^{\pi/2}(-1)\,\dd t = -\pi/2 \approx -1.57$.
The field opposes the motion at every point (the dot product is $-1$
throughout), so the work is negative. This field is the rotational
field of \cref{fig:vol02:ch08:div-curl}; its circulation will return
in the Green's-theorem treatment below, where the same line integral
around a closed loop measures the enclosed curl. A bead pushed
against this field around a full circle requires work equal to the
loop's enclosed area times the field's curl.
\end{example}

\begin{example}[Path dependence in a non-conservative field]
The same field $\vect{F}(x, y) = (y, -x)$ moves a particle from
$(1, 0)$ to $(0, 1)$ by a second route: the straight segment
$\vect{r}(t) = (1 - t, t)$ for $t \in [0, 1]$. Now $\vect{r}\,'(t) =
(-1, 1)$ and $\vect{F}(\vect{r}(t)) = (t, -(1 - t)) = (t, t - 1)$, so
the integrand is $\vect{F}\cdot\vect{r}\,' = (t)(-1) + (t - 1)(1) =
-t + t - 1 = -1$, giving work $-1$. The arc route gave $-\pi/2
\approx -1.57$; the straight route gives $-1$. The two endpoints are
the same, yet the work differs. The field is not conservative: its
line integral depends on the path. The gap $-\pi/2 - (-1) = 1 -
\pi/2 \approx -0.57$ is exactly the circulation around the closed
loop formed by going out along the arc and back along the segment,
and Green's theorem will tie that circulation to the curl
$\partial(-x)/\partial x - \partial(y)/\partial y = -1 - 1 = -2$
integrated over the enclosed region. Path dependence is the
operational signature of a non-zero curl.
\end{example}

\section{Surfaces: parametrisation, normals, area}\pathtag{standard}

A \engterm{parametrised surface} in $\R^{3}$ is the image of a
continuous map $\vect{r}: D \to \R^{3}$, $(u, v) \mapsto
\vect{r}(u, v)$, where $D$ is a region in the $(u, v)$ parameter
plane. The domain $D$ is rectangular for simple surfaces and more
elaborate for closed surfaces such as the sphere, where the
parametrisation either covers the surface in patches or accepts a
controlled singularity at the poles.

\subsection{Tangent vectors and the surface normal}

At each interior point of $D$, the partial derivatives
$\vect{r}_{u} = \partial\vect{r}/\partial u$ and $\vect{r}_{v} =
\partial\vect{r}/\partial v$ are tangent vectors to the surface. When
they are linearly independent, the surface is called \engterm{regular}
at that point, and the normal vector is
\[
\vect{N} = \vect{r}_{u} \times \vect{r}_{v}.
\]
The unit normal is $\unitvec{n} = \vect{N} / \norm{\vect{N}}$. The
order of the cross product fixes the direction of the normal: swapping
$u$ and $v$, or swapping the order of the cross product, reverses
the direction. The reader who has chosen a parametrisation has
implicitly chosen an orientation; the chapter's failure section is
about what goes wrong when the implicit choice and the intended
choice disagree.

\subsection{Surface area}

The infinitesimal area element is $\dd A = \norm{\vect{r}_{u} \times
\vect{r}_{v}} \,\dd u \,\dd v$, and the surface area of the
parametrised surface is
\[
A = \iint_{D} \norm{\vect{r}_{u} \times \vect{r}_{v}} \,\dd u \,\dd v.
\]
The cross-product magnitude is the area of the parallelogram spanned
by the two tangent vectors at each point; the integral accumulates
those infinitesimal parallelograms across the surface. The formula
generalises the cross-product-as-parallelogram-area identity from
section~8.2 to a continuous family of parallelograms.

\begin{example}[Lateral surface area of a cone]
A conical funnel has base radius $a = 0.10\,\si{\meter}$ and height
$h = 0.25\,\si{\meter}$. Parametrise the lateral surface by
$\vect{r}(u, v) = \big(\tfrac{a u}{h}\cos v, \tfrac{a u}{h}\sin v,
u\big)$ for $u \in [0, h]$ (height) and $v \in [0, 2\pi]$. The
tangents are $\vect{r}_{u} = (\tfrac{a}{h}\cos v, \tfrac{a}{h}\sin v,
1)$ and $\vect{r}_{v} = (-\tfrac{a u}{h}\sin v, \tfrac{a u}{h}\cos v,
0)$, and their cross product has magnitude $\norm{\vect{r}_{u}\times
\vect{r}_{v}} = \tfrac{a u}{h}\sqrt{1 + a^{2}/h^{2}}$. The area is
\[
A = \int_{0}^{2\pi}\!\int_{0}^{h}\frac{a u}{h}\sqrt{1 + \frac{a^{2}}
{h^{2}}}\,\dd u\,\dd v = 2\pi\,\frac{a}{h}\sqrt{1 + \frac{a^{2}}{h^{2}}}
\cdot\frac{h^{2}}{2} = \pi a\sqrt{a^{2} + h^{2}}.
\]
This recovers the textbook lateral-area formula $\pi a \ell$ with slant
height $\ell = \sqrt{a^{2} + h^{2}}$. Numerically, $\ell = \sqrt{0.01
+ 0.0625} = 0.269\,\si{\meter}$ and $A = \pi(0.10)(0.269) \approx
0.085\,\si{\meter\squared}$. The parametrisation and the surface-area
integral reproduce a memorised formula, which is the cross-check an
engineer runs to trust the integral machinery before applying it to a
shape with no closed-form area.
\end{example}

\subsection{Surface integrals}

The \engterm{scalar surface integral} of a function $f$ over a
parametrised surface $S$ is
\[
\iint_{S} f \,\dd A = \iint_{D} f(\vect{r}(u, v))
\norm{\vect{r}_{u} \times \vect{r}_{v}} \,\dd u \,\dd v,
\]
and reports a quantity-times-area, such as the surface mass of a
membrane with variable density. The \engterm{flux integral} of a
vector field $\vect{F}$ through an oriented surface $S$ is
\[
\iint_{S} \vect{F} \cdot \unitvec{n} \,\dd A =
\iint_{D} \vect{F}(\vect{r}(u, v)) \cdot
(\vect{r}_{u} \times \vect{r}_{v}) \,\dd u \,\dd v,
\]
and reports the rate of transport of $\vect{F}$ across $S$ in the
direction of $\unitvec{n}$. The flux integral is the working object
of the divergence theorem and the chapter's project.

\begin{example}[Flux of a uniform field through a tilted plate]
A uniform field $\vect{F} = (0, 0, 5)$, in units of the transported
quantity per square metre per second, crosses a flat rectangular
plate of area $2\,\si{\meter\squared}$ whose unit normal is
$\unitvec{n} = (0, \sin 30^{\circ}, \cos 30^{\circ}) = (0, 0.5,
0.866)$. The flux is
\[
\iint_{S}\vect{F}\cdot\unitvec{n}\,\dd A = (\vect{F}\cdot\unitvec{n})A
= (5\times 0.866)(2) = 8.66
\]
of the quantity per second. Only the component of $\vect{F}$ along
$\unitvec{n}$ crosses the plate; the plate tilted $30^{\circ}$ from
horizontal intercepts $\cos 30^{\circ} = 0.866$ of the flux a
horizontal plate would. This is the rule behind a solar panel's
output falling with the cosine of its tilt away from the incoming
radiation, and behind a sail's thrust depending on its angle to the
wind. A plate edge-on to the field ($\unitvec{n}\perp\vect{F}$)
catches zero flux, as the dot product reports.
\end{example}

\begin{example}[Flux through a curved surface patch]
The field $\vect{F} = (x, y, 0)$ crosses the curved side of a
cylinder of radius $a = 2$, height $h = 3$, axis along $\unitvec{z}$.
Parametrise the side as $\vect{r}(\phi, z) = (a\cos\phi, a\sin\phi,
z)$ for $\phi \in [0, 2\pi]$, $z \in [0, 3]$. The tangents are
$\vect{r}_{\phi} = (-a\sin\phi, a\cos\phi, 0)$ and $\vect{r}_{z} =
(0, 0, 1)$, so $\vect{r}_{\phi}\times\vect{r}_{z} = (a\cos\phi,
a\sin\phi, 0)$, which points radially outward, as required. On the
surface $\vect{F} = (a\cos\phi, a\sin\phi, 0)$, so the integrand is
$\vect{F}\cdot(\vect{r}_{\phi}\times\vect{r}_{z}) = a^{2}\cos^{2}\phi
+ a^{2}\sin^{2}\phi = a^{2}$. The flux is
\[
\iint_{S}\vect{F}\cdot\unitvec{n}\,\dd A = \int_{0}^{3}\int_{0}^{2\pi}
a^{2}\,\dd\phi\,\dd z = a^{2}(2\pi)(3) = 4\cdot 6\pi = 24\pi
\approx 75.4.
\]
A cross-check by the divergence theorem confirms this below:
$\divg\vect{F} = 2$, and the cylinder volume is $\pi a^{2}h = 12\pi$,
so $\iiint 2\,\dd V = 24\pi$; the two flat caps carry zero flux of
this field because $\vect{F}$ has no $z$ component. The surface
integral and the volume integral agree, which is the divergence
theorem at work.
\end{example}

\begin{example}[Flux through a hemispherical cap, two ways]
The field $\vect{F} = (0, 0, z)$ crosses the upper unit hemisphere
$S$, the surface $x^{2} + y^{2} + z^{2} = 1$ with $z \geq 0$, oriented
with an upward (outward) normal. The direct surface integral uses the
spherical parametrisation $\vect{r}(\theta, \phi) = (\sin\theta\cos\phi,
\sin\theta\sin\phi, \cos\theta)$ for $\theta \in [0, \pi/2]$, $\phi \in
[0, 2\pi]$. The outward unit normal on the unit sphere is the radial
direction itself, $\unitvec{n} = \vect{r}$, so on the surface
$\vect{F}\cdot\unitvec{n} = z\cdot z = \cos^{2}\theta$, and the area
element is $\dd A = \sin\theta\,\dd\theta\,\dd\phi$. The flux is
\[
\iint_{S}\vect{F}\cdot\unitvec{n}\,\dd A = \int_{0}^{2\pi}\!\int_{0}
^{\pi/2}\cos^{2}\theta\sin\theta\,\dd\theta\,\dd\phi = 2\pi\left[
-\frac{\cos^{3}\theta}{3}\right]_{0}^{\pi/2} = 2\pi\cdot\frac{1}{3} =
\frac{2\pi}{3}.
\]
The cross-check closes the hemisphere with its flat disc cap at
$z = 0$ and applies the divergence theorem to the solid half-ball.
The divergence is $\divg\vect{F} = \partial z/\partial z = 1$, so the
volume integral is just the volume of the half-ball,
$\iiint_{V}1\,\dd V = \tfrac{1}{2}\cdot\tfrac{4}{3}\pi = \tfrac{2}{3}
\pi$. The closed surface is the hemisphere plus the disc; on the disc
$z = 0$ the outward normal is $-\unitvec{z}$ and $\vect{F} = (0, 0, 0)$,
so the disc carries zero flux. The hemisphere therefore carries the
entire $2\pi/3$, matching the direct integral. The two unrelated
computations, one a surface integral in spherical angles and the other
a volume of a half-ball, agree to the digit, which is the most direct
confirmation a reader can run that the curved-surface flux integral
and the divergence theorem describe the same quantity.
\end{example}

\begin{example}[Surface area of a torus from its parametrisation]
The torus is the surface a flux integral meets in cooling coils,
toroidal pressure vessels, and the magnetic-confinement geometry of
Volume~IX, and it has no elementary axis of symmetry that collapses
the area to a memorised formula. Parametrise the torus of tube radius
$a$ and centreline radius $R > a$ by
\[
\vect{r}(u, v) = \big((R + a\cos u)\cos v,\ (R + a\cos u)\sin v,\
a\sin u\big),
\]
with $u \in [0, 2\pi]$ the angle around the tube cross-section and
$v \in [0, 2\pi]$ the angle around the centreline. The tangents are
\[
\vect{r}_{u} = \big(-a\sin u\cos v,\ -a\sin u\sin v,\ a\cos u\big),
\qquad
\vect{r}_{v} = \big(-(R + a\cos u)\sin v,\ (R + a\cos u)\cos v,\ 0\big).
\]
Their cross product has the magnitude
$\norm{\vect{r}_{u}\times\vect{r}_{v}} = a(R + a\cos u)$, the product
of the two radii of curvature swept by the parameter pair: $a$ around
the tube and $R + a\cos u$ around the centreline at tube-angle $u$.
The factor $R + a\cos u$ is larger on the outer rim ($u = 0$) than on
the inner rim ($u = \pi$), which is the geometric statement that the
outer half of the tube carries more surface than the inner half. The
area is
\[
A = \int_{0}^{2\pi}\!\int_{0}^{2\pi} a(R + a\cos u)\,\dd u\,\dd v =
2\pi\,a\int_{0}^{2\pi}(R + a\cos u)\,\dd u = 2\pi a\,(2\pi R) =
4\pi^{2}aR.
\]
The $\cos u$ term integrates to zero around the tube: the extra area
on the outer rim exactly compensates the deficit on the inner rim, so
the total area is that of a cylinder of length $2\pi R$ (the
centreline) and circumference $2\pi a$ (the tube). For a cooling coil
with $R = 0.15\,\si{\meter}$ and $a = 0.02\,\si{\meter}$, the area is
$A = 4\pi^{2}(0.02)(0.15) \approx 0.118\,\si{\meter\squared}$, the
surface available for convective heat exchange. The parametrisation
turned a shape with no elementary area formula into a single integral,
and the integral's $\cos u$ cancellation is Pappus's centroid theorem
recovered from the cross-product area element.
\end{example}

\subsection{Closed surfaces and outward normals}

A \engterm{closed surface} bounds a three-dimensional region. The
working convention is that the unit normal $\unitvec{n}$ on a closed
surface points outward, away from the bounded region. With this
convention, $\oint_{\partial V} \vect{F} \cdot \unitvec{n} \,\dd A$ is
the net outward flux of $\vect{F}$ through the boundary of the
volume $V$. A reader using a parametrisation whose induced normal
points inward must either reverse the parametrisation or change the
sign of the result. The discipline is to verify the orientation
explicitly, and the failure section returns to the operational test
that catches the error.

\section{Gradient, divergence, curl}\pathtag{core}

Three differential operators on fields recur in working engineering.
Each takes a field defined on a region of $\R^{3}$ and returns
another field. Each has an interpretation as a local rate of
something the engineer is tracking.

\subsection{The gradient}

A \engterm{scalar field} $\Phi(\vect{r})$ assigns a number to each
point. The \engterm{gradient} $\grad\Phi$ is the vector field whose
components in Cartesian coordinates are
\[
\grad\Phi = \left(\pdv{\Phi}{x}, \pdv{\Phi}{y}, \pdv{\Phi}{z}\right).
\]
The gradient at a point has two interpretations.

First, it is the direction of steepest increase of $\Phi$ at the
point. A reader walking on a heated plate with temperature field
$T(\vect{r})$ feels the temperature increase fastest in the direction
$\grad T$ and decrease fastest in the direction $-\grad T$. The
magnitude $\norm{\grad T}$ is the rate of increase per unit distance
in that direction.

Second, the gradient is the local linear approximation to the field's
variation: for small $\dd\vect{r}$,
\[
\dd\Phi \approx \grad\Phi \cdot \dd\vect{r}.
\]
This is the multivariable analogue of the scalar derivative
relation $\dd f \approx f'(x) \dd x$. The right-hand side is the
\engterm{directional derivative} of $\Phi$ in the direction
$\dd\vect{r}$.

A vector field that is the gradient of some scalar field is called
\engterm{conservative}. Conservative fields have the property that
their line integral between two endpoints does not depend on the
path connecting them: $\int_{C} \grad\Phi \cdot \dd\vect{r} =
\Phi(\vect{r}_{\text{end}}) - \Phi(\vect{r}_{\text{start}})$. Volume~III
takes up conservative force fields and the closely related notion of
potential energy.

\begin{example}[Steepest ascent on a temperature plate]
A heated plate has temperature field $T(x, y) = 100 - 2x^{2} - y^{2}$,
in degrees Celsius with $x, y$ in metres. The gradient is
$\grad T = (-4x, -2y)$. At the point $(2, 3)$, $\grad T = (-8, -6)$,
with magnitude $\norm{\grad T} = \sqrt{64 + 36} = 10\,\si{\kelvin\per
\meter}$, pointing toward $(-8, -6)/10 = (-0.8, -0.6)$, that is, back
toward the hot centre. A walker at $(2, 3)$ moving in the direction
$\unitvec{u} = (1, 0)$ feels the temperature change at the rate
$\grad T\cdot\unitvec{u} = -8\,\si{\kelvin\per\meter}$: cooling,
because moving in $+x$ moves away from the hot centre. Moving along
$\unitvec{u} = (-0.8, -0.6)$ gives $\grad T\cdot\unitvec{u} =
(-8)(-0.8) + (-6)(-0.6) = 6.4 + 3.6 = 10\,\si{\kelvin\per\meter}$,
the maximum rate, because that direction aligns with the gradient.
Moving perpendicular to the gradient, along the level set, gives
zero: the temperature is constant along a contour. The directional
derivative is the gradient projected onto the direction of travel,
and it is largest when the travel direction is the gradient itself.
\end{example}

\begin{example}[Recovering a potential from a conservative field]
A field $\vect{F}(x, y) = (2xy + 3, x^{2} - 1)$ is claimed to be
conservative. The test is whether a potential $\Phi$ exists with
$\grad\Phi = \vect{F}$. Integrate the first component in $x$:
$\Phi = \int(2xy + 3)\,\dd x = x^{2}y + 3x + g(y)$, where $g(y)$ is an
unknown function of $y$ alone. Differentiate in $y$ and match the
second component: $\partial\Phi/\partial y = x^{2} + g'(y)$ must
equal $x^{2} - 1$, so $g'(y) = -1$ and $g(y) = -y + C$. The potential
is $\Phi(x, y) = x^{2}y + 3x - y + C$. The reconstruction succeeds, so
the field is conservative, and the work it does from $(0, 0)$ to
$(1, 2)$ is $\Phi(1, 2) - \Phi(0, 0) = (2 + 3 - 2) - 0 = 3$,
independent of path. Had the matching step produced a $g'$ that
depended on $x$, no potential would exist and the field would not be
conservative; the failed match is the operational signal that the
curl is non-zero.
\end{example}

\begin{example}[Potential reconstruction in three dimensions]
The three-dimensional case adds one matching step but follows the same
procedure. A field $\vect{F}(x, y, z) = (yz + 2x,\ xz + 1,\ xy - 3)$
is claimed conservative. The curl test runs first and cheapest:
$(\curl\vect{F})_{x} = \partial(xy - 3)/\partial y - \partial(xz + 1)/
\partial z = x - x = 0$, and the other two components vanish by the
same cancellation, so a potential exists. To find it, integrate the
first component in $x$:
\[
\Phi = \int(yz + 2x)\,\dd x = xyz + x^{2} + g(y, z),
\]
where $g$ is an unknown function of $y$ and $z$. Differentiate in $y$
and match the second component: $\partial\Phi/\partial y = xz +
\partial g/\partial y$ must equal $xz + 1$, so $\partial g/\partial y
= 1$ and $g(y, z) = y + h(z)$, with $h$ a function of $z$ alone.
Differentiate the running potential $\Phi = xyz + x^{2} + y + h(z)$ in
$z$ and match the third component: $\partial\Phi/\partial z = xy +
h'(z)$ must equal $xy - 3$, so $h'(z) = -3$ and $h(z) = -3z + C$. The
potential is
\[
\Phi(x, y, z) = xyz + x^{2} + y - 3z + C.
\]
The work the field does from $(0, 0, 0)$ to $(1, 1, 1)$ is
$\Phi(1, 1, 1) - \Phi(0, 0, 0) = (1 + 1 + 1 - 3) - 0 = 0$, independent
of path. The three-step ladder, integrate in $x$, match in $y$, match
in $z$, recovers the potential of any three-dimensional conservative
field, and each matching step either succeeds (leaving an integration
constant that depends on fewer variables) or fails (producing a
required derivative that depends on a variable already integrated out,
which signals a non-zero curl the test should already have caught).
\end{example}

\begin{mastery}
\textit{The Helmholtz decomposition.} A sufficiently smooth vector
field that decays fast enough at infinity splits uniquely into a
curl-free part and a divergence-free part,
\[
\vect{F} = -\grad\Phi + \curl\vect{A},
\]
the \engterm{Helmholtz decomposition}. The first term, the gradient of
a scalar potential $\Phi$, carries all of the field's divergence:
$\divg\vect{F} = -\nabla^{2}\Phi$, since $\divg(\curl\vect{A}) = 0$.
The second term, the curl of a vector potential $\vect{A}$, carries
all of the field's curl: $\curl\vect{F} = \curl(\curl\vect{A})$, since
$\curl(\grad\Phi) = \vect{0}$. The two scalar identities
$\divg(\curl\,\cdot) = 0$ and $\curl(\grad\,\cdot) = \vect{0}$ from the
previous subsection are exactly what make the split clean: each
potential feeds one operator and is invisible to the other. The
decomposition is the structural reason a fluid velocity field separates
into a potential flow (irrotational, $\vect{v} = \grad\Phi$) plus a
rotational remainder, and the reason electromagnetism is written with a
scalar potential $V$ and a vector potential $\vect{A}$: the electric
and magnetic fields are the gradient and curl parts of a single
four-potential. A reader who recovers a scalar potential by the ladder
above has computed the curl-free part of a field whose curl already
vanished; the Helmholtz decomposition is what to reach for when the
curl does not vanish and the field still needs to be split into pieces
that the gradient and the curl can each handle. The construction of
$\Phi$ and $\vect{A}$ from $\vect{F}$ is by Poisson-equation solves,
$\nabla^{2}\Phi = -\divg\vect{F}$ and $\nabla^{2}\vect{A} =
-\curl\vect{F}$, which Volume~II Chapter~12 takes up as the working
field-equation solvers.
\end{mastery}

\begin{example}[Splitting a field into gradient and curl parts]
A two-dimensional flow field $\vect{F}(x, y) = (x - y,\ x + y)$ has
both a divergence and a curl, so neither a pure potential nor a pure
stream function describes it; the Helmholtz decomposition splits it
into the two pieces each tool handles. The divergence is
$\divg\vect{F} = \partial(x - y)/\partial x + \partial(x + y)/
\partial y = 1 + 1 = 2$, and the curl is $(\curl\vect{F})_{z} =
\partial(x + y)/\partial x - \partial(x - y)/\partial y = 1 - (-1) =
2$. The decomposition seeks $\vect{F} = \grad\Phi + \curl(\psi
\unitvec{z})$, where $\Phi$ carries the divergence and the scalar
$\psi$ (the stream function) carries the curl. The scalar potential
solves $\nabla^{2}\Phi = \divg\vect{F} = 2$, satisfied by
$\Phi = \tfrac{1}{2}(x^{2} + y^{2})$, whose gradient is the radial
part $\grad\Phi = (x, y)$. The stream function solves $\nabla^{2}\psi
= -(\curl\vect{F})_{z} = -2$, satisfied by $\psi = -\tfrac{1}{2}
(x^{2} + y^{2})$, whose curl part $\curl(\psi\unitvec{z}) =
(\partial\psi/\partial y,\ -\partial\psi/\partial x) = (-y, x)$ is the
rotational field. The two pieces sum to
\[
\grad\Phi + \curl(\psi\unitvec{z}) = (x, y) + (-y, x) = (x - y,\
x + y) = \vect{F},
\]
recovering the original field. The radial part $(x, y)$ is an
irrotational source flow (all divergence, no curl); the part
$(-y, x)$ is an incompressible rotation (all curl, no divergence). A
reader who decomposes a measured flow this way separates the part that
a potential-flow solver handles from the part that carries the
vorticity, which is the working preprocessing step behind every
vortex-panel and source-panel aerodynamic method
\cite{text:anderson-aerodynamics}.
\end{example}

\begin{figure}[h]
  \centering
  % Gradient field of a scalar potential.
% Contour lines of Phi(x,y) = x^2 + 2 y^2 with arrows for grad Phi.
\begin{tikzpicture}[>=Stealth, scale=1.0]
  \begin{axis}[
    width=10cm, height=7cm,
    xlabel={$x$}, ylabel={$y$},
    xmin=-2.2, xmax=2.2, ymin=-1.6, ymax=1.6,
    axis lines=center,
    label style={font=\footnotesize},
    tick label style={font=\footnotesize},
    samples=21,
  ]
    % Contours of Phi = x^2 + 2 y^2.
    \addplot[engblue, very thin, domain=0:360, samples=80] ({0.6*cos(x)}, {0.6/sqrt(2)*sin(x)});
    \addplot[engblue, very thin, domain=0:360, samples=80] ({1.0*cos(x)}, {1.0/sqrt(2)*sin(x)});
    \addplot[engblue, very thin, domain=0:360, samples=80] ({1.4*cos(x)}, {1.4/sqrt(2)*sin(x)});
    \addplot[engblue, very thin, domain=0:360, samples=80] ({1.8*cos(x)}, {1.8/sqrt(2)*sin(x)});

    % Gradient arrows: grad Phi = (2x, 4y); arrows scaled.
    \addplot[
      black,
      quiver={u=2*x/12, v=4*y/12, scale arrows=1.0, every arrow/.append style={-Stealth, thin}},
    ] coordinates {
      (-1.5,-1) (-1,-1) (-0.5,-1) (0.0,-1) (0.5,-1) (1,-1) (1.5,-1)
      (-1.5,-0.5) (-1,-0.5) (-0.5,-0.5) (0.0,-0.5) (0.5,-0.5) (1,-0.5) (1.5,-0.5)
      (-1.5,0.5) (-1,0.5) (-0.5,0.5) (0.0,0.5) (0.5,0.5) (1,0.5) (1.5,0.5)
      (-1.5,1) (-1,1) (-0.5,1) (0.0,1) (0.5,1) (1,1) (1.5,1)
    };

    \node[font=\footnotesize, anchor=west, engblue] at (axis cs:1.5,0.05) {level sets of $\Phi$};
    \node[font=\footnotesize, anchor=west] at (axis cs:0.7,1.25) {$\grad\Phi$ (perpendicular to level set)};
  \end{axis}
\end{tikzpicture}

  \caption[Gradient field of $\Phi(x, y) = x^{2} + 2 y^{2}$.]{Level
  sets of the scalar field $\Phi = x^{2} + 2 y^{2}$ are nested
  ellipses; the gradient $\grad\Phi = (2x, 4y)$ is perpendicular to
  the level set at every point and grows in magnitude with distance
  from the origin. Source: authors' computation.}
  \label{fig:vol02:ch08:gradient-field}
\end{figure}

\begin{archetype}[Transport, gradient form]
A transport quantity moves down a gradient. Heat flows from hot to
cold, with flux density $\vect{q} = -k \grad T$ for thermal
conductivity $k$. Mass flows down a concentration gradient, with
flux density $\vect{J} = -D\grad c$ for diffusivity $D$. Charge in
a conductor flows down a voltage gradient, with current density
$\vect{j} = -\sigma_{e}\grad V$ for electrical conductivity
$\sigma_{e}$. The minus sign reports that flux opposes the gradient:
the quantity moves in the direction of decreasing $\Phi$, which is
the direction $-\grad\Phi$. Volumes III, IV, and V take up the
specific cases; the archetype is the same.
\end{archetype}

\begin{example}[Heat flux down a temperature gradient]
A wall has a steady temperature field $T(x) = 60 - 200x$, in degrees
Celsius with $x$ in metres measured from the warm inner face, across a
$0.20\,\si{\meter}$ wall whose outer face is at $60 - 200(0.20) =
20\,\degC$. The gradient is $\grad T = (-200, 0, 0)\,\si{\kelvin\per
\meter}$, pointing toward the warm side (decreasing $x$). For a wall
of thermal conductivity $k = 0.8\,\si{\watt\per\meter\per\kelvin}$,
Fourier's law gives the heat flux
\[
\vect{q} = -k\grad T = -0.8\times(-200, 0, 0) = (160, 0, 0)\,\si{\watt
\per\meter\squared},
\]
pointing in $+x$, from the warm face toward the cold face. The minus
sign in Fourier's law is what turns the inward-pointing gradient into
an outward-pointing flux: heat flows down the temperature gradient,
not up it. Through a wall of area $12\,\si{\meter\squared}$ the total
conducted power is $160\times 12 = 1.9\,\si{\kilo\watt}$. The same
gradient-driven-transport structure governs the concentration-driven
diffusion of a contaminant and the voltage-driven current in a
resistor; the archetype is one equation with three sets of units.
\end{example}

\subsection{The divergence}

A \engterm{vector field} $\vect{F}(\vect{r})$ assigns a vector to each
point. The \engterm{divergence} $\divg\vect{F}$ is the scalar field
whose value in Cartesian coordinates is
\[
\divg\vect{F} = \pdv{F_{x}}{x} + \pdv{F_{y}}{y} + \pdv{F_{z}}{z}.
\]
The divergence at a point has the interpretation of net outward flux
per unit volume. Imagine a small box centred at the point; the rate
at which $\vect{F}$ flows out of the box, divided by the volume of
the box, is $\divg\vect{F}$ in the limit as the box shrinks to the
point. A positive divergence is a local source; a negative divergence
is a local sink; zero divergence is a point that conserves the
quantity carried by $\vect{F}$.

\begin{example}[Divergence of a radial outflow]
A point sprinkler at the origin drives a radial velocity field
$\vect{v} = (kx, ky, kz)$ with $k = 0.5\,\si{\per\second}$, a field
that grows linearly with distance from the source. Its divergence is
$\divg\vect{v} = k + k + k = 3k = 1.5\,\si{\per\second}$, a uniform
positive value everywhere. The interpretation is direct: every
infinitesimal box loses more volume out of its outer faces than it
gains through its inner faces, at the rate $1.5$ per second per unit
volume. A box of volume $V_{0}$ centred anywhere expands at the rate
$\dd V/\dd t = (\divg\vect{v})V_{0} = 1.5\,V_{0}$ per second. This is
the divergence as a local volume-creation rate; an incompressible
flow would forbid it, so this field describes a compressible
expansion, such as a gas escaping into vacuum near the source.
\end{example}

For an incompressible fluid, the velocity field $\vect{v}$ satisfies
$\divg\vect{v} = 0$ everywhere: no point is a source or sink of
volume. For a temperature field with a local heating element, the
heat flux $\vect{q}$ has $\divg\vect{q} = -\rho c \,\partial T /
\partial t$ at points away from the element and $\divg\vect{q}$ equal
to the local volumetric heat source at the element. The reader will
meet this balance as the heat equation in Volume~II Chapter~12.

\subsection{The material derivative}

A field measured by a fixed sensor and a field measured by a sensor
riding along with the flow report two different rates of change. The
distinction is the most common source of sign and magnitude error in
the transport equations of Volumes III and V, and it is a pure
application of the gradient. Let $\Phi(\vect{r}, t)$ be a scalar field,
say the temperature of a moving fluid, and let a fluid parcel follow
the path $\vect{r}(t)$ with velocity $\vect{v} = \dd\vect{r}/\dd t$.
The rate of change of $\Phi$ that the parcel experiences is the total
derivative along its path, computed by the chain rule,
\[
\frac{\mathrm{D}\Phi}{\mathrm{D}t} = \pdv{\Phi}{t} +
\dv{\vect{r}}{t}\cdot\grad\Phi = \pdv{\Phi}{t} +
(\vect{v}\cdot\grad)\Phi.
\]
The operator $\mathrm{D}/\mathrm{D}t = \partial/\partial t +
\vect{v}\cdot\grad$ is the \engterm{material derivative} (also the
convective or substantial derivative). Its first term is the
\engterm{local rate}, the change a fixed sensor reads because the
field itself is changing in time. Its second term is the
\engterm{advective rate}, the change a moving parcel reads because it
has been carried into a region where the field has a different value.
A parcel can experience a changing temperature in a steady field
($\partial\Phi/\partial t = 0$) purely by moving down a spatial
gradient, and a fixed sensor can read a changing temperature in a
field with zero spatial gradient purely by the field's own time
evolution. The two rates coincide only when the parcel is at rest or
the field is uniform.

Applied to the velocity field itself, the material derivative gives
the acceleration of a fluid parcel,
\[
\vect{a} = \frac{\mathrm{D}\vect{v}}{\mathrm{D}t} = \pdv{\vect{v}}{t}
+ (\vect{v}\cdot\grad)\vect{v},
\]
the left-hand side of the Navier-Stokes momentum equation that
Volume~V develops \cite{text:white-fluid-mechanics}. The advective
term $(\vect{v}\cdot\grad)\vect{v}$ is nonlinear in $\vect{v}$, which
is the algebraic root of the difficulty of fluid mechanics: a parcel
that accelerates because it is carried into a faster region feeds back
on the velocity field that carries it.

\begin{example}[Heating a parcel in a steady flow]
Air flows steadily through a heated duct with velocity $\vect{v} =
(2, 0, 0)\,\si{\meter\per\second}$ and a steady temperature field that
rises along the duct, $T(x) = 20 + 50x$ in degrees Celsius with $x$ in
metres. The field is steady, so the local rate $\partial T/\partial t
= 0$: a thermometer bolted to the duct wall at a fixed $x$ reads a
constant temperature. A parcel of air carried along by the flow reads
something different. The material derivative is
\[
\frac{\mathrm{D}T}{\mathrm{D}t} = \pdv{T}{t} + (\vect{v}\cdot\grad)T
= 0 + (2)\pdv{T}{x} = 2\times 50 = 100\,\si{\kelvin\per\second}.
\]
The parcel heats at $100\,\si{\kelvin\per\second}$ even though every
fixed point in the duct holds a constant temperature, because the
parcel is swept down the temperature gradient at $2\,\si{\meter\per
\second}$ into ever-hotter regions. An engineer sizing the heater
needs the parcel's rate, not the wall sensor's rate: the energy the
flow carries away is set by how fast each parcel heats, which the
material derivative reports and the fixed-sensor reading hides.
\end{example}

\subsection{The curl}

The \engterm{curl} $\curl\vect{F}$ of a vector field is the vector
field whose Cartesian components are
\[
\curl\vect{F} =
\begin{pmatrix}
\partial F_{z}/\partial y - \partial F_{y}/\partial z \\
\partial F_{x}/\partial z - \partial F_{z}/\partial x \\
\partial F_{y}/\partial x - \partial F_{x}/\partial y
\end{pmatrix},
\]
or as the symbolic determinant
\[
\curl\vect{F} = \det
\begin{pmatrix}
\unitvec{x} & \unitvec{y} & \unitvec{z} \\
\partial/\partial x & \partial/\partial y & \partial/\partial z \\
F_{x} & F_{y} & F_{z}
\end{pmatrix}.
\]
The curl at a point has the interpretation of twice the local
angular velocity of an infinitesimal paddle wheel placed at the
point and free to rotate about an axis aligned with $\curl\vect{F}$.
A vector field with $\curl\vect{F} = \vect{0}$ everywhere is called
\engterm{irrotational}; a fluid flow whose velocity is irrotational
has no local rotation. A non-zero curl reports local circulation.

\begin{example}[Curl of a shear flow]
A river flows faster at the surface than near the bed, modelled by a
horizontal shear $\vect{v} = (cy, 0, 0)$ with $c = 2\,\si{\per
\second}$ and $y$ the height above the bed. The only non-zero curl
component is
\[
(\curl\vect{v})_{z} = \pdv{v_{y}}{x} - \pdv{v_{x}}{y} = 0 - c = -2\,
\si{\per\second}.
\]
A paddle wheel floating in this flow spins clockwise (about
$-\unitvec{z}$) at angular velocity $\tfrac{1}{2}\abs{(\curl
\vect{v})_{z}} = 1\,\si{\radian\per\second}$, because the faster water
at its top pushes harder than the slower water at its bottom. The
flow has straight, parallel streamlines and yet a non-zero curl: a
field can have curl without any visible swirl in its streamlines.
Curl measures local differential rotation, which the shear supplies
even when every fluid element travels in a straight line.
\end{example}

\begin{figure}[h]
  \centering
  % Conservative vs non-conservative field.
% Left: gradient field of Phi = x^2 + y^2 (arrows radial, level sets
% circular, arrows cross contours at right angles).
% Right: rotational field (-y, x) with closed-loop circulation.
\begin{tikzpicture}[>=Stealth, scale=1.0]

  % --- Panel A: conservative (gradient) field ---
  \begin{scope}[xshift=0cm]
    % level sets (circles)
    \draw[enggray] (0,0) circle (0.6);
    \draw[enggray] (0,0) circle (1.1);
    \draw[enggray] (0,0) circle (1.5);
    % gradient arrows: radial, grow with radius
    \foreach \ang in {0,45,90,135,180,225,270,315} {
      \draw[->, engblue, thick]
        ({0.7*cos(\ang)},{0.7*sin(\ang)}) -- ({1.45*cos(\ang)},{1.45*sin(\ang)});
    }
    \node[font=\footnotesize, align=center] at (0,-2.0)
      {Conservative: $\vect{F}=\grad\Phi$,\\path-independent work};
  \end{scope}

  % --- Panel B: non-conservative (rotational) field ---
  \begin{scope}[xshift=5.5cm]
    \foreach \ang in {0,45,90,135,180,225,270,315} {
      \pgfmathsetmacro{\xs}{1.1*cos(\ang)}
      \pgfmathsetmacro{\ys}{1.1*sin(\ang)}
      \pgfmathsetmacro{\xt}{\xs + 0.55*(-sin(\ang))}
      \pgfmathsetmacro{\yt}{\ys + 0.55*cos(\ang)}
      \draw[->, engred, thick] (\xs,\ys) -- (\xt,\yt);
    }
    % a closed loop with net circulation
    \draw[->, engamber, thick] (0,0) circle (0.55);
    \node[engamber, font=\footnotesize] at (0,0) {$\oint\neq 0$};
    \node[font=\footnotesize, align=center] at (0,-2.0)
      {Non-conservative: $\curl\vect{F}\neq\vect{0}$,\\path-dependent work};
  \end{scope}

\end{tikzpicture}

  \caption[Conservative versus non-conservative fields.]{A
  conservative field (left) is the gradient of a scalar potential; its
  arrows cross the level sets at right angles and its line integral
  between two points is path-independent. A non-conservative field
  (right) has closed-loop circulation, so the work done from one point
  to another depends on the path taken. The curl distinguishes the
  two: zero curl on a simply connected region certifies a conservative
  field. Source: authors' diagram.}
  \label{fig:vol02:ch08:conservative}
\end{figure}

\begin{figure}[h]
  \centering
  % Divergence (radial source) vs curl (rotational) sketch.
% Two side-by-side panels with arrows on a small region.
\begin{tikzpicture}[>=Stealth, scale=1.0]

  % --- Panel A: divergence (source field F = r-hat / r) ---
  \begin{scope}[xshift=0cm]
    \draw[enggray] (0,0) circle (1.4);
    \foreach \ang in {0,30,60,90,120,150,180,210,240,270,300,330} {
      \draw[->, engblue, thick]
        ({0.4*cos(\ang)},{0.4*sin(\ang)}) -- ({1.3*cos(\ang)},{1.3*sin(\ang)});
    }
    \fill (0,0) circle (0.05);
    \node[font=\footnotesize, below=1.55cm] at (0,0) {$\divg\vect{F} > 0$ (source)};
  \end{scope}

  % --- Panel B: curl (rotational field F = (-y, x)) ---
  \begin{scope}[xshift=5.5cm]
    \draw[enggray] (0,0) circle (1.4);
    \foreach \ang in {0,30,60,90,120,150,180,210,240,270,300,330} {
      \pgfmathsetmacro{\xs}{1.0*cos(\ang)}
      \pgfmathsetmacro{\ys}{1.0*sin(\ang)}
      \pgfmathsetmacro{\xt}{\xs + 0.5*(-sin(\ang))}
      \pgfmathsetmacro{\yt}{\ys + 0.5*cos(\ang)}
      \draw[->, engred, thick] (\xs,\ys) -- (\xt,\yt);
    }
    \draw[->, very thick] (0,0) -- (0,0.5) node[above, font=\footnotesize]{$\curl\vect{F}$};
    \node[font=\footnotesize, below=1.55cm] at (0,0) {$\curl\vect{F} \neq \vect{0}$ (rotation)};
  \end{scope}

\end{tikzpicture}

  \caption[Divergence and curl in two model fields.]{A radial source
  field $\vect{F} = \hat{\vect{r}}/r$ has positive divergence at the
  origin (left); a rotational field $\vect{F} = (-y, x)$ has uniform
  non-zero curl along $+\unitvec{z}$ (right). A reader who imagines a
  small paddle wheel placed in either field sees the divergence as a
  net outflow and the curl as a rotation rate. Source: authors'
  diagram.}
  \label{fig:vol02:ch08:div-curl}
\end{figure}

\subsection{Two identities the engineer keeps to hand}

Two identities link the three operators. The first,
\[
\divg(\curl\vect{F}) = 0
\quad\text{for any sufficiently smooth } \vect{F},
\]
says the curl of a vector field has zero divergence. The second,
\[
\curl(\grad\Phi) = \vect{0}
\quad\text{for any sufficiently smooth } \Phi,
\]
says the gradient of a scalar field has zero curl. The two identities
are the backbone of Helmholtz decomposition, which factors a
sufficiently nice vector field into a curl-free part (a gradient)
and a divergence-free part (a curl). The Volume~II Chapter~12
treatment of the Laplace and Poisson equations rests on these two
identities.

The second identity gives the working test for conservativeness on a
simply connected region: a field is the gradient of some potential if
and only if its curl vanishes there. The field $\vect{F} = (2xy + 3,
x^{2} - 1, 0)$ recovered as a gradient above has curl $(\curl
\vect{F})_{z} = \partial(x^{2} - 1)/\partial x - \partial(2xy + 3)/
\partial y = 2x - 2x = 0$, confirming the potential exists before any
integration is attempted. The field $(y, -x, 0)$ of the line-integral
examples has $(\curl\vect{F})_{z} = -1 - 1 = -2 \neq 0$, certifying
that no potential exists and that its work integral must depend on the
path, as the two routes computed above confirmed. The curl test
spares the engineer the integration when the field is not
conservative.

\begin{mastery}
\textit{The Laplacian and the second-order operators.} Composing the
first-order operators produces the second-order operators the
engineer meets in the field equations of Volumes IV, V, and IX. The
\engterm{Laplacian} of a scalar field is the divergence of its
gradient,
\[
\nabla^{2}\Phi = \divg(\grad\Phi) =
\frac{\partial^{2}\Phi}{\partial x^{2}} +
\frac{\partial^{2}\Phi}{\partial y^{2}} +
\frac{\partial^{2}\Phi}{\partial z^{2}},
\]
and measures how much the field's value at a point departs from the
average of its neighbours: positive where the point sits below its
surroundings, negative where it sits above. The heat equation
$\partial T/\partial t = \alpha\nabla^{2}T$, the wave equation
$\partial^{2}u/\partial t^{2} = c^{2}\nabla^{2}u$, and the Laplace
equation $\nabla^{2}\Phi = 0$ for steady fields are the three
canonical second-order field equations, each built on this operator.
The vector Laplacian satisfies
\[
\curl(\curl\vect{F}) = \grad(\divg\vect{F}) - \nabla^{2}\vect{F},
\]
the working tool for reducing the curl-curl form of the
electromagnetic wave equation to a componentwise Laplacian. A reader
who derives this identity once from the BAC-CAB rule applied to the
symbolic $\nabla$ never has to rederive it; a reader who trusts the
mnemonic without the once-through derivation eventually applies it in
a coordinate system where the componentwise reading fails, because
the vector Laplacian acts componentwise only in Cartesian
coordinates.
\end{mastery}

\begin{mastery}
\textit{The differential identities by index notation.} The same
$\varepsilon$-$\delta$ machinery that collapsed the algebraic
identities of section~8.2 derives the differential identities of this
section, with $\partial_{i} = \partial/\partial x_{i}$ playing the
role $u_{i}$ played there. Write the divergence as $\divg\vect{F} =
\partial_{i}F_{i}$, the gradient as $(\grad\Phi)_{i} =
\partial_{i}\Phi$, and the curl as $(\curl\vect{F})_{i} =
\varepsilon_{ijk}\partial_{j}F_{k}$. The two identities that close the
operator algebra fall out in two lines each. The curl of a gradient,
\[
(\curl\grad\Phi)_{i} = \varepsilon_{ijk}\partial_{j}\partial_{k}\Phi
= 0,
\]
vanishes because $\varepsilon_{ijk}$ is antisymmetric in $j, k$ while
$\partial_{j}\partial_{k}\Phi$ is symmetric in $j, k$ (Clairaut's
theorem), and the contraction of a symmetric pair with an
antisymmetric pair is identically zero. The divergence of a curl,
\[
\divg(\curl\vect{F}) = \partial_{i}\varepsilon_{ijk}\partial_{j}F_{k}
= \varepsilon_{ijk}\partial_{i}\partial_{j}F_{k} = 0,
\]
vanishes for the same reason, antisymmetry of $\varepsilon$ in $i, j$
against symmetry of $\partial_{i}\partial_{j}$. The curl-of-curl
identity follows from the $\varepsilon$-$\delta$ contraction in one
pass:
\[
(\curl\curl\vect{F})_{i} = \varepsilon_{ijk}\partial_{j}
\varepsilon_{klm}\partial_{l}F_{m} = (\delta_{il}\delta_{jm} -
\delta_{im}\delta_{jl})\partial_{j}\partial_{l}F_{m} =
\partial_{i}(\partial_{m}F_{m}) - \partial_{j}\partial_{j}F_{i},
\]
which reads $\grad(\divg\vect{F}) - \nabla^{2}\vect{F}$, the identity
the previous box stated as a mnemonic. A reader who proves it once
this way sees why the vector Laplacian $\nabla^{2}\vect{F} =
\partial_{j}\partial_{j}F_{i}$ acts componentwise only in Cartesian
coordinates: the step that pulled $\partial_{j}\partial_{j}$ out as a
scalar operator on $F_{i}$ used the constancy of the Cartesian basis,
which the $\unitvec{\rho}, \unitvec{\phi}$ and $\unitvec{r},
\unitvec{\theta}, \unitvec{\phi}$ bases lack. The index notation is
the one apparatus that derives every vector identity, algebraic and
differential, from two symbols and one contraction rule, and it is the
language the constitutive laws of Volumes III and V are written in.
\end{mastery}

\begin{mastery}
\textit{Tensors, the stress field, and the divergence of a tensor.}
A force on a beam is a vector; the internal state of stress in the
beam's material is a \engterm{tensor}, a linear map that returns the
traction vector on any plane through a point as a function of that
plane's normal. The \engterm{Cauchy stress tensor} $\tens{\sigma}$ has
nine components $\sigma_{ij}$, where $\sigma_{ij}$ is the $i$-component
of force per unit area on the plane whose normal is $\unitvec{x}_{j}$.
The traction on a plane of normal $\unitvec{n}$ is $\vect{t} =
\tens{\sigma}\,\unitvec{n}$, or in index notation $t_{i} =
\sigma_{ij}n_{j}$. The strain-gauge rosette of section~8.1 measured
the companion strain tensor $\tens{\varepsilon}$ one projection at a
time, which is why three gauges resolved a three-component plane
tensor. The diagonal components $\sigma_{11}, \sigma_{22},
\sigma_{33}$ are the normal stresses (tension and compression); the
off-diagonal components are the shears. Force balance on an
infinitesimal material element gives the equilibrium equation
\[
\divg\tens{\sigma} + \vect{f} = \vect{0}, \qquad
(\divg\tens{\sigma})_{i} = \partial_{j}\sigma_{ij},
\]
where $\vect{f}$ is the body force per unit volume (gravity, for
instance). The \engterm{divergence of a tensor} contracts one index of
the gradient against the tensor, returning a vector, exactly the
divergence theorem applied component by component: the surface
integral of the traction $\tens{\sigma}\,\unitvec{n}$ over a closed
surface equals the volume integral of $\divg\tens{\sigma}$, which is
the net internal force the material exerts on the region. Volume~III
develops the stress tensor as the working object of solid mechanics
\cite{text:timoshenko-plates-shells}; the point here is that the
divergence theorem of this chapter, read for a tensor rather than a
vector, is the equilibrium equation of every loaded structure. The
metric tensor $g_{ij} = \unitvec{e}_{i}\cdot\unitvec{e}_{j}$ is the
tensor that records the lengths and angles a coordinate system
measures with; in Cartesian coordinates $g_{ij} = \delta_{ij}$ and the
distinction between upper and lower indices collapses, which is why
this chapter could write every dot product as $u_{i}v_{i}$ without
tracking the metric. In the curvilinear systems of the next
subsection, and in the curved spacetime of general relativity, the
metric carries the scale factors $h_{i}$ and the dot product becomes
$g_{ij}u^{i}v^{j}$, which is where the tensor apparatus earns its
generality.
\end{mastery}

\subsection{Operators in cylindrical and spherical coordinates}

The componentwise expressions in cylindrical and spherical
coordinates are not the same as in Cartesian, because the basis
vectors themselves vary in space. The forms below are the working
results; the derivations are exercises in the chain rule that the
reader can do once and reread when needed.

\paragraph{Cylindrical $(\rho, \phi, z)$.}
\begin{align*}
\grad\Phi &= \pdv{\Phi}{\rho}\unitvec{\rho} +
  \frac{1}{\rho}\pdv{\Phi}{\phi}\unitvec{\phi} +
  \pdv{\Phi}{z}\unitvec{z}, \\
\divg\vect{F} &= \frac{1}{\rho}\pdv{(\rho F_{\rho})}{\rho} +
  \frac{1}{\rho}\pdv{F_{\phi}}{\phi} + \pdv{F_{z}}{z}.
\end{align*}

\paragraph{Spherical $(r, \theta, \phi)$.}
\begin{align*}
\grad\Phi &= \pdv{\Phi}{r}\unitvec{r} +
  \frac{1}{r}\pdv{\Phi}{\theta}\unitvec{\theta} +
  \frac{1}{r\sin\theta}\pdv{\Phi}{\phi}\unitvec{\phi}, \\
\divg\vect{F} &= \frac{1}{r^{2}}\pdv{(r^{2} F_{r})}{r} +
  \frac{1}{r\sin\theta}\pdv{(\sin\theta\, F_{\theta})}{\theta} +
  \frac{1}{r\sin\theta}\pdv{F_{\phi}}{\phi}.
\end{align*}

The radial factors $\rho$, $r^{2}$, and $r\sin\theta$ in the
divergence expressions are the Jacobians of the surface elements at
constant $\rho$, $r$, and $\theta$ respectively. They are the same
factors that appeared in the volume and area elements above; the
agreement is not a coincidence but a statement of the divergence
theorem at infinitesimal scale.

\begin{example}[Divergence of a radial field in cylindrical coordinates]
A line source along the $z$-axis drives a purely radial flow
$\vect{F} = (C/\rho)\unitvec{\rho}$, the two-dimensional analogue of a
point source, with $C$ a constant. The Cartesian instinct is to
compute three partial derivatives; the cylindrical form is one line.
With $F_{\rho} = C/\rho$, $F_{\phi} = 0$, $F_{z} = 0$,
\[
\divg\vect{F} = \frac{1}{\rho}\pdv{(\rho F_{\rho})}{\rho} =
\frac{1}{\rho}\pdv{(C)}{\rho} = 0 \quad(\rho > 0).
\]
The divergence vanishes away from the axis: the $1/\rho$ falloff of
the field is exactly the rate that keeps the outward flux through
every concentric cylinder equal. The flux per unit length through a
cylinder of radius $a$ is $F_{\rho}(a)\cdot 2\pi a = (C/a)(2\pi a) =
2\pi C$, independent of $a$, confirming the zero divergence away from
the source. The source itself sits on the axis, the cylindrical
coordinate singularity, where the formula does not apply. The
cylindrical divergence formula does the bookkeeping the Cartesian
form would obscure behind a chain of product-rule terms.
\end{example}

The curl carries the same curvilinear factors. In cylindrical
coordinates the three components are
\begin{align*}
(\curl\vect{F})_{\rho} &= \frac{1}{\rho}\pdv{F_{z}}{\phi} -
  \pdv{F_{\phi}}{z}, \\
(\curl\vect{F})_{\phi} &= \pdv{F_{\rho}}{z} - \pdv{F_{z}}{\rho}, \\
(\curl\vect{F})_{z} &= \frac{1}{\rho}\pdv{(\rho F_{\phi})}{\rho} -
  \frac{1}{\rho}\pdv{F_{\rho}}{\phi},
\end{align*}
and in spherical coordinates,
\begin{align*}
(\curl\vect{F})_{r} &= \frac{1}{r\sin\theta}\left[
  \pdv{(\sin\theta\,F_{\phi})}{\theta} - \pdv{F_{\theta}}{\phi}\right], \\
(\curl\vect{F})_{\theta} &= \frac{1}{r}\left[\frac{1}{\sin\theta}
  \pdv{F_{r}}{\phi} - \pdv{(r F_{\phi})}{r}\right], \\
(\curl\vect{F})_{\phi} &= \frac{1}{r}\left[\pdv{(r F_{\theta})}{r} -
  \pdv{F_{r}}{\theta}\right].
\end{align*}
The factors of $1/\rho$, $1/r$, and $1/(r\sin\theta)$ are not
ornamental. They are the inverse of the distance a unit change in the
angular coordinate sweeps out at the point of evaluation, and a
calculation that drops them returns a curl with the wrong units.

\begin{example}[Curl of a rigid rotation in cylindrical coordinates]
A fluid in solid-body rotation about the $z$-axis at angular rate
$\Omega$ has velocity $\vect{v} = \Omega\rho\,\unitvec{\phi}$: every
point circles the axis with speed proportional to its radius. The
Cartesian form is $\vect{v} = (-\Omega y, \Omega x, 0)$, whose curl is
the familiar $(\curl\vect{v})_{z} = \partial(\Omega x)/\partial x -
\partial(-\Omega y)/\partial y = 2\Omega$. The cylindrical form
reaches the same number in one term. With $F_{\rho} = 0$,
$F_{\phi} = \Omega\rho$, $F_{z} = 0$ independent of $\phi$ and $z$,
\[
(\curl\vect{v})_{z} = \frac{1}{\rho}\pdv{(\rho\cdot\Omega\rho)}{\rho}
= \frac{1}{\rho}\pdv{(\Omega\rho^{2})}{\rho} = \frac{2\Omega\rho}{\rho}
= 2\Omega,
\]
uniform everywhere. The curl of a rigid rotation is twice the angular
velocity, the interpretation stated when the curl was introduced: a
paddle wheel placed anywhere in solid-body rotation spins at the bulk
rate $\Omega$, and the curl reports $2\Omega$. The $1/\rho$ factor and
the $\rho$ inside the derivative conspire to leave the radius out of
the answer, which is the geometric statement that solid-body rotation
has the same local rotation rate at every radius. A reader who omits
the $\rho$ factors gets a curl that grows with radius, contradicting
the uniform-rotation picture.
\end{example}

\begin{example}[Divergence of an inverse-square field in spherical
coordinates]
The point-source field $\vect{F} = (C/r^{2})\unitvec{r}$ is the model
for gravity and electrostatics. The Cartesian divergence requires the
quotient rule on three components; the spherical form is one term.
With $F_{r} = C/r^{2}$, $F_{\theta} = F_{\phi} = 0$,
\[
\divg\vect{F} = \frac{1}{r^{2}}\pdv{(r^{2}\cdot C/r^{2})}{r} =
\frac{1}{r^{2}}\pdv{(C)}{r} = 0 \quad(r > 0).
\]
The divergence vanishes everywhere the formula applies, exactly as the
Cartesian computation of the worked divergence-theorem example below
reports. The $r^{2}$ inside the derivative cancels the $1/r^{2}$ of
the field, leaving a constant whose radial derivative is zero. The
flux through a sphere of radius $a$ is $F_{r}(a)\cdot 4\pi a^{2} =
(C/a^{2})(4\pi a^{2}) = 4\pi C$, independent of $a$, which is the
radius-independent flux that the divergence theorem and Gauss's law
both record. The source again hides at the coordinate singularity
$r = 0$, where the smooth formula does not reach; the failure section
returns to this trap.
\end{example}

\begin{example}[Curl of the magnetic field around a wire]
A long straight wire along the $z$-axis carrying current $I$ produces
a magnetic field that circles the wire, $\vect{B} = \frac{\mu_{0}I}
{2\pi\rho}\unitvec{\phi}$, with $\mu_{0}$ the permeability of free
space. The field is purely azimuthal and falls as $1/\rho$, the same
falloff the radial line-source field had, now wrapped around the axis
rather than pointing away from it. The cylindrical curl, with
$B_{\rho} = 0$, $B_{\phi} = \mu_{0}I/(2\pi\rho)$, $B_{z} = 0$, all
independent of $\phi$ and $z$, has only the $z$-component to check
through the $\rho$-factored term,
\[
(\curl\vect{B})_{z} = \frac{1}{\rho}\pdv{(\rho B_{\phi})}{\rho} =
\frac{1}{\rho}\pdv{}{\rho}\!\Big(\rho\cdot\frac{\mu_{0}I}{2\pi\rho}
\Big) = \frac{1}{\rho}\pdv{}{\rho}\!\Big(\frac{\mu_{0}I}{2\pi}\Big)
= 0 \quad(\rho > 0),
\]
and the $\rho$- and $\phi$-components vanish because $B_{z} = 0$ and
nothing depends on $\phi$ or $z$. The curl is zero everywhere off the
axis, yet the circulation of $\vect{B}$ around any loop enclosing the
wire is $\oint\vect{B}\cdot\dd\vect{r} = B_{\phi}(\rho)\cdot 2\pi\rho =
\mu_{0}I$, independent of the loop radius. By Stokes's theorem the
circulation equals the curl flux through the loop, which is zero off
the axis, so the entire $\mu_{0}I$ comes from the axis itself, where
the current density is concentrated and the smooth curl formula does
not apply. This is Ampère's law, and it is the magnetostatic twin of
the inverse-square example: a field whose local operator vanishes
everywhere it is computed, yet whose integral around (or through) the
source is fixed by the source, because the source hides at the
coordinate singularity. Volume~IV builds the magnetostatic device on
this single calculation; the curl formula and the $\rho$ scale factor
are the tools of this chapter
\cite{text:v2ch08-griffiths-em-2017}.
\end{example}

\begin{mastery}
\textit{Scale factors and the general curvilinear operators.} The
cylindrical and spherical formulas above are special cases of one
pattern. An orthogonal curvilinear system with coordinates
$(q_{1}, q_{2}, q_{3})$ has \engterm{scale factors}
$h_{i} = \norm{\partial\vect{r}/\partial q_{i}}$, the distance a unit
change in $q_{i}$ moves the point. The arc-length element is
$\dd s^{2} = h_{1}^{2}\dd q_{1}^{2} + h_{2}^{2}\dd q_{2}^{2} +
h_{3}^{2}\dd q_{3}^{2}$, and the volume element is
$\dd V = h_{1}h_{2}h_{3}\,\dd q_{1}\dd q_{2}\dd q_{3}$. For Cartesian,
$(h_{1}, h_{2}, h_{3}) = (1, 1, 1)$; for cylindrical $(\rho, \phi, z)$,
$(1, \rho, 1)$; for spherical $(r, \theta, \phi)$, $(1, r, r\sin\theta)$.
The three differential operators read, in any orthogonal system,
\[
\grad\Phi = \sum_{i}\frac{1}{h_{i}}\pdv{\Phi}{q_{i}}\unitvec{q}_{i},
\qquad
\divg\vect{F} = \frac{1}{h_{1}h_{2}h_{3}}\sum_{i}
  \pdv{}{q_{i}}\!\left(\frac{h_{1}h_{2}h_{3}}{h_{i}}F_{i}\right),
\]
and the curl is the determinant
\[
\curl\vect{F} = \frac{1}{h_{1}h_{2}h_{3}}\det
\begin{pmatrix}
h_{1}\unitvec{q}_{1} & h_{2}\unitvec{q}_{2} & h_{3}\unitvec{q}_{3} \\
\partial/\partial q_{1} & \partial/\partial q_{2} &
  \partial/\partial q_{3} \\
h_{1}F_{1} & h_{2}F_{2} & h_{3}F_{3}
\end{pmatrix}.
\]
Substituting the spherical scale factors $(1, r, r\sin\theta)$
reproduces the spherical divergence and curl above without any
chain-rule labour; the Jacobian $h_{1}h_{2}h_{3} = r^{2}\sin\theta$ is
the volume-element factor and the divergence-theorem factor at once. A
reader who memorises the three scale factors for each system, rather
than the six operator formulas, reconstructs every formula on demand
and reads the field equations of Volumes IV, V, and IX in whatever
coordinate system the geometry forces \cite{text:v2ch08-marsden-tromba-2012}.
The Laplacian follows by composing: $\nabla^{2}\Phi = \divg(\grad\Phi)$
with the two formulas above, which is how the separable solutions of
the heat and wave equations in cylindrical and spherical geometry
acquire their Bessel-function and Legendre-polynomial form in
Volume~II Chapter~12.
\end{mastery}

\section{Stokes, Gauss, Green: the three integral theorems}\pathtag{core}

The three integral theorems of vector calculus link the local
behaviour of a field to its accumulated boundary behaviour. Each is
a balance statement: an interior accumulation equals a boundary flux.
The three theorems are the chapter's formal home for the balance
archetype.

\subsection{Green's theorem}

Green's theorem connects a line integral around a simple closed curve
in the plane to a double integral over the region the curve encloses.
For a vector field $\vect{F} = (P, Q)$ in $\R^{2}$ and a simple
closed curve $C$ traversed counterclockwise, enclosing the region
$D$,
\[
\oint_{C} \vect{F} \cdot \dd\vect{r} =
\oint_{C} (P \,\dd x + Q \,\dd y) =
\iint_{D} \left(\pdv{Q}{x} - \pdv{P}{y}\right) \dd A.
\]
The left side is the circulation of $\vect{F}$ around $C$; the right
side is the integral over $D$ of the scalar that a two-dimensional
reader recognises as the $z$-component of the curl. Green's theorem
is Stokes's theorem in two dimensions.

A working application: the area of a planar region $D$ bounded by
$C$ equals
\[
\text{Area}(D) = \oint_{C} x \,\dd y = -\oint_{C} y \,\dd x =
\frac{1}{2}\oint_{C} (x \,\dd y - y \,\dd x).
\]
The third form is the antisymmetric \engterm{shoelace formula}: a
land surveyor walking the boundary of a parcel and recording $x$ and
$y$ at each turn computes the parcel's area by accumulating
$\tfrac{1}{2}(x_{i} y_{i+1} - x_{i+1} y_{i})$ around the boundary.

\begin{example}[Green's theorem, both sides]
Take $\vect{F} = (P, Q) = (-y, x)$ and the unit disc $D$ of radius
$1$, with boundary $C$ the unit circle traversed counterclockwise.
The right side of Green's theorem is
\[
\iint_{D}\Big(\pdv{Q}{x} - \pdv{P}{y}\Big)\dd A =
\iint_{D}\big(1 - (-1)\big)\dd A = 2\,\text{Area}(D) = 2\pi.
\]
The left side, computed directly: parametrise $C$ by $(\cos t,
\sin t)$, $t \in [0, 2\pi]$, so $\dd\vect{r} = (-\sin t, \cos t)\dd t$
and $\vect{F} = (-\sin t, \cos t)$. Then
\[
\oint_{C}\vect{F}\cdot\dd\vect{r} = \int_{0}^{2\pi}(\sin^{2} t +
\cos^{2} t)\,\dd t = \int_{0}^{2\pi}\dd t = 2\pi.
\]
The two sides agree. The field $(-y, x)$ circulates counterclockwise
with circulation $2\pi$ around the unit circle, and Green's theorem
reads that circulation as twice the enclosed area, because the
$z$-curl of the field is the constant $2$. A surveyor or a planimeter
exploits the same identity in reverse: the boundary line integral
returns the enclosed area without the interior ever being sampled.
\end{example}

\begin{example}[Area of a polygon by the shoelace formula]
A quadrilateral parcel has corners $(0, 0)$, $(40, 0)$, $(50, 30)$,
$(10, 20)$, in metres, listed counterclockwise. The shoelace sum is
\[
2A = \sum_{i}(x_{i}y_{i+1} - x_{i+1}y_{i}) =
(0\cdot 0 - 40\cdot 0) + (40\cdot 30 - 50\cdot 0) +
(50\cdot 20 - 10\cdot 30) + (10\cdot 0 - 0\cdot 20),
\]
which evaluates to $0 + 1200 + (1000 - 300) + 0 = 1900$, so
$A = 950\,\si{\meter\squared}$. Listing the corners clockwise would
return $-950$; the sign reports the traversal direction, and the area
is the absolute value. The shoelace formula is Green's theorem
specialised to a piecewise-linear boundary, with each edge
contributing one term of the sum. A spreadsheet column of corner
coordinates and a single shoelace formula compute any polygon's area
without trigonometry.
\end{example}

\subsection{The divergence theorem (Gauss)}

The \engterm{divergence theorem} connects the volume integral of the
divergence of a vector field over a region to the flux of the field
through the region's boundary. For a sufficiently smooth vector field
$\vect{F}$ on a region $V \subset \R^{3}$ with closed, oriented
boundary $\partial V$ whose unit normal points outward,
\[
\iiint_{V} \divg\vect{F} \,\dd V =
\oint_{\partial V} \vect{F} \cdot \unitvec{n} \,\dd A.
\]
The integrand on the left is the local rate of net outflow per unit
volume, integrated over the volume. The right side is the integrated
net flux through the boundary. The theorem says the two are equal:
whatever is created or destroyed inside $V$ must, in the steady
state, cross $\partial V$.

\begin{archetype}[Balance, divergence theorem form]
A balance is an accounting equation. Volume~I introduced the form
``what is inside equals what came in minus what went out, plus what
was created.'' The divergence theorem is the same statement in the
calculus of fields: the total interior source rate equals the total
boundary outflow rate. Mass conservation in fluid mechanics, charge
conservation in electromagnetism, energy conservation in
thermodynamics, and species conservation in reactive transport are
all instances of this single archetype, with $\vect{F}$ being the
appropriate flux density and $\divg\vect{F}$ the appropriate source
density.
\end{archetype}

\begin{example}[The continuity equation from the divergence theorem]
The divergence theorem turns the verbal mass balance into a
differential equation. Let $\rho(\vect{r}, t)$ be the mass density and
$\vect{v}$ the velocity, so $\vect{J} = \rho\vect{v}$ is the mass-flux
density. For any fixed volume $V$, the mass inside is $\iiint_{V}\rho
\,\dd V$, and its rate of change equals the negative net outward mass
flux through the boundary (mass leaving the volume decreases the
interior mass):
\[
\dv{}{t}\iiint_{V}\rho\,\dd V = -\oint_{\partial V}\rho\vect{v}\cdot
\unitvec{n}\,\dd A.
\]
Apply the divergence theorem to the right side to convert the surface
integral to a volume integral, and move the time derivative inside the
left integral over the fixed volume:
\[
\iiint_{V}\pdv{\rho}{t}\,\dd V = -\iiint_{V}\divg(\rho\vect{v})\,\dd V.
\]
Because the volume $V$ is arbitrary, the integrands must be equal at
every point, giving the \engterm{continuity equation}
\[
\pdv{\rho}{t} + \divg(\rho\vect{v}) = 0.
\]
The arbitrary-volume argument is the workhorse of field theory: a
balance that holds over every volume forces a pointwise relation
between the local rate of accumulation and the local divergence of the
flux. For an incompressible fluid $\rho$ is constant, and the equation
reduces to $\divg\vect{v} = 0$, the incompressibility condition used
throughout this chapter. The same derivation, with charge density in
place of mass density, produces the charge-continuity equation of
Volume~IV; with thermal energy density, the heat equation of
Volume~II Chapter~12.
\end{example}

A worked specialisation of the theorem produces the chapter's
project formula. Take $\vect{F}(\vect{r}) = \vect{r}$, the position
vector field. Its divergence in Cartesian coordinates is
\[
\divg\vect{r} = \pdv{x}{x} + \pdv{y}{y} + \pdv{z}{z} = 3.
\]
The divergence theorem then reads
\[
\iiint_{V} 3 \,\dd V = \oint_{\partial V} \vect{r} \cdot \unitvec{n}
\,\dd A,
\]
which rearranges to
\[
V = \frac{1}{3} \oint_{\partial V} \vect{r} \cdot \unitvec{n} \,\dd A.
\]
The volume of any region equals one-third of the surface integral of
$\vect{r} \cdot \unitvec{n}$ over its closed boundary. This is the
formula the project asks the reader to apply to a sampled household
object.

\begin{figure}[h]
  \centering
  % Flux of a vector field through a closed surface (divergence theorem).
% A closed blob with an interior source; field lines emanate and cross
% the boundary; outward normals shown on the boundary.
\begin{tikzpicture}[>=Stealth, scale=1.0]

  % closed surface (a smooth blob)
  \draw[thick, engblack]
    plot[smooth cycle, tension=0.8]
    coordinates {(-2,0) (-1.2,1.4) (0.4,1.7) (1.8,0.9) (2.0,-0.6)
                 (0.7,-1.6) (-1.1,-1.3)};

  % interior source
  \fill[engred] (-0.1,0.1) circle (0.08);
  \node[engred, font=\footnotesize, below] at (-0.1,0.05)
    {$\divg\vect{F}>0$};

  % field lines from source crossing boundary
  \foreach \ang in {20,70,130,200,260,320} {
    \draw[->, engblue]
      (-0.1,0.1) -- ({-0.1+2.4*cos(\ang)},{0.1+2.0*sin(\ang)});
  }

  % a few outward normals on the boundary
  \draw[->, engamber, thick] (1.8,0.9) -- ++(0.55,0.32)
    node[right, font=\footnotesize]{$\unitvec{n}$};
  \draw[->, engamber, thick] (-2,0) -- ++(-0.55,0.0)
    node[left, font=\footnotesize]{$\unitvec{n}$};
  \draw[->, engamber, thick] (0.7,-1.6) -- ++(0.18,-0.55)
    node[below, font=\footnotesize]{$\unitvec{n}$};

  \node[font=\footnotesize] at (1.5,1.7) {$\partial V$};

  \node[font=\footnotesize, align=center] at (0,-2.6)
    {$\displaystyle\oint_{\partial V}\vect{F}\cdot\unitvec{n}\,\dd A
      = \iiint_{V}\divg\vect{F}\,\dd V$};

\end{tikzpicture}

  \caption[Flux of a vector field through a closed surface.]{The
  divergence theorem equates the net outward flux of $\vect{F}$
  through the closed surface $\partial V$ to the integral of
  $\divg\vect{F}$ over the enclosed volume. Field lines that begin
  inside $V$ (at a source, $\divg\vect{F} > 0$) must cross the
  boundary; the surface flux counts exactly the interior source
  strength. The outward unit normal $\unitvec{n}$ on each surface
  patch sets the sign of the local flux contribution. Source:
  authors' diagram.}
  \label{fig:vol02:ch08:flux-surface}
\end{figure}

\begin{example}[Divergence theorem on a box]
The field $\vect{F} = (x^{2}, y^{2}, z^{2})$ crosses the boundary of
the unit cube $[0, 1]^{3}$. Computing the flux face by face would
require six surface integrals. The divergence theorem replaces them
with one volume integral. The divergence is $\divg\vect{F} = 2x +
2y + 2z$, so
\[
\oint_{\partial V}\vect{F}\cdot\unitvec{n}\,\dd A =
\iiint_{V}(2x + 2y + 2z)\,\dd V = \int_{0}^{1}\!\int_{0}^{1}\!
\int_{0}^{1}(2x + 2y + 2z)\,\dd x\,\dd y\,\dd z.
\]
By symmetry each of the three terms integrates to the same value:
$\int_{0}^{1}2x\,\dd x = 1$ times the unit measure of the other two
axes, so the total is $1 + 1 + 1 = 3$. A direct face-by-face check
confirms it: the face at $x = 1$ contributes $\int\!\int 1^{2}\,
\dd y\,\dd z = 1$ and the face at $x = 0$ contributes $0$, and the
same for $y$ and $z$, summing to $3$. The volume integral does in one
step what six surface integrals do in six, which is the practical
reason the divergence theorem is the working tool for flux
accounting.
\end{example}

\begin{example}[Flux of an inverse-square field: the source matters]
The radial field $\vect{F} = \unitvec{r}/r^{2} = \vect{r}/r^{3}$ is
the model for a point source (a gravitational or electrostatic field).
Its divergence is zero everywhere except the origin: a direct
computation gives $\divg(\vect{r}/r^{3}) = 0$ for $r > 0$. The flux
through any sphere of radius $a$ centred at the origin is, by direct
surface integration with $\unitvec{n} = \unitvec{r}$ and
$\vect{F}\cdot\unitvec{n} = 1/a^{2}$,
\[
\oint_{S}\vect{F}\cdot\unitvec{n}\,\dd A = \frac{1}{a^{2}}\cdot
4\pi a^{2} = 4\pi,
\]
independent of the radius $a$. The flux is $4\pi$ through every sphere,
yet the divergence integrates to zero over the shell between two
spheres, consistent with $\divg\vect{F} = 0$ there. The naive
application of the divergence theorem to the solid ball would give
$\iiint 0\,\dd V = 0$, contradicting the surface flux of $4\pi$. The
resolution is that $\vect{F}$ is singular at the origin, where
$\divg\vect{F}$ is not zero but a concentrated source (a Dirac delta
of strength $4\pi$). The divergence theorem requires the field to be
smooth throughout the enclosed volume; a hidden singularity is the
most common way the theorem is misapplied, and the failure section
returns to the coordinate-singularity version of this trap.
\end{example}

\begin{example}[Volume of a cone by the position-field formula]
The chapter project computes a household object's volume from
$V = \tfrac{1}{3}\oint_{\partial V}\vect{r}\cdot\unitvec{n}\,\dd A$.
A cone of base radius $a$ and height $h$, apex at the origin and base
at $z = h$, gives a clean test of the formula before it meets a real
object. The cone has two surfaces: the flat base at $z = h$ and the
lateral surface. On the base, the outward normal is $\unitvec{n} =
+\unitvec{z}$ and $\vect{r}\cdot\unitvec{n} = z = h$ everywhere, so the
base contributes $h$ times its area $\pi a^{2}$, that is, $\pi a^{2}h$.
On the lateral surface the position vector lies in the tangent plane
of the cone through the apex, so $\vect{r}\cdot\unitvec{n} = 0$ at
every point and the lateral surface contributes nothing. The closed
surface integral is therefore
\[
\oint_{\partial V}\vect{r}\cdot\unitvec{n}\,\dd A = \pi a^{2}h,
\qquad V = \frac{1}{3}\,\pi a^{2}h,
\]
the textbook cone volume. The point is not the formula, which the reader knows, but
the method: the closed-surface flux of the position field, divided by
three, returns the volume of any shape, smooth or sampled, with no
interior integration. The project replaces the cone's two analytic
surfaces with a triangulated mesh and replaces the integral with a
sum of signed tetrahedra, but the identity being exercised is the one
verified here on a shape whose answer is known.
\end{example}

\begin{estimation}
\textbf{Question.} An insulated water tank with internal volume
$0.20\,\si{\meter\cubed}$ holds water at $60\,\degC$ in a basement at
$15\,\degC$. The tank's walls are a $50\,\si{\milli\meter}$ layer of
polyurethane foam, thermal conductivity $k \approx 0.025\,\si{\watt
\per\meter\per\kelvin}$. Use the divergence theorem to estimate the
steady-state heat-loss rate from the tank, and the daily energy that
the heater must replace. State a factor-of-three confidence range.

\smallskip
\textit{Estimate, before reading on.}

\smallskip
The divergence theorem applied to the heat flux $\vect{q} = -k\grad T$
relates the volumetric source rate to the surface flux:
$\iiint_V \divg\vect{q}\,\dd V = \oint_{\partial V} \vect{q}\cdot
\unitvec{n}\,\dd A$. At steady state with the heater supplying the
interior source, the surface integral is the total wall heat loss.
Approximating the tank as a cube with side $\ell = V^{1/3} \approx
0.585\,\si{\meter}$, the wall area is $A \approx 6\ell^{2} \approx
2.05\,\si{\meter\squared}$. The wall behaves as a one-dimensional
conductor; the heat flux through each face is $\norm{\vect{q}} \approx
k \Delta T / d = 0.025 \times 45 / 0.050 = 22.5\,\si{\watt\per\meter
\squared}$. Total loss $\dot{Q} \approx 22.5 \times 2.05 \approx
46\,\si{\watt}$.

Daily energy is $\dot{Q} \times 86400\,\si{\second} \approx 4.0
\times 10^{6}\,\si{\joule} \approx 1.1\,\si{\kilo\watt\hour}$, the
order of one inexpensive lamp's daily consumption. A factor-of-three
band, $0.4$ to $3.3\,\si{\kilo\watt\hour\per\day}$, brackets the
plausible range when uncertainties on the wall thickness, the foam
conductivity, and the basement temperature compound. The dominant
uncertainty is the foam conductivity, which depends on density,
moisture, and ageing and can degrade by a factor of two over a decade
\cite{text:incropera2007}.

The estimate's structural value is the same as in the garden-hose
case in section~8.1: the divergence theorem reduces the total heat
loss to a one-dimensional surface flux times area, and the area times
flux to a single product the engineer reads off in seconds. The same
identity supports the design of attic insulation, the rating of a
domestic refrigerator, and the heat-budget calculation of a small
spacecraft, each with the same accounting and different numbers.
\end{estimation}

\subsection{Stokes's theorem}

\engterm{Stokes's theorem} connects the surface integral of the curl
of a vector field over an oriented surface to the line integral of
the field around the surface's boundary. For a sufficiently smooth
vector field $\vect{F}$ on a surface $S$ with positively oriented
boundary $\partial S$,
\[
\iint_{S} (\curl\vect{F}) \cdot \unitvec{n} \,\dd A =
\oint_{\partial S} \vect{F} \cdot \dd\vect{r}.
\]
The orientation convention is the right-hand rule: with the right
hand wrapped around the boundary in the direction of $\dd\vect{r}$,
the thumb points along $\unitvec{n}$. The left side is the total
``twist'' of $\vect{F}$ over $S$; the right side is the
circulation of $\vect{F}$ around the boundary $\partial S$. The
theorem says the two are equal: a surface integral of curl reports
the same number as a line integral around the surface's edge.

\begin{figure}[h]
  \centering
  % Stokes orientation: surface S with boundary curve, right-hand rule for normal.
\begin{tikzpicture}[>=Stealth, scale=1.0]
  % Surface: tilted ellipse representing the spanning surface
  \draw[fill=engblue!10, draw=engblue, thick]
    (0,0) ellipse [x radius=2.4, y radius=0.9];

  % Boundary curve with orientation arrows around the ellipse
  \foreach \ang in {15, 105, 195, 285} {
    \pgfmathsetmacro{\xs}{2.4*cos(\ang)}
    \pgfmathsetmacro{\ys}{0.9*sin(\ang)}
    \pgfmathsetmacro{\xt}{2.4*cos(\ang + 6)}
    \pgfmathsetmacro{\yt}{0.9*sin(\ang + 6)}
    \draw[->, very thick, engred] (\xs,\ys) -- (\xt,\yt);
  }

  % Normal vector
  \draw[->, very thick] (0.4,0.1) -- ++(0.4,1.6) node[right, font=\footnotesize]{$\unitvec{n}$};

  \node[font=\footnotesize, align=center] at (3.6,0.0) {$\partial S$ traced\\with right hand};
  \node[font=\footnotesize, align=center, anchor=west] at (1.0,2.0) {thumb $\to$ $\unitvec{n}$};

  \node[font=\footnotesize] at (-2.2,-0.6) {$S$};
  \node[font=\footnotesize, engred] at (0,-1.25) {$\partial S$};
\end{tikzpicture}

  \caption[Stokes orientation: right-hand rule between boundary and
  normal.]{Stokes's theorem links the surface integral of $\curl
  \vect{F}$ over $S$ to the line integral of $\vect{F}$ around
  $\partial S$. The orientation convention is the right-hand rule:
  fingers along the boundary in the parametrised direction, thumb
  along $\unitvec{n}$. Source: authors' diagram.}
  \label{fig:vol02:ch08:stokes-orientation}
\end{figure}

\begin{example}[Stokes's theorem, both sides]
Take $\vect{F} = (-y, x, 0)$ and let $S$ be the flat unit disc in the
$z = 0$ plane, with upward normal $\unitvec{n} = \unitvec{z}$ and
boundary $\partial S$ the unit circle traversed counterclockwise. The
curl is $\curl\vect{F} = (0, 0, 2)$, so the surface side is
\[
\iint_{S}(\curl\vect{F})\cdot\unitvec{n}\,\dd A = \iint_{S}2\,\dd A =
2\pi,
\]
twice the disc area. The boundary side, parametrising the circle by
$(\cos t, \sin t, 0)$, is the same line integral computed in the
Green's-theorem example, $\oint_{\partial S}\vect{F}\cdot\dd\vect{r} =
2\pi$. The two sides agree, as Stokes's theorem requires. Now replace
the flat disc with the upper unit hemisphere, which has the same
boundary circle. The curl flux through the hemisphere is also $2\pi$,
because $\divg(\curl\vect{F}) = 0$ guarantees the curl flux depends
only on the boundary, not the spanning surface. A reader who computes
the hemisphere integral the hard way, and a reader who computes the
flat-disc integral the easy way, get the same $2\pi$; Stokes's theorem
licences choosing whichever surface is easier.
\end{example}

\begin{example}[Stokes on a paraboloid cap, surface side computed
the hard way]
The previous example invoked surface-independence to avoid the curved
integral. This one computes the curved integral directly, to show the
machinery works and returns the boundary value. Take the field
$\vect{F} = (-y, x, xyz)$ and let $S$ be the paraboloid cap
$z = 1 - x^{2} - y^{2}$ for $z \geq 0$, oriented with an upward
normal, whose boundary $\partial S$ is the unit circle in the
$z = 0$ plane. The curl is
\[
\curl\vect{F} = \big(\partial(xyz)/\partial y - 0,\ 0 -
\partial(xyz)/\partial x,\ \partial x/\partial x - \partial(-y)/
\partial y\big) = (xz,\ -yz,\ 2).
\]
Parametrise the cap over the unit disc $D$ in the $(x, y)$ plane by
$\vect{r}(x, y) = (x, y, 1 - x^{2} - y^{2})$. The upward area-vector
element is $\vect{r}_{x}\times\vect{r}_{y} = (2x, 2y, 1)\,\dd x\,\dd y$.
The surface integral is
\[
\iint_{S}(\curl\vect{F})\cdot\unitvec{n}\,\dd A = \iint_{D}(xz, -yz, 2)
\cdot(2x, 2y, 1)\,\dd x\,\dd y = \iint_{D}\big(2x^{2}z - 2y^{2}z + 2
\big)\,\dd x\,\dd y,
\]
with $z = 1 - x^{2} - y^{2}$ on the cap. The term $2x^{2}z - 2y^{2}z =
2z(x^{2} - y^{2})$ integrates to zero over the disc by symmetry: the
disc is invariant under the swap $x \leftrightarrow y$, which sends
$x^{2} - y^{2}$ to its negative while leaving $z$ unchanged, so the
positive and negative contributions cancel. Only the constant term
survives,
\[
\iint_{D}2\,\dd x\,\dd y = 2\,\text{Area}(D) = 2\pi.
\]
The boundary side: on $\partial S$, $z = 0$, so $\vect{F} = (-y, x, 0)$,
and the line integral around the unit circle is the $2\pi$ already
computed. The curved surface integral and the boundary line integral
agree, and the agreement did not rely on surface-independence: the
$xyz$ term in the field contributed a genuine curved-surface integrand
that vanished only after the symmetry argument, while the part of the
curl that came from the planar circulation supplied the whole answer.
A reader who carries the full curl through the paraboloid integral, and
checks that the symmetric terms cancel rather than assuming they do,
has exercised the surface integral as a working tool rather than as a
corollary of a theorem.
\end{example}

\begin{estimation}
\textbf{Question.} A light aircraft wing has chord $c = 1.5\,\si{\meter}$,
span $b = 10\,\si{\meter}$, and flies at $v = 50\,\si{\meter\per
\second}$ through air of density $\rho = 1.2\,\si{\kilogram\per\meter
\cubed}$. The Kutta-Joukowski theorem states that the lift per unit
span is $L' = \rho v\,\Gamma$, where $\Gamma = \oint\vect{v}\cdot
\dd\vect{r}$ is the circulation around an airfoil section. Estimate the
circulation required to lift a $900\,\si{\kilogram}$ aircraft, and
state a factor-of-three confidence range.

\smallskip
\textit{Estimate, before reading on.}

\smallskip
The weight to support is $W = mg = 900\times 9.8 \approx 8.8\times
10^{3}\,\si{\newton}$. Treating the lift as uniform along the span,
the lift per unit span is $L' = W/b = 8800/10 = 880\,\si{\newton\per
\meter}$. The Kutta-Joukowski relation inverts to
\[
\Gamma = \frac{L'}{\rho v} = \frac{880}{1.2\times 50} \approx
15\,\si{\meter\squared\per\second}.
\]
The circulation is a line integral of velocity around the airfoil, a
direct application of the circulation integral of this chapter; by
Stokes's theorem it equals the flux of vorticity through any section
bounded by that loop. A factor-of-three band, $5$ to $44\,\si{\meter
\squared\per\second}$, brackets the spanwise lift distribution (which
is elliptical, not uniform, so the root section carries more than the
estimate and the tip less) and the unknown flight condition (climb
versus cruise changes the required lift). The dominant uncertainty is
the spanwise distribution; a working aerodynamicist computes $\Gamma$
from the actual lift distribution rather than the span-average used
here. The estimate's value is the structure: lift is circulation
times $\rho v$, and circulation is a line integral the engineer can
reason about with the tools of this chapter.
\end{estimation}

The working consequence is that a circulation depends only on the
curl of the field through any surface spanning the curve, not on the
particular surface chosen. Two different surfaces with the same
boundary give the same circulation. The reader meets this
path-independence again in Volume~III, where it is the working
basis of magnetostatics, and in Volume~IV, where it is the
working basis of electromagnetism.

\subsection{Why three theorems and not one}

The three theorems are three faces of the same fundamental statement.
The general form, the \engterm{generalised Stokes theorem}, is
\[
\int_{M} \dd\omega = \int_{\partial M} \omega,
\]
for a sufficiently nice differential form $\omega$ on a sufficiently
nice manifold $M$ \cite{text:strang2016}. Green is the case $M$ is a
region in $\R^{2}$ and $\omega$ is a 1-form. Gauss is the case $M$ is
a region in $\R^{3}$ and $\omega$ is a 2-form. Stokes is the case
$M$ is a surface in $\R^{3}$ and $\omega$ is a 1-form. The reader
who studies differential forms in a later course will recognise the
three theorems as the same theorem in three dimensions and codimensions;

\begin{mastery}
\textit{Differential forms in three lines.} The unification is worth a
concrete sketch even for the reader who will not take the later course.
A \engterm{differential $k$-form} is the integrand of a $k$-dimensional
integral: a $0$-form is a function $f$, a $1$-form is $P\,\dd x +
Q\,\dd y + R\,\dd z$ (the integrand of a line integral), a $2$-form is
$A\,\dd y\wedge\dd z + B\,\dd z\wedge\dd x + C\,\dd x\wedge\dd y$ (the
integrand of a flux integral), and a $3$-form is $g\,\dd x\wedge\dd y
\wedge\dd z$ (the integrand of a volume integral). The
\engterm{wedge product} $\wedge$ is antisymmetric, $\dd x\wedge\dd y =
-\dd y\wedge\dd x$, which is where every orientation sign in the
chapter comes from: the antisymmetry is the algebra of the right-hand
rule. The \engterm{exterior derivative} $\dd$ raises a $k$-form to a
$(k{+}1)$-form, and reading what $\dd$ does to each degree recovers the
three operators of this chapter: $\dd$ of a $0$-form is the gradient,
$\dd$ of a $1$-form is the curl, and $\dd$ of a $2$-form is the
divergence. The two identities $\curl(\grad\,\cdot) = \vect{0}$ and
$\divg(\curl\,\cdot) = 0$ are the single statement $\dd(\dd\omega) = 0$,
``the boundary of a boundary is empty,'' read in two degrees. The
generalised Stokes theorem $\int_{M}\dd\omega = \int_{\partial M}
\omega$ then specialises to all three integral theorems by choosing the
degree of $\omega$. A reader who carries this dictionary, $\dd =
\{\grad, \curl, \divg\}$ by degree and $\wedge =$ orientation, reads
the entire chapter as one theorem and one operator, which is the
viewpoint that generalises to the curved spacetimes of relativity and
the gauge fields of particle physics \cite{text:strang2016}. The cost
is one antisymmetric product and one degree-raising derivative; the
return is that the three theorems, three operators, and two identities
collapse to one of each.
\end{mastery}

\noindent
The working engineer who does not study forms still has the three
working tools, with the same orientation discipline, the same
boundary discipline, and the same balance interpretation.

\subsection{The working consequence: balance is a theorem}

\begin{principle}[Balance is a theorem]
A balance equation, written for any quantity that admits a flux
density and a source density, is the divergence theorem applied to
that quantity. The interior accumulation rate is the volume integral
of the source density. The boundary flux rate is the surface integral
of the flux density dotted with the outward normal. The two are
equal by the divergence theorem. A reader who has internalised this
fact has reduced ten domain-specific conservation laws to one
mathematical identity applied ten times.
\end{principle}

The principle is what the chapter offers Volume~III and beyond. The
mass-conservation equation of fluid mechanics, the charge-conservation
equation of electromagnetism, the energy-conservation equation of
thermodynamics, the species-conservation equation of chemical
engineering, and the population-conservation equation of biological
modelling are all the same identity. Each volume restates it in the
domain's working notation; the reader who has met the identity here
recognises the restatement as a notational variant rather than as a
new principle.

\subsection{A short-form historical example: Maxwell and the integral
theorems}

The most consequential application of the three integral theorems in
the engineering record is James Clerk Maxwell's 1865 reformulation of
electromagnetism. Maxwell's original \textit{Treatise on Electricity
and Magnetism} stated electromagnetism in twenty scalar equations.
Oliver Heaviside compressed them, between 1884 and 1885, into the
four vector equations the working engineer now reads as Maxwell's
equations \cite{text:v2ch08-griffiths-em-2017}. Two of the four are
integrals of $\divg\vect{E}$ and $\divg\vect{B}$ over a closed
surface, equated by the divergence theorem to surface fluxes; the
other two are integrals of $\curl\vect{E}$ and $\curl\vect{B}$ over
an open surface, equated by Stokes's theorem to line integrals around
the bounding curve. Volume~IV develops the four equations in their
working form. The point here is structural: the unification of
electric and magnetic phenomena into a single field theory rests on
the divergence theorem and Stokes's theorem applied to fields whose
local rates the experimentalist measures and whose integral
behaviour the engineer designs against. The same three theorems
that govern the chapter project are what made twentieth-century
electrical engineering possible.

A second short-form case: the finite-volume method, the workhorse of
modern computational fluid dynamics, computes the divergence of a
flux field on each cell by summing the flux through the cell's faces
and dividing by the cell volume. The method's name records what it is
doing: enforcing the divergence theorem at cell scale instead of at
infinitesimal scale \cite{text:v2ch08-marsden-tromba-2012}. A reader
who later opens a CFD or finite-element code finds the divergence
theorem, written as a balance over a finite cell, at the bottom of
every solver loop.

A reader who wants to see the three theorems hold numerically before
trusting them analytically can run the companion script at
\repopath{volumes/02-form/ch08-vectors-and-vector-calculus/code/verify_theorems.py},
which evaluates both sides of Green's, Gauss's, and Stokes's theorems
for the three worked examples above and reports the discrepancy. The
discrepancy falls toward floating-point noise as the discretisation
refines; the tabulated convergence sweep sits in
\repopath{data/theorem_verification.csv}. The exercise is worth doing
once: watching the boundary integral and the interior integral
converge to the same number, from two unrelated discretisations, is
the most direct evidence a reader can produce that the local and the
global descriptions of a field are the same statement.

\subsection{Case study: the field, capacitance, and breakdown margin
of a coaxial cable}\pathtag{standard}

The tools of this chapter converge on one device the reader has
handled without thinking about it: the coaxial cable that carries a
television signal, a radio-frequency feed, or an Ethernet-over-coax
run. A coaxial cable is an inner conductor of radius $a$ surrounded by
a concentric outer conductor of inner radius $b$, with an insulating
dielectric filling the gap. The geometry is cylindrically symmetric,
the field is radial, and every quantity the cable designer cares
about, the field strength at the inner conductor, the capacitance per
unit length, the voltage at which the dielectric breaks down, follows
from the divergence theorem applied in cylindrical coordinates. This
case study works the device end to end, from the symmetry argument to
the breakdown margin, using only the operators and theorems of this
chapter.

\paragraph{Step 1: the symmetry argument fixes the field's form.}
The inner conductor carries charge per unit length $\lambda$ (current
as of the static analysis; the radio-frequency case is Volume~IV). The
system is invariant under rotation about the axis and under translation
along it, away from the cable's ends. A field that respects those two
symmetries cannot depend on $\phi$ or $z$, and cannot have a $\phi$ or
$z$ component without breaking the symmetry. The electric field is
therefore purely radial and depends only on $\rho$,
\[
\vect{E} = E_{\rho}(\rho)\,\unitvec{\rho}.
\]
This is the payoff of the coordinate-choice discipline of section~8.3:
choosing cylindrical coordinates because the geometry has cylindrical
symmetry collapses a three-variable vector field to a single function
of one variable, before any equation is solved. A reader who attacked
the problem in Cartesian coordinates would carry $\sqrt{x^{2} + y^{2}}$
through every step and recover the same answer with ten times the
algebra.

\paragraph{Step 2: the divergence theorem gives the field.} Gauss's
law states that the flux of $\vect{E}$ through any closed surface
equals the enclosed charge divided by the permittivity
$\varepsilon$ of the dielectric, which is the divergence theorem
applied to $\vect{E}$ with $\divg\vect{E} = \rho_{q}/\varepsilon$ for
charge density $\rho_{q}$. Take the closed surface to be a coaxial
cylinder of radius $\rho$ (with $a < \rho < b$) and length $L$, the
\engterm{Gaussian surface}. The flux has three pieces: the two flat
end caps, where $\vect{E} = E_{\rho}\unitvec{\rho}$ is perpendicular to
the cap normal $\pm\unitvec{z}$ and contributes nothing, and the
curved side, where $\vect{E}$ is parallel to the outward normal
$\unitvec{\rho}$ and contributes $E_{\rho}(\rho)$ times the side area
$2\pi\rho L$. The enclosed charge is $\lambda L$. The divergence
theorem reads
\[
\oint_{\partial V}\vect{E}\cdot\unitvec{n}\,\dd A =
E_{\rho}(\rho)\,(2\pi\rho L) = \frac{\lambda L}{\varepsilon},
\qquad\text{so}\qquad
E_{\rho}(\rho) = \frac{\lambda}{2\pi\varepsilon\rho}.
\]
The field falls as $1/\rho$, the same falloff the cylindrical
line-source example produced for a different reason: it is the falloff
that keeps the flux through every coaxial cylinder equal to the fixed
enclosed charge. A direct check confirms the field is divergence-free
in the gap, where there is no charge: using the cylindrical divergence
formula, $\divg\vect{E} = \frac{1}{\rho}\partial(\rho\cdot\lambda/
(2\pi\varepsilon\rho))/\partial\rho = \frac{1}{\rho}\partial(\lambda/
(2\pi\varepsilon))/\partial\rho = 0$, as a charge-free region demands.

\begin{figure}[h]
  \centering
  % Coaxial cable cross-section: inner conductor radius a, outer radius b,
% radial field in the dielectric gap, and a Gaussian cylinder of radius rho.
\begin{tikzpicture}[>=Stealth, scale=1.0]

  % outer conductor (annulus boundary)
  \draw[thick, engblack] (0,0) circle (2.4);
  \node[engblack, font=\footnotesize] at (1.75,1.75) {$b$};
  \draw[engblack, dashed] (0,0) -- (60:2.4);

  % inner conductor
  \fill[enggray!40] (0,0) circle (0.6);
  \draw[thick, engblack] (0,0) circle (0.6);
  \node[engblack, font=\footnotesize] at (0.30,0.30) {$a$};
  \draw[engblack, dashed] (0,0) -- (30:0.6);

  % Gaussian surface (intermediate cylinder, edge-on as a circle)
  \draw[thick, engblue, dashed] (0,0) circle (1.5);
  \node[engblue, font=\footnotesize] at (-1.5,1.15) {$\rho$};

  % radial field arrows in the gap
  \foreach \ang in {0,45,90,135,180,225,270,315} {
    \draw[->, engred] (\ang:0.65) -- (\ang:1.45);
  }
  \node[engred, font=\footnotesize] at (3.2,0)
    {$\vect{E}=\dfrac{\lambda}{2\pi\varepsilon\rho}\,\unitvec{\rho}$};

  % label dielectric
  \node[font=\footnotesize, align=center] at (0,-3.0)
    {$E_{\rho}(\rho)\,(2\pi\rho L)=\dfrac{\lambda L}{\varepsilon}$};

\end{tikzpicture}

  \caption[Coaxial cable: radial field and Gaussian surface.]{Cross
  section of a coaxial cable. The inner conductor of radius $a$ carries
  charge per unit length $\lambda$; the field in the dielectric gap is
  radial and falls as $1/\rho$. The divergence theorem on the coaxial
  Gaussian cylinder of radius $\rho$ (blue dashed) equates the side
  flux $E_{\rho}(\rho)(2\pi\rho L)$ to the enclosed charge $\lambda L
  /\varepsilon$, fixing the field in one line. The end caps carry no
  flux because $\vect{E}\perp\unitvec{z}$ there. Source: authors'
  diagram.}
  \label{fig:vol02:ch08:coaxial-field}
\end{figure}

\paragraph{Step 3: the gradient gives the voltage.} The field is the
negative gradient of the potential, $\vect{E} = -\grad V$, which in
cylindrical coordinates with a radial field reads
$E_{\rho} = -\partial V/\partial\rho$. The voltage between the
conductors is the line integral of the field from the inner conductor
to the outer one, along any radial path (the field is conservative, so
the path does not matter):
\[
V_{ab} = \int_{a}^{b}E_{\rho}(\rho)\,\dd\rho = \frac{\lambda}
{2\pi\varepsilon}\int_{a}^{b}\frac{\dd\rho}{\rho} = \frac{\lambda}
{2\pi\varepsilon}\ln\!\frac{b}{a}.
\]
The logarithm is the signature of the cylindrical geometry: the field's
$1/\rho$ falloff integrates to a logarithm, so the voltage depends on
the \emph{ratio} $b/a$, not on the absolute dimensions. A cable scaled
up by a factor of ten in both radii has the same voltage for the same
charge per unit length.

\paragraph{Step 4: the capacitance per unit length.} The capacitance
per unit length is the charge per unit length divided by the voltage,
\[
C' = \frac{\lambda}{V_{ab}} = \frac{2\pi\varepsilon}{\ln(b/a)}.
\]
For a common cable with $a = 0.5\,\si{\milli\meter}$, $b = 1.75\,
\si{\milli\meter}$, and a polyethylene dielectric of relative
permittivity $\varepsilon_{r} = 2.25$ (so $\varepsilon =
\varepsilon_{r}\varepsilon_{0} = 2.25\times 8.85\times 10^{-12} =
1.99\times 10^{-11}\,\si{\farad\per\meter}$), the ratio is
$b/a = 3.5$, $\ln 3.5 = 1.25$, and
\[
C' = \frac{2\pi\times 1.99\times 10^{-11}}{1.25} = 1.00\times 10^{-10}
\,\si{\farad\per\meter} = 100\,\si{\pico\farad\per\meter}.
\]
The figure of order $100\,\si{\pico\farad\per\meter}$ is what a cable
datasheet reports, and a reader who has reached it from the divergence
theorem has rebuilt the datasheet entry from first principles. The
characteristic impedance the radio engineer designs against,
$Z_{0} = \sqrt{L'/C'}$ with $L'$ the inductance per unit length, rests
on this same $C'$; the $50\,\si{\ohm}$ and $75\,\si{\ohm}$ standard
impedances correspond to particular choices of the ratio $b/a$.

\paragraph{Step 5: the breakdown margin.} The field is strongest where
$\rho$ is smallest, at the inner conductor's surface $\rho = a$:
\[
E_{\max} = E_{\rho}(a) = \frac{\lambda}{2\pi\varepsilon a} =
\frac{V_{ab}}{a\ln(b/a)}.
\]
The inner conductor is the failure surface. For the cable above
carrying a working voltage $V_{ab} = 5\,\si{\kilo\volt}$ (a
cable-television trunk amplifier feed),
\[
E_{\max} = \frac{5000}{(0.5\times 10^{-3})(1.25)} = 8.0\times 10^{6}
\,\si{\volt\per\meter} = 8.0\,\si{\mega\volt\per\meter}.
\]
Polyethylene breaks down at roughly $20\,\si{\mega\volt\per\meter}$
(current as of 2026; the figure depends on purity, temperature, and
ageing, and a design uses a derated value \cite{text:v2ch08-griffiths-em-2017}).
The breakdown margin is $20/8.0 = 2.5$, which is thin for a safety-
critical run: a working cable designer targets a margin of three to
five and would either raise $a$ or choose a dielectric with a higher
breakdown strength. The design lever is the inner radius: because
$E_{\max}$ scales as $1/a$ for fixed voltage and fixed ratio,
enlarging the inner conductor lowers the peak field, which is why
high-voltage coaxial feeds use a thick inner conductor that looks
oversized for the current it carries.

\paragraph{Step 6: the optimisation the geometry hides.} For a fixed
outer radius $b$ and a fixed voltage $V_{ab}$, there is an inner radius
that minimises the peak field. Write $E_{\max} = V_{ab}/(a\ln(b/a))$
and minimise over $a$; setting the derivative of $a\ln(b/a)$ to zero
gives $\ln(b/a) = 1$, that is, $b/a = e \approx 2.72$. A cable built at
this ratio carries the most voltage before breakdown for a given outer
size, and the optimum is a pure consequence of the logarithmic
potential the cylindrical divergence theorem produced. The same
ratio-of-$e$ result governs the design of high-voltage bushings and
the spacing of coaxial vacuum feedthroughs; it is one calculation,
reusable across every cylindrical high-field device, and it falls out
of the chapter's three operators applied in the coordinate system the
geometry selected.

The case study used every tool of the chapter in sequence: the
symmetry argument to fix the field's form, the divergence theorem to
find the field, the gradient relation and a line integral to find the
voltage, an algebraic ratio to find the capacitance, and a derivative
to find the optimum. The reader who follows the chain from symmetry to
breakdown margin has seen the working logic of field engineering, in
which the coordinate choice, the integral theorem, and the local
operator are three steps of one calculation rather than three separate
topics. Volume~IV develops the electromagnetic device in full; the
skeleton is the vector calculus of this chapter.

\subsection{Case study: drag and lift on a body from a control-volume
flux balance}\pathtag{standard}

The coaxial-cable study read a field off a closed surface with the
divergence theorem. The second study reads two forces, drag and lift,
off the same surface integral machinery applied to the flux of
momentum and the circulation of velocity. The setting is a body, a
cylinder or an airfoil section, held in a steady stream of air or
water, and the question is the force the fluid exerts on it. A working
aerodynamicist or hydrodynamicist computes that force two ways that
this chapter has both built: a momentum-flux balance over a control
volume for the drag, and a circulation line integral for the lift.
The case study works the body end to end, from the control-volume
bookkeeping to the numerical force, using only the flux integral, the
divergence theorem, and the circulation integral.

\paragraph{Step 1: the control volume and the mass balance.}
Enclose the body in a fixed rectangular \engterm{control volume} $V$
whose upstream face is far enough ahead that the flow there is the
undisturbed free stream $\vect{v} = (U, 0)$, and whose downstream face
is in the body's wake, where the flow has slowed to a profile
$\vect{v} = (u(y), 0)$ with $u(y) < U$ behind the body. The top and
bottom faces are far enough from the body that the flow there is
nearly the free stream. Mass conservation over $V$ is the divergence
theorem applied to the velocity field with $\divg\vect{v} = 0$ for the
incompressible flow:
\[
\oint_{\partial V}\vect{v}\cdot\unitvec{n}\,\dd A = \iiint_{V}
\divg\vect{v}\,\dd V = 0.
\]
The net volume flux through the closed control surface is zero. The
upstream face admits $U$ per unit area; the wake face discharges the
smaller $u(y)$; the deficit must leave through the top and bottom
faces. The mass balance fixes that lateral outflow before any force is
computed, and a reader who omits it double-counts the momentum the
lateral faces carry.

\begin{figure}[h]
  \centering
  % Control volume around a body: uniform inflow, momentum-deficit wake outflow.
\begin{tikzpicture}[>=Stealth, scale=1.0]
  % Control-volume box
  \draw[draw=engblue, thick] (0,0) rectangle (7,4);

  % The body (cylinder cross-section)
  \fill[gray!55] (2.6,2.0) circle [radius=0.35];
  \node[font=\footnotesize] at (2.6,1.45) {body};

  % Upstream uniform inflow arrows (left face)
  \foreach \yy in {0.5, 1.25, 2.0, 2.75, 3.5} {
    \draw[->, thick] (-0.9,\yy) -- (-0.1,\yy);
  }
  \node[font=\footnotesize, align=center] at (-1.0,4.4) {free stream\\$\vect{v}=(U,0)$};

  % Downstream wake outflow arrows (right face), shorter near centreline
  \draw[->, thick] (7.1,0.5) -- (7.9,0.5);
  \draw[->, thick] (7.1,1.25) -- (7.7,1.25);
  \draw[->, thick] (7.1,2.0) -- (7.45,2.0);
  \draw[->, thick] (7.1,2.75) -- (7.7,2.75);
  \draw[->, thick] (7.1,3.5) -- (7.9,3.5);
  \node[font=\footnotesize, align=center] at (8.1,4.4) {wake\\$u(y)<U$};

  % Wake profile sketch on the right face
  \draw[engred, thick] plot[smooth, domain=0:4, samples=40]
    ({7.6 - 0.9*exp(-((\x-2.0)*(\x-2.0))/0.5)}, \x);
  \node[font=\footnotesize, engred, anchor=west] at (6.0,2.0) {$u(y)$};

  % Lateral outflow arrows (top and bottom faces)
  \draw[->, thick, engblue] (5.0,4.0) -- (5.0,4.7);
  \draw[->, thick, engblue] (5.0,0.0) -- (5.0,-0.7);
  \node[font=\footnotesize, engblue, anchor=west] at (5.1,4.6) {lateral mass flux};

  % Control surface label
  \node[font=\footnotesize, engblue, anchor=south west] at (0.05,0.05) {$\partial V$};
  \node[font=\footnotesize] at (3.5,-0.4) {$\divg\vect{v}=0$ over $V$};
\end{tikzpicture}

  \caption[Control volume around a body: momentum-deficit wake.]{A
  fixed rectangular control volume encloses a body in a steady stream.
  The upstream face (left) carries the uniform free stream $U$; the
  downstream face (right) carries the slowed wake profile $u(y)$. The
  momentum deficit between the two faces, plus the lateral mass flux
  through the top and bottom, equals the drag the body exerts on the
  fluid, which by reaction is the drag the fluid exerts on the body.
  The same surface integral that the divergence theorem uses for mass
  here carries momentum. Source: authors' diagram.}
  \label{fig:vol02:ch08:control-volume-drag}
\end{figure}

\paragraph{Step 2: the momentum flux gives the drag.} The drag is the
net rate at which the fluid loses streamwise momentum, which is the
flux of $x$-momentum through the control surface. The $x$-momentum
flux density is $\rho u\,(\vect{v}\cdot\unitvec{n})$, the product of
the streamwise velocity and the volume flux through each surface
element. Integrating over the closed surface and invoking the steady
momentum balance (the net momentum outflow equals the net force on the
fluid), the drag force the body exerts on the fluid is
\[
D = \rho\oint_{\partial V} u\,(\vect{v}\cdot\unitvec{n})\,\dd A =
\rho\int_{\text{wake}} u(y)\big(u(y) - U\big)\,\dd y,
\]
after the upstream, top, and bottom faces are accounted for by the
mass balance of Step~1. The integrand $u(u - U)$ is negative wherever
the wake is slowed ($u < U$), so the body removes momentum from the
fluid and the fluid pushes back on the body with an equal and opposite
drag, the working content of the \engterm{momentum-deficit} method
\cite{text:white-fluid-mechanics}. The method reads the drag from a
single downstream velocity traverse, with no pressure measurement on
the body and no integration over the body's own surface.

\paragraph{Step 3: a worked drag number.} A circular cylinder of
diameter $d = 0.05\,\si{\meter}$ sits in a water tunnel with free
stream $U = 1.0\,\si{\meter\per\second}$ and density $\rho = 1000\,
\si{\kilogram\per\meter\cubed}$. A downstream traverse finds a wake
that is well approximated, over its width $w \approx 2d$, by the
profile $u(y) = U\big(1 - \tfrac{1}{2}\mathrm{e}^{-(2y/d)^{2}}\big)$,
a Gaussian momentum deficit with a centreline velocity of $U/2$. The
drag per unit span is
\[
D' = \rho\int_{-\infty}^{\infty}u(u - U)\,\dd y = \rho U^{2}
\int_{-\infty}^{\infty}\Big(1 - \tfrac{1}{2}\mathrm{e}^{-(2y/d)^{2}}
\Big)\Big(-\tfrac{1}{2}\mathrm{e}^{-(2y/d)^{2}}\Big)\,\dd y.
\]
The integral evaluates with $\int_{-\infty}^{\infty}\mathrm{e}^{-(2y/
d)^{2}}\dd y = \tfrac{d}{2}\sqrt{\pi}$ and $\int_{-\infty}^{\infty}
\mathrm{e}^{-2(2y/d)^{2}}\dd y = \tfrac{d}{2}\sqrt{\pi/2}$, giving
\[
D' = \rho U^{2}\Big(-\tfrac{1}{2}\cdot\tfrac{d\sqrt{\pi}}{2} +
\tfrac{1}{4}\cdot\tfrac{d}{2}\sqrt{\tfrac{\pi}{2}}\Big) = \rho U^{2}d
\Big(-\tfrac{\sqrt{\pi}}{4} + \tfrac{1}{8}\sqrt{\tfrac{\pi}{2}}\Big)
\approx -0.32\,\rho U^{2}d.
\]
The magnitude is $\abs{D'} = 0.32\rho U^{2}d = 0.32\times 1000\times
1.0^{2}\times 0.05 = 16\,\si{\newton\per\meter}$. Expressed as a drag
coefficient $C_{D} = \abs{D'}/(\tfrac{1}{2}\rho U^{2}d) = 0.32/0.5 =
0.64$, this sits within the measured band for a circular cylinder at
this Reynolds number ($Re = Ud/\nu \approx 5\times 10^{4}$, where the
measured $C_{D} \approx 1.0$ to $1.2$, current as of 2026); the
factor-of-two gap is the model wake profile's coarseness, not a fault
of the momentum method \cite{text:white-fluid-mechanics}. A reader who
measures the actual wake profile and integrates it numerically closes
the gap, and the structure of the calculation, drag as a momentum-flux
integral over a control surface, is unchanged.

\begin{mastery}
\textit{The Reynolds transport theorem.} The momentum balance of
Step~2 assumed the control volume was fixed and the flow steady. The
general bookkeeping that relates the rate of change of a quantity for a
moving parcel of fluid to the rate of change over a fixed control
volume is the \engterm{Reynolds transport theorem}, the divergence
theorem applied to an extensive quantity carried by a flow. For any
field $\phi$ per unit volume,
\[
\dv{}{t}\iiint_{V(t)}\phi\,\dd V = \iiint_{V}\pdv{\phi}{t}\,\dd V +
\oint_{\partial V}\phi\,(\vect{v}\cdot\unitvec{n})\,\dd A,
\]
where the left side follows a material volume $V(t)$ that moves with
the flow and the right side is evaluated over the fixed volume $V$ that
coincides with it at the instant. The first right-hand term is the
local accumulation; the second is the convective flux through the
boundary, the same surface integral the divergence theorem converts to
$\iiint\divg(\phi\vect{v})\,\dd V$. Setting $\phi = \rho$ recovers the
continuity equation; setting $\phi = \rho\vect{v}$ recovers the
momentum equation whose steady form gave the drag above; setting
$\phi = \rho e$ for energy density $e$ recovers the energy equation of
Volume~V. The Reynolds transport theorem is the material derivative of
section~8.6 written for an extensive quantity over a volume rather than
for an intensive quantity at a point, and the two are the same chain
rule applied at two scales \cite{text:white-fluid-mechanics}.
\end{mastery}

\paragraph{Step 4: circulation gives the lift.} The lift comes from
the circulation, the line integral of velocity around a loop enclosing
the body,
\[
\Gamma = \oint_{C}\vect{v}\cdot\dd\vect{r},
\]
which by Stokes's theorem equals the flux of vorticity
$\curl\vect{v}$ through any surface the loop bounds. For a lifting
body the circulation is non-zero, and the Kutta-Joukowski theorem,
introduced in the estimation block above, gives the lift per unit
span as $L' = \rho U\,\Gamma$. Take a thin airfoil at a small angle
of attack with a measured circulation $\Gamma = 6.0\,\si{\meter
\squared\per\second}$ in the same water stream ($U = 1.0\,\si{\meter
\per\second}$, $\rho = 1000\,\si{\kilogram\per\meter\cubed}$). The
lift per unit span is
\[
L' = \rho U\Gamma = 1000\times 1.0\times 6.0 = 6.0\times 10^{3}\,
\si{\newton\per\meter}.
\]
The circulation is the same loop integral the chapter has computed for
abstract fields, now closed around a physical wing section; Stokes's
theorem guarantees that the loop may be drawn at any distance from the
body, since the vorticity outside the thin boundary layer and wake is
zero and the circulation through any enclosing loop is the same. A
reader who measures the velocity around a loop with a set of probes,
sums $\vect{v}\cdot\dd\vect{r}$ around it, and multiplies by $\rho U$
has measured the lift without ever touching the wing surface.

\paragraph{Step 5: the divergence-free check and the failure margin.}
Both forces rested on the incompressible assumption $\divg\vect{v} =
0$, which the reader checks on the measured field before trusting the
balances. A wake traverse that returns a net mass imbalance through
the control surface signals either a compressibility effect (the Mach
number is not small), a three-dimensional flow leaking momentum out of
the measurement plane, or a probe-calibration error. The check is the
divergence theorem run in reverse: integrate $\vect{v}\cdot\unitvec{n}$
over the closed control surface and confirm it returns zero within the
traverse uncertainty. For the water-tunnel cylinder the free-stream
Mach number is effectively zero and the assumption is exact; for a
transonic aircraft wing the same momentum method requires the
compressible momentum flux $\rho u(\vect{v}\cdot\unitvec{n})$ with a
varying $\rho$, and the incompressible shortcut fails. The failure
margin is therefore the Mach number: the control-volume drag method in
the form above is trustworthy below roughly $\mathrm{Ma} = 0.3$, where
density variations stay below a few percent, and requires the
compressible form above it \cite{text:anderson-aerodynamics}.

The case study used the flux integral for the momentum balance, the
divergence theorem for the mass balance and the divergence-free check,
the circulation line integral for the lift, and Stokes's theorem to
free the lift loop from the body's surface. Two forces that a reader
might expect to require a pressure integral over the body's own
surface instead fell out of surface integrals over a control volume
drawn in the undisturbed flow, which is the working economy the
integral theorems buy: the force on a body is read from the fluid far
from it, where the flow is simple, rather than on the body, where it is
not. Volumes V and IX develop the fluid and aerodynamic devices in
full; the bookkeeping is the vector calculus of this chapter.

\section{Failure: orientation errors and the sign-of-the-normal
problem}\pathtag{core}

Vector calculus has one failure mode that recurs in working
engineering more often than any other: the orientation error. The
sign of a flux integral, the direction of a torque, the chirality of
a coordinate frame, and the inward-versus-outward convention on a
closed surface all depend on a binary choice. A working engineer who
gets the sign wrong has produced a calculation that looks correct in
its magnitudes and is reversed in its direction; a working engineer
who never checks the sign explicitly is one orientation error away
from a sign-reversal in a load, a flow, or a control output.

\subsection{The mechanism}

The orientation error has three common forms.

\engterm{The hand-rule swap}. Cross products and curl results depend
on the right-hand rule. A reader who applies the left-hand rule, or
who applies the right-hand rule with a left-handed coordinate frame,
produces a vector that points in the opposite direction. The error
propagates: a torque computed with the wrong hand rule reverses the
direction of the angular acceleration; a magnetic-field direction
computed with the wrong hand rule reverses the force on a current.
Some engineering communities use left-handed coordinate frames as a
local convention (the screen-coordinate frame in many graphics
systems is left-handed); a working engineer crossing between
communities checks the convention before trusting the formula.

\engterm{The parametrisation flip}. A surface parametrisation
$\vect{r}(u, v)$ implies a normal direction $\unitvec{n} =
(\vect{r}_{u} \times \vect{r}_{v}) / \norm{\vect{r}_{u} \times
\vect{r}_{v}}$. Swapping the roles of $u$ and $v$, or reversing the
orientation of either parameter, reverses $\unitvec{n}$. A flux
integral computed with the reversed normal is the negative of the
intended one. The reader who chooses a parametrisation without
checking which direction the normal points has implicitly chosen a
sign for every flux integral over the surface.

\engterm{The closed-surface orientation lapse}. The divergence
theorem requires the unit normal on a closed surface to point
outward. A reader who parametrises a closed surface piecewise
(top, bottom, sides) and uses the same parametrisation rule for
each piece often discovers, midway through the surface integral,
that one piece's normal points inward. The flux contributions of
the pieces with reversed normals are subtracted instead of added,
and the result is wrong by an unknown amount that depends on the
proportion of misoriented pieces.

\engterm{The coordinate-singularity mistake}. Curvilinear coordinate
systems carry singular points where the basis vectors are undefined or
the Jacobian degenerates: the polar axis in cylindrical coordinates
(where $\unitvec{\rho}$ and $\unitvec{\phi}$ are undefined at
$\rho = 0$), the poles in spherical coordinates (where the
$1/\sin\theta$ factors blow up), and the origin of any inverse-power
field. A reader who applies the divergence theorem to a volume that
contains such a singularity, or who divides by $\rho$ or $\sin\theta$
at the singular point, produces a result that is finite in the
arithmetic and wrong in the physics. The inverse-square field of the
worked example above is the canonical case: its divergence is zero
everywhere the formula is valid, yet its flux through any enclosing
sphere is $4\pi$, because the source hides at the coordinate
singularity the smooth formula cannot see. The discipline is to
excise a small ball around every singularity, apply the theorem to the
punctured region where the field is smooth, and account for the
excised flux separately.

\subsection{The operational test}

The discipline that catches all three forms is the same: an explicit
operational test of the orientation, performed before any
quantitative calculation, and rerun whenever the parametrisation
changes.

For a closed surface, the test is a sample point. Choose a point on
the surface, evaluate $\unitvec{n}$ at that point, and check by
inspection whether the resulting vector points away from the
enclosed volume. If it does, the orientation is outward. If it does
not, swap the parametrisation order or negate $\unitvec{n}$ before
proceeding.

For a curl-and-circulation calculation, the test is a sample
direction. Trace the boundary curve in the parametrised direction;
align the right hand so the fingers curl with the boundary. The
thumb's direction is the surface normal as Stokes's theorem demands
it. If the chosen $\unitvec{n}$ disagrees, reverse the boundary
parametrisation or negate the normal.

For a torque or angular-velocity calculation, the test is a sample
load. Apply a positive force along $\vect{F}$ at position $\vect{r}$;
imagine which way the system actually spins. The torque vector
$\vect{r} \times \vect{F}$ should point along the axis whose
right-hand rotation matches the imagined spin. If the formal
computation produces the opposite axis, the cross-product order or
the sign convention is wrong.

\subsection{The rule the engineer applies}

\begin{principle}[Verify orientation explicitly]
Every flux integral, surface integral, line integral, cross product,
curl, and torque carries an implicit sign that depends on a binary
orientation choice. The engineer states the choice explicitly,
verifies it against an operational sample at the time the choice is
made, and rechecks it whenever the parametrisation, the basis, or the
hand rule changes. A calculation whose magnitudes are correct and
whose sign is wrong has the same effect as a calculation that is
wrong by a factor of two: the error is loud, traceable, and
unforgivable in a system that depends on the sign for safety,
direction, or stability.
\end{principle}

The principle has a corollary. A reader who is asked to inherit a
calculation from a previous engineer should rerun the orientation
test before extending the calculation. The previous engineer's
parametrisation is the previous engineer's responsibility; the
calculation about to be performed is the present engineer's
responsibility, and the orientation discipline does not transfer
across the inheritance. Volumes III, IV, and IX take up specific
cases of inherited-orientation failure in mechanics, electromagnetism,
and aerospace; the rule above is the general one.

\subsection{A worked example of the failure}

A reader computes the flux of $\vect{F} = (0, 0, z)$ outward through
a closed surface composed of the unit cube's six faces, by
parametrising each face as a function of two of the Cartesian
coordinates with a fixed value of the third. On the top face,
$z = 1$, the reader writes $\vect{r}(x, y) = (x, y, 1)$ for $(x, y)
\in [0, 1]^{2}$. The induced normal is $\vect{r}_{x} \times
\vect{r}_{y} = \unitvec{x} \times \unitvec{y} = +\unitvec{z}$, which
is outward. On the bottom face, $z = 0$, the reader writes
$\vect{r}(x, y) = (x, y, 0)$ with the same parametrisation rule. The
induced normal is again $+\unitvec{z}$, but on the bottom face
$+\unitvec{z}$ points inward, not outward.

\begin{figure}[h]
  \centering
  % The unit-cube orientation trap from section 8.8.
% Shows the cube; on the bottom face the induced normal points inward (red);
% on the top face it points outward (blue, correct).
\begin{tikzpicture}[>=Stealth, scale=1.4]
  % cube vertices
  \coordinate (A) at (0,0);
  \coordinate (B) at (2,0);
  \coordinate (C) at (2.7,0.7);
  \coordinate (D) at (0.7,0.7);
  \coordinate (E) at (0,2);
  \coordinate (F) at (2,2);
  \coordinate (G) at (2.7,2.7);
  \coordinate (H) at (0.7,2.7);

  % back faces (dashed)
  \draw[dashed, enggray] (A)--(D)--(C);
  \draw[dashed, enggray] (D)--(H);

  % front + top + right faces
  \draw[thick] (A)--(B)--(F)--(E)--cycle;
  \draw[thick] (B)--(C)--(G)--(F);
  \draw[thick] (E)--(H)--(G);

  % top-face outward normal (correct, blue, up)
  \draw[->, very thick, engblue]
    ($(E)!0.5!(G)$) -- ++(0,0.9) node[above, font=\footnotesize]{$+\unitvec{z}$ outward (top, correct)};

  % bottom-face induced normal (wrong, red, also up but pointing into cube)
  \draw[->, very thick, engred]
    ($(A)!0.5!(C)$) -- ++(0,0.9) node[above right, font=\footnotesize, align=left]
    {$+\unitvec{z}$ induced (bottom)\\points \emph{into} the cube};

  % labels
  \node[font=\footnotesize] at ($(A)!0.5!(B)+(0,-0.18)$) {bottom face $z=0$};
  \node[font=\footnotesize] at ($(H)!0.5!(G)+(0,0.18)$) {top face $z=1$};
\end{tikzpicture}

  \caption[Unit-cube orientation trap.]{On the top face the induced
  normal $+\unitvec{z}$ points outward, as the divergence theorem
  requires (blue). On the bottom face the same parametrisation rule
  produces an induced normal $+\unitvec{z}$ that points inward,
  reversing the sign of every flux contribution from that face (red).
  The reader either reverses the bottom-face parametrisation or
  negates the contribution. Source: authors' diagram.}
  \label{fig:vol02:ch08:cube-trap}
\end{figure}

A reader who carries this orientation through the surface integral
adds the bottom face's flux $\iint_{\text{bottom}} \vect{F} \cdot
(+\unitvec{z}) \,\dd A = 0$ instead of subtracting it. In this
specific case, $\vect{F} = (0, 0, z)$ vanishes on the bottom face,
and the error of sign produces no error in magnitude. A different
field, $\vect{F} = (0, 0, 1)$, exposes the failure: each side face
contributes zero, the top contributes $+1$ (correctly), and the
bottom contributes $+1$ in the reader's convention versus the
correct $-1$. The reader's result is $+2$; the correct result is $0$.
The divergence theorem confirms: $\divg\vect{F} = 0$, so
$\iiint_{V} \divg\vect{F} \,\dd V = 0$.

The discipline this example asks of the reader is mechanical. Before
the surface integral is computed, evaluate $\unitvec{n}$ on each face
at a sample point and confirm it points outward. The five-second
check at parametrisation time prevents the half-hour debugging hunt
at result-checking time, and prevents the much longer professional
embarrassment when a result reaches a colleague before the sign
error is caught.

\subsection{A worked example of the coordinate-singularity trap}

A reader computes the outward flux of the radial field $\vect{F} =
\vect{r}/r^{3}$ through a cube centred at the origin, intending to use
the divergence theorem to convert the awkward surface integral into a
volume integral. The reader computes $\divg\vect{F} = 0$ for $r > 0$,
concludes $\iiint_{V}\divg\vect{F}\,\dd V = 0$, and reports zero net
flux. The report is wrong: the true flux is $4\pi$, the same value the
sphere calculation produced, because the cube also encloses the
origin and the flux through any closed surface enclosing the source is
$4\pi$ regardless of shape.

The mechanism is the singularity at $r = 0$, where $\divg\vect{F}$ is
not zero but a concentrated source. The smooth formula $\divg\vect{F}
= 0$ holds only away from the origin; the divergence theorem in the
form used requires the field to be smooth throughout the enclosed
volume, and the cube violates that requirement at its centre. The
operational test that catches the error is a sanity check the reader
should run before trusting any divergence-theorem result: confirm the
field is finite and differentiable at every interior point of the
volume. The radial field fails the check at the origin in one second
of inspection, which is cheaper than the surface integral the reader
was trying to avoid.

The corrective procedure is the excision discipline named above.
Surround the origin with a small ball $B_{\epsilon}$ of radius
$\epsilon$, apply the theorem to the region between the cube and the
ball where $\vect{F}$ is smooth (which gives zero, since
$\divg\vect{F} = 0$ there), and conclude that the outward cube flux
equals the outward flux through the small sphere, which is $4\pi$ for
any $\epsilon$. The same excision turns the apparent contradiction
into Gauss's law for a point charge, the working tool of Volume~IV.
The reader who omits the smoothness check inherits the contradiction;
the reader who runs it inherits the correct $4\pi$.

\section{Project}\pathtag{core}

\begin{figure}[p]
  \centering
  % Full-page schematic of the divergence theorem on a triangulated household object.
% Three panels: (a) the object with sampled cross-sections; (b) the closed
% triangulation; (c) a single oriented triangle contributing one signed tetrahedron.
\begin{tikzpicture}[>=Stealth, scale=1.0]

  % --- Panel (a): object with horizontal cross-sections ---
  \begin{scope}[xshift=0cm, yshift=0cm]
    \node[font=\footnotesize, anchor=south] at (0,4.0) {(a) Sampled cross-sections};
    \draw[fill=enggray!20, draw=enggray] (-1.2,0) .. controls (-1.6,1.0) and (-1.4,2.0) ..
      (-0.8,3.0) .. controls (-0.3,3.6) and (0.4,3.5) ..
      (1.0,3.0) .. controls (1.5,2.2) and (1.3,1.0) ..
      (0.9,0) -- cycle;
    \foreach \z in {0.4, 1.0, 1.7, 2.4, 3.0} {
      \draw[engblue, thin] (-1.5,\z) -- (1.4,\z);
      \foreach \fact in {-1,-0.5,0,0.5,1} {
        \fill[engblue] (\fact*1.0, \z) circle (0.03);
      }
    }
    \node[font=\footnotesize, anchor=north] at (0,-0.2) {30+ surface points};
  \end{scope}

  % --- Panel (b): closed triangulation ---
  \begin{scope}[xshift=5.2cm, yshift=0cm]
    \node[font=\footnotesize, anchor=south] at (0,4.0) {(b) Closed triangulation};
    \coordinate (P1) at (-1.0, 0.4);
    \coordinate (P2) at (0.0, 0.0);
    \coordinate (P3) at (0.9, 0.5);
    \coordinate (P4) at (-1.1, 1.2);
    \coordinate (P5) at (-0.1, 1.0);
    \coordinate (P6) at (1.0, 1.3);
    \coordinate (P7) at (-1.0, 2.0);
    \coordinate (P8) at (0.0, 1.9);
    \coordinate (P9) at (0.9, 2.1);
    \coordinate (PT) at (-0.1, 3.4);
    \draw[engblue] (P1)--(P2)--(P3) (P4)--(P5)--(P6) (P7)--(P8)--(P9);
    \draw[engblue] (P1)--(P4)--(P7)--(PT)--(P9)--(P6)--(P3) (P2)--(P5)--(P8)--(PT);
    \draw[engblue] (P1)--(P5) (P2)--(P6) (P4)--(P8) (P5)--(P9);
    \foreach \p in {P1,P2,P3,P4,P5,P6,P7,P8,P9,PT} {
      \fill[engblue] (\p) circle (0.04);
    }
    \node[font=\footnotesize, anchor=north] at (0,-0.2) {consistent outward normals};
  \end{scope}

  % --- Panel (c): single oriented triangle and its tetrahedron ---
  \begin{scope}[xshift=10.6cm, yshift=0.0cm]
    \node[font=\footnotesize, anchor=south] at (0,4.0) {(c) Per-triangle signed volume};
    \coordinate (O) at (0,0);
    \coordinate (Q1) at (-0.6, 1.5);
    \coordinate (Q2) at (0.8, 1.2);
    \coordinate (Q3) at (0.2, 2.6);
    \fill[englime!30] (Q1)--(Q2)--(Q3)--cycle;
    \draw[thick] (Q1)--(Q2)--(Q3)--cycle;
    \draw[dashed, enggray] (O)--(Q1) (O)--(Q2) (O)--(Q3);
    \fill (O) circle (0.05);
    \node[font=\footnotesize, anchor=east] at (-0.05,0) {origin};
    \node[font=\footnotesize] at (-0.85,1.55) {$\vect{p}_1$};
    \node[font=\footnotesize] at (1.05,1.2) {$\vect{p}_2$};
    \node[font=\footnotesize] at (0.45,2.7) {$\vect{p}_3$};
    \draw[->, very thick, engred] (0.15,1.75) -- ++(0.7,0.4)
      node[right, font=\footnotesize]{$\vect{p}_2\times\vect{p}_3$};
    \node[font=\footnotesize, anchor=north, align=center] at (0,-0.2)
      {$V_{\triangle} = \tfrac{1}{6}\,\vect{p}_1\!\cdot\!(\vect{p}_2\!\times\!\vect{p}_3)$};
  \end{scope}

\end{tikzpicture}

  \caption[Project workflow: surface sample, triangulation, per-triangle
  signed volume.]{The three working steps of the chapter project. Panel
  (a) shows the object with thirty-plus surface samples drawn from
  several horizontal cross-sections. Panel (b) shows the closed
  triangulation built from the samples, with consistent outward
  normals. Panel (c) shows a single oriented triangle and the signed
  tetrahedron $V_{\triangle} = \tfrac{1}{6}\vect{p}_{1}\cdot
  (\vect{p}_{2}\times\vect{p}_{3})$ whose sum over all triangles is
  the divergence-theorem volume of the object. The companion code at
  \repopath{volumes/02-form/ch08-vectors-and-vector-calculus/code/}
  implements the per-triangle sum and reports the convergence sweep
  whose tabulated reference appears in
  \repopath{data/sweep_reference.csv}. Source: authors' diagram,
  algorithm after \textcite{paper:v2ch08-zhang-chen-2001}.}
  \label{fig:vol02:ch08:divergence-theorem}
\end{figure}

\begin{project}[The volume of an irregular object, by the divergence
theorem]
The reader takes one irregular household object whose volume the
reader already measured by water displacement in the Volume~I
Chapter~7 project, and computes that volume a second way: by the
divergence theorem applied to a sampled triangulation of the object's
surface. The two values are compared. Where they agree within their
combined uncertainty, the divergence theorem and the empirical
measurement reinforce each other; where they disagree, the reader
diagnoses whether the surface sample is too coarse, the divergence
formula is being applied with a sign error, or the water-displacement
measurement was biased. The chapter has just finished arguing that
balance is a theorem; the project is the reader's first occasion to
exercise that theorem on a measured object.

\textbf{Track}. Household and analysis (hybrid). The reader works at
the kitchen table or workbench. \textbf{Hazard class}: none. The
chosen object should not be sharp, electrical, valuable, or
water-damaged. \textbf{Tools}. The same household tools as
Volume~I Chapter~7 (kitchen scale, ruler, beaker or graduated
cylinder), plus paper and pencil for the cross-section sampling and
either a programming environment (Python, MATLAB, Julia, or
equivalent) or a spreadsheet for the surface-mesh triangulation and
the divergence-formula calculation. A 3D scanner or smartphone
photogrammetry app is optional and useful but not required.

\textbf{Choosing the object}. Reuse the irregular object measured in
Volume~I Chapter~7 if possible: the comparison against the original
water-displacement value is the project's reconciliation step. If
the original object is no longer available, choose a comparable
object (a smooth stone, a small piece of fruit, a small hand tool)
and rerun the water-displacement measurement first.

\textbf{Procedure}. Apply the formula
\[
V = \frac{1}{3} \oint_{\partial V} \vect{r} \cdot \unitvec{n} \,\dd A
\]
to a sampled approximation of the object's surface, in five steps.

\begin{enumerate}[leftmargin=*]
  \item \engterm{Establish a coordinate frame}. Place the object on
    a sheet of paper in a stable orientation. Mark a chosen origin
    on the paper directly under a chosen reference point on the
    object (a stable resting point or a marked feature). Define the
    $z$-axis pointing up; choose two perpendicular axes in the paper
    plane and label them $x$ and $y$. Verify the frame is right-
    handed: $\unitvec{x} \times \unitvec{y}$ should point along
    $+\unitvec{z}$.

  \item \engterm{Sample at least 30 surface points}. Use a calliper,
    ruler, or other measuring tool to record the $(x, y, z)$
    coordinates of at least 30 points distributed across the
    object's surface. Distribute the points so that no part of the
    surface is unsampled by more than approximately one-third of the
    object's smallest characteristic dimension. The points should be
    organised either by horizontal cross-sections at several heights
    or by latitude-and-longitude rings around the object. Record the
    sampling resolution as part of the measurement record.

  \item \engterm{Build a closed triangulated surface}. Connect the
    sampled points into a closed triangle mesh that approximates the
    object's surface. The simplest construction stacks the
    horizontal cross-sections, connects each to the next by a band
    of triangles, and caps the top and bottom with fan triangulations
    centred on the highest and lowest sampled points. Each triangle
    is specified by an ordered triple of vertex labels; the order
    determines the triangle's normal direction. Use the same
    orientation convention for every triangle (for instance,
    counterclockwise as seen from outside the object), so the
    resulting mesh has consistently outward-pointing normals.
    Verify the convention by sampling one triangle and checking by
    inspection that the right-hand rule applied to its vertex order
    gives a normal that points away from the enclosed volume.

  \item \engterm{Compute the divergence-theorem volume}. For each
    triangle with ordered vertices $\vect{p}_{1}, \vect{p}_{2},
    \vect{p}_{3}$, the contribution to $V$ is
    \[
    V_{\triangle} = \frac{1}{6}
    \vect{p}_{1} \cdot (\vect{p}_{2} \times \vect{p}_{3}).
    \]
    This is the signed volume of the tetrahedron with one vertex at
    the origin and the opposite face the triangle. The total volume
    is the sum of these contributions over all triangles:
    $V = \sum_{\triangle} V_{\triangle}$. The derivation is in the
    chapter's worked example above; the formula follows from the
    divergence-theorem expression $V = \tfrac{1}{3} \oint
    \vect{r} \cdot \unitvec{n} \,\dd A$ specialised to a flat
    triangle, where the triangle's contribution to the integral is
    $\tfrac{1}{3}\vect{r}_{\text{centroid}} \cdot \unitvec{n} \,
    A_{\triangle}$ and a short calculation reduces it to the
    one-sixth scalar triple product above. Compute the sum.

  \item \engterm{Reconcile against water displacement}. Compare the
    divergence-theorem volume $V_{\text{DT}}$ to the
    water-displacement volume $V_{\text{WD}}$ from Volume~I
    Chapter~7, with both reported with their standard uncertainties.
    Compute the compatibility ratio
    \[
    r = \frac{\abs{V_{\text{DT}} - V_{\text{WD}}}}
             {\sqrt{u(V_{\text{DT}})^{2} + u(V_{\text{WD}})^{2}}}
    \]
    using the chapter~7 cross-checking discipline. A value of $r$
    near $1$ or below means the two methods agree within their
    combined uncertainty.
\end{enumerate}

\textbf{Uncertainty budget for the divergence-theorem volume}. The
divergence-theorem volume inherits uncertainty from three sources.
The first is the position uncertainty on each sampled point, which
propagates through the scalar triple product to a per-triangle
volume uncertainty; the second is the truncation error of replacing
a smooth surface by a piecewise-flat one, which decreases as the
sample is refined; the third is the orientation discipline, which
either contributes zero error if the convention is consistent or
contributes a sign-flipped triangle's worth of error if it is not.
The reader reports each source separately and combines them in
quadrature for an expanded uncertainty at $k = 2$.

\textbf{Reflection (1500 words)}. The reader writes a 1500-word
narrative addressing:

\begin{itemize}[leftmargin=*]
  \item How the divergence-theorem volume compared to the
    Volume~I water-displacement volume, both in magnitude and in
    the compatibility ratio.
  \item Which of the three uncertainty sources dominated the
    divergence-theorem result, and what step would most reduce the
    dominant source.
  \item Whether refining the surface sampling from $30$ points
    to a finer grid would close any remaining gap, or whether the
    remaining disagreement points to a different error source.
  \item One physical engineering question (mass of a casting,
    drag on a body, heat loss from a surface) where the same
    surface-to-volume bookkeeping would be useful, and whether the
    divergence-theorem method would be the working tool or a
    pedagogical aid.
  \item What the divergence theorem buys over the empirical
    water-displacement method: in particular, where the theorem
    earns its place as a working tool and where water displacement
    is the cheaper and more defensible measurement.
\end{itemize}

\textbf{Deliverable}. The cross-section data, the triangulation, the
per-triangle volume contributions, the total divergence-theorem
volume with its uncertainty budget, the compatibility ratio against
Volume~I, and the reflection. The general editor's reference set of
divergence-theorem computations for a representative stone is
published online; the reader is encouraged to compare \emph{after}
completing their own.

\textbf{Time}. Two to four hours for sampling and triangulation, one
to two hours for the divergence calculation and uncertainty budget,
two to three hours for the reflection. Total approximately five to
nine hours, distributed across one or two days.

\textbf{Companion code and data}. A reference implementation of the
per-triangle signed-volume sum
$V = \tfrac{1}{6}\sum_{\triangle}\vect{p}_{1}\cdot(\vect{p}_{2}\times
\vect{p}_{3})$ lives at
\repopath{volumes/02-form/ch08-vectors-and-vector-calculus/code/mesh_volume.py}
and accepts a minimal Wavefront ``.obj'' file. The reader who reaches
the divergence-calculation step with the triangulation in hand
exports the mesh in that format (or in a spreadsheet of vertex
triples and triangle index triples) and runs the script as
\texttt{uv run mesh\_volume.py mymesh.obj}. A convergence sweep on
three reference shapes (unit cube, regular tetrahedron, geodesic
icosphere at three refinement levels) is in
\repopath{code/triple_product_sweep.py}; its tabulated output sits
in \repopath{data/sweep_reference.csv} so a reader can confirm a
clean install before applying the script to the household object.
The script prints a warning on a negative signed volume, which is
the test for the uniformly inverted orientation discussed in
section~8.8.

\textbf{Phases and success criteria}. The project succeeds when the
reader has produced (i)~a triangulation whose every triangle has been
checked for outward orientation against a sample point; (ii)~a
divergence-theorem volume $V_{\text{DT}}$ with a stated uncertainty
budget; (iii)~a compatibility ratio against the Volume~I
water-displacement volume below approximately two; and (iv)~a
1500-word reflection that names which uncertainty source dominated
and what step would most reduce it. A reader who reports
$V_{\text{DT}}$ without a sign check and without an uncertainty
budget has not completed the project, even if the number happens to
agree with the water-displacement result. The discipline being
exercised is the explicit verification step, not the arithmetic.

\textbf{Common failure modes}. Three recur. First, the bottom-cap
triangulation inherits the same vertex-order rule as the top cap,
which inverts its normal; a reader who has not run the operational
test on at least one bottom-cap triangle will see a divergence-theorem
volume that is too small by twice the bottom cap's contribution.
Second, a sampling grid that is too coarse on a high-curvature region
(a handle, a spout, a sharp edge) returns a volume biased downward
because the piecewise-flat triangulation under-represents the surface
area available to the integrand $\vect{r}\cdot\unitvec{n}$. Third, the
Volume~I water-displacement value carries its own systematic biases
(meniscus reading, surface bubbles, beaker calibration); a reader
who treats $V_{\text{WD}}$ as ground truth and assumes any
disagreement is the divergence-theorem method's fault has inverted
the uncertainty discipline of Volume~I Chapter~7.
\end{project}

\section{Exercises}

The exercises test calculation, derivation, estimation, simulation,
design, diagnosis, failure analysis, and judgment. The editor
should attach solutions and rubrics before release.

\subsection*{Calculation}\pathtag{core}

\begin{exercise}
Two vectors are $\vect{u} = (3, -1, 2)$ and $\vect{v} = (1, 4, -2)$.
Compute $\vect{u} + \vect{v}$, $\vect{u} - \vect{v}$,
$\norm{\vect{u}}$, $\norm{\vect{v}}$, and the unit vector
$\unitvec{u}$.

\paragraph{Solution sketch.} Componentwise, $\vect{u} + \vect{v} =
(4, 3, 0)$ and $\vect{u} - \vect{v} = (2, -5, 4)$. The norms are
$\norm{\vect{u}} = \sqrt{9 + 1 + 4} = \sqrt{14} \approx 3.74$ and
$\norm{\vect{v}} = \sqrt{1 + 16 + 4} = \sqrt{21} \approx 4.58$. The
unit vector is $\unitvec{u} = (3, -1, 2)/\sqrt{14} \approx
(0.802, -0.267, 0.535)$. Check: $\unitvec{u}\cdot\unitvec{u} = 1$
to three significant figures.
\end{exercise}

\begin{exercise}
For the vectors of the previous exercise, compute $\vect{u} \cdot
\vect{v}$, $\vect{u} \times \vect{v}$, the angle $\theta$ between
them, and the area of the parallelogram they span.

\paragraph{Solution sketch.} The dot product is $\vect{u}\cdot
\vect{v} = 3 - 4 - 4 = -5$. The cross product is
$\vect{u}\times\vect{v} = ((-1)(-2) - (2)(4),\ (2)(1) - (3)(-2),\
(3)(4) - (-1)(1)) = (-6, 8, 13)$. The angle satisfies $\cos\theta =
\vect{u}\cdot\vect{v}/(\norm{\vect{u}}\norm{\vect{v}}) =
-5/\sqrt{14\cdot 21} = -5/\sqrt{294} \approx -0.292$, so $\theta
\approx 107^{\circ}$ (obtuse, as the negative dot product reports).
The parallelogram area is $\norm{\vect{u}\times\vect{v}} =
\sqrt{36 + 64 + 169} = \sqrt{269} \approx 16.4$. A cross-check:
Lagrange's identity gives $\norm{\vect{u}}^{2}\norm{\vect{v}}^{2} -
(\vect{u}\cdot\vect{v})^{2} = 14\cdot 21 - 25 = 269$, agreeing.
\end{exercise}

\begin{exercise}
Three vectors are $\vect{u} = (1, 2, 0)$, $\vect{v} = (0, 1, 3)$,
$\vect{w} = (2, 0, 1)$. Compute the scalar triple product
$\vect{u} \cdot (\vect{v} \times \vect{w})$ and report the volume
of the parallelepiped they span.

\paragraph{Solution sketch.} Compute $\vect{v}\times\vect{w} =
((1)(1) - (3)(0),\ (3)(2) - (0)(1),\ (0)(0) - (1)(2)) = (1, 6, -2)$.
Then $\vect{u}\cdot(\vect{v}\times\vect{w}) = (1)(1) + (2)(6) +
(0)(-2) = 13$. The parallelepiped volume is the absolute value,
$\abs{13} = 13$. The positive sign reports that $(\vect{u}, \vect{v},
\vect{w})$ is right-handed (same handedness as the Cartesian basis).
\end{exercise}

\begin{exercise}
A point at Cartesian coordinates $(2.0, 2.0, 3.0) \,\si{\meter}$ is
expressed in cylindrical coordinates $(\rho, \phi, z)$ and in
spherical coordinates $(r, \theta, \phi)$. Compute both
representations to three significant figures.

\paragraph{Solution sketch.} Cylindrical: $\rho = \sqrt{2^{2} + 2^{2}}
= 2\sqrt{2} \approx 2.83\,\si{\meter}$, $\phi = \arctan(2/2) =
\pi/4 = 45^{\circ}$, $z = 3.0\,\si{\meter}$. Spherical: $r =
\sqrt{2^{2} + 2^{2} + 3^{2}} = \sqrt{17} \approx 4.12\,\si{\meter}$,
$\theta = \arccos(z/r) = \arccos(3/\sqrt{17}) \approx 43.3^{\circ}$,
$\phi = 45^{\circ}$ (same as cylindrical). Cross-check: $r\sin\theta
= \rho$ to three significant figures, $4.12\cdot 0.686 \approx 2.83$.
\end{exercise}

\begin{exercise}
A scalar field $\Phi(x, y, z) = x^{2} y + y z^{2}$ has a gradient
$\grad\Phi$ at every point. Compute $\grad\Phi$ at the point
$(1, 2, -1)$ and the directional derivative of $\Phi$ at that
point in the direction of the unit vector $\unitvec{u} =
(1, 1, 1)/\sqrt{3}$.

\paragraph{Solution sketch.} The partials are $\partial\Phi/
\partial x = 2 x y$, $\partial\Phi/\partial y = x^{2} + z^{2}$,
$\partial\Phi/\partial z = 2 y z$. At $(1, 2, -1)$, $\grad\Phi =
(4, 2, -4)$. The directional derivative is $\grad\Phi\cdot
\unitvec{u} = (4 + 2 - 4)/\sqrt{3} = 2/\sqrt{3} \approx 1.155$.
The magnitude $\norm{\grad\Phi} = \sqrt{16 + 4 + 16} = 6$, so the
direction of steepest ascent at the point is $(4, 2, -4)/6 =
(2/3, 1/3, -2/3)$; the directional derivative along $\unitvec{u}$
is smaller than the maximum because $\unitvec{u}$ is not aligned
with $\grad\Phi$.
\end{exercise}

\begin{exercise}
A vector field is $\vect{F}(x, y, z) = (xy, y^{2} z, x z^{2})$.
Compute $\divg\vect{F}$ and $\curl\vect{F}$ at the point
$(1, 1, 1)$.

\paragraph{Solution sketch.} The divergence is $\divg\vect{F} =
\partial(xy)/\partial x + \partial(y^{2}z)/\partial y + \partial
(xz^{2})/\partial z = y + 2 y z + 2 x z$; at $(1, 1, 1)$ this is
$1 + 2 + 2 = 5$. The curl components are: $(\curl\vect{F})_{x} =
\partial(xz^{2})/\partial y - \partial(y^{2}z)/\partial z = 0 -
y^{2} = -1$ at the point; $(\curl\vect{F})_{y} = \partial(xy)/
\partial z - \partial(xz^{2})/\partial x = 0 - z^{2} = -1$;
$(\curl\vect{F})_{z} = \partial(y^{2}z)/\partial x - \partial(xy)/
\partial y = 0 - x = -1$. So $\curl\vect{F}(1, 1, 1) = (-1, -1, -1)$.
Cross-check: the magnitudes are order unity, consistent with the
polynomial degree of $\vect{F}$.
\end{exercise}

\begin{exercise}
A curve is parametrised by $\vect{r}(t) = (\cos t, \sin t,
t)$ for $t \in [0, 2\pi]$. Compute the tangent vector
$\vect{r}\,'(t)$, its magnitude, the arc length of the curve, and
the curvature at $t = \pi/2$.

\paragraph{Solution sketch.} The tangent is $\vect{r}\,'(t) =
(-\sin t, \cos t, 1)$ with constant magnitude $\norm{\vect{r}\,'(t)}
= \sqrt{1 + 1} = \sqrt{2}$. The arc length is $s = \int_{0}^{2\pi}
\sqrt{2}\,\dd t = 2\pi\sqrt{2} \approx 8.886$. The second derivative
is $\vect{r}\,''(t) = (-\cos t, -\sin t, 0)$, with $\norm{\vect{r}\,'
\times\vect{r}\,''} = \norm{(\sin t, -\cos t, 1)} = \sqrt{2}$ (the
cross product reduces to $(\cos t \cdot 0 - 1\cdot(-\sin t),\
1\cdot(-\cos t) - (-\sin t)\cdot 0,\ (-\sin t)(-\sin t) - \cos t
\cdot(-\cos t)) = (\sin t, -\cos t, 1)$). Hence $\kappa = \sqrt{2}/
(\sqrt{2})^{3} = 1/2$, constant along the helix, with radius of
curvature $R = 2$. The result is a standard helix consistency check:
a circle of radius $1$ wound with unit pitch has $\kappa = 1/2$.
\end{exercise}

\begin{exercise}
A surface is parametrised by $\vect{r}(u, v) = (u\cos v, u\sin v,
u^{2})$ for $u \in [0, 1]$ and $v \in [0, 2\pi]$. Compute the
tangent vectors $\vect{r}_{u}$, $\vect{r}_{v}$, the normal vector
$\vect{N} = \vect{r}_{u} \times \vect{r}_{v}$ at $(u, v) = (1/2,
\pi/4)$, and the surface area.

\paragraph{Solution sketch.} The tangents are $\vect{r}_{u} =
(\cos v, \sin v, 2 u)$ and $\vect{r}_{v} = (-u\sin v, u\cos v, 0)$.
The cross product is $\vect{r}_{u}\times\vect{r}_{v} = (\sin v\cdot
0 - 2 u\cdot u\cos v,\ 2 u\cdot(-u\sin v) - \cos v\cdot 0,\
\cos v\cdot u\cos v - \sin v\cdot(-u\sin v)) = (-2 u^{2}\cos v,
-2 u^{2}\sin v, u)$. At $(1/2, \pi/4)$, $\vect{N} = (-2\cdot(1/4)
\cdot\sqrt{2}/2,\ -2\cdot(1/4)\cdot\sqrt{2}/2,\ 1/2) = (-\sqrt{2}/4,
-\sqrt{2}/4, 1/2)$. The norm is $\norm{\vect{N}} = u\sqrt{4 u^{2}
+ 1}$, so the surface area is $A = \int_{0}^{2\pi}\int_{0}^{1}
u\sqrt{4 u^{2} + 1}\,\dd u\,\dd v = 2\pi\cdot[(4 u^{2} + 1)^{3/2}
/12]_{0}^{1} = \pi(5\sqrt{5} - 1)/6 \approx 5.33$. The surface is
the paraboloid $z = x^{2} + y^{2}$ between $z = 0$ and $z = 1$.
\end{exercise}

\subsection*{Derivation}\pathtag{standard}

\begin{exercise}
Starting from the componentwise definition of the dot product, derive
the geometric formula $\vect{u} \cdot \vect{v} = \norm{\vect{u}}
\norm{\vect{v}} \cos\theta$ using the law of cosines applied to the
triangle with sides $\vect{u}$, $\vect{v}$, and $\vect{u} - \vect{v}$.

\paragraph{Solution sketch.} The triangle with sides $\vect{u}$,
$\vect{v}$, $\vect{u} - \vect{v}$ has the angle $\theta$ opposite the
side $\vect{u} - \vect{v}$. The law of cosines states $\norm{\vect{u}
- \vect{v}}^{2} = \norm{\vect{u}}^{2} + \norm{\vect{v}}^{2} - 2
\norm{\vect{u}}\norm{\vect{v}}\cos\theta$. Expand the left side
componentwise: $\norm{\vect{u} - \vect{v}}^{2} = \sum_{i}(u_{i} -
v_{i})^{2} = \sum_{i} u_{i}^{2} - 2\sum_{i} u_{i} v_{i} + \sum_{i}
v_{i}^{2} = \norm{\vect{u}}^{2} - 2 \vect{u}\cdot\vect{v} +
\norm{\vect{v}}^{2}$. Equating the two expressions and cancelling
$\norm{\vect{u}}^{2} + \norm{\vect{v}}^{2}$ gives $-2 \vect{u}\cdot
\vect{v} = -2 \norm{\vect{u}}\norm{\vect{v}}\cos\theta$, hence
$\vect{u}\cdot\vect{v} = \norm{\vect{u}}\norm{\vect{v}}\cos\theta$.
The derivation requires only that $\vect{u}$, $\vect{v}$, and
$\vect{u} - \vect{v}$ form a planar triangle, which they do whenever
$\vect{u}$ and $\vect{v}$ are non-zero.
\end{exercise}

\begin{exercise}
Derive Lagrange's identity $\norm{\vect{u} \times \vect{v}}^{2} =
\norm{\vect{u}}^{2}\norm{\vect{v}}^{2} - (\vect{u} \cdot
\vect{v})^{2}$ by expanding the components of the cross product and
collecting terms. State the identity's interpretation as the
algebraic content of $\sin^{2}\theta + \cos^{2}\theta = 1$.

\paragraph{Solution sketch.} Write $\vect{u}\times\vect{v} = (u_{2}
v_{3} - u_{3} v_{2}, u_{3} v_{1} - u_{1} v_{3}, u_{1} v_{2} - u_{2}
v_{1})$ and expand each squared component, summing to get
$\sum_{i\neq j} u_{i}^{2} v_{j}^{2} - 2 (u_{1} v_{1} u_{2} v_{2} +
u_{1} v_{1} u_{3} v_{3} + u_{2} v_{2} u_{3} v_{3})$. The first sum
equals $\left(\sum_{i} u_{i}^{2}\right)\left(\sum_{j} v_{j}^{2}
\right) - \sum_{i} u_{i}^{2} v_{i}^{2}$, and the cross term equals
$\left(\sum_{i} u_{i} v_{i}\right)^{2} - \sum_{i} u_{i}^{2} v_{i}^{2}$
when expanded. Collecting gives $\norm{\vect{u}\times\vect{v}}^{2} =
\norm{\vect{u}}^{2}\norm{\vect{v}}^{2} - (\vect{u}\cdot\vect{v})^{2}$.
Interpretation: substituting $\vect{u}\cdot\vect{v} = \norm{\vect{u}}
\norm{\vect{v}}\cos\theta$ and $\norm{\vect{u}\times\vect{v}} =
\norm{\vect{u}}\norm{\vect{v}}\sin\theta$ on both sides reduces the
identity to $\sin^{2}\theta = 1 - \cos^{2}\theta$, the Pythagorean
identity. Lagrange's identity is the coordinate-free statement of
the same fact.
\end{exercise}

\begin{exercise}
Derive the BAC-CAB rule, $\vect{u} \times (\vect{v} \times \vect{w})
= \vect{v}(\vect{u} \cdot \vect{w}) - \vect{w}(\vect{u} \cdot
\vect{v})$, by expanding both sides in Cartesian components and
showing they are equal term by term.

\paragraph{Solution sketch.} Let $\vect{a} = \vect{v}\times\vect{w}$.
The $x$-component of $\vect{u}\times\vect{a}$ is $u_{2} a_{3} -
u_{3} a_{2} = u_{2}(v_{1} w_{2} - v_{2} w_{1}) - u_{3}(v_{3} w_{1}
- v_{1} w_{3})$. Add and subtract $u_{1} v_{1} w_{1}$:
$= v_{1}(u_{1} w_{1} + u_{2} w_{2} + u_{3} w_{3}) - w_{1}(u_{1}
v_{1} + u_{2} v_{2} + u_{3} v_{3}) = v_{1}(\vect{u}\cdot\vect{w}) -
w_{1}(\vect{u}\cdot\vect{v})$. This is the $x$-component of
$\vect{v}(\vect{u}\cdot\vect{w}) - \vect{w}(\vect{u}\cdot\vect{v})$.
The $y$- and $z$-components follow by cyclic permutation of indices.
The identity is the most useful three-vector identity in mechanics
and electromagnetism; it unwinds every triple cross product into
dot-product-weighted vectors that the engineer can manipulate
componentwise.
\end{exercise}

\begin{exercise}
Derive the divergence-theorem formula for the volume of a region,
$V = \tfrac{1}{3} \oint \vect{r} \cdot \unitvec{n} \,\dd A$, by
applying the divergence theorem to the position vector field
$\vect{F}(\vect{r}) = \vect{r}$. State the regularity assumption on
the boundary that the derivation requires.

\paragraph{Solution sketch.} Compute $\divg\vect{r} = \partial x/
\partial x + \partial y/\partial y + \partial z/\partial z = 3$.
Substitute into the divergence theorem: $\iiint_{V} 3\,\dd V =
\oint_{\partial V}\vect{r}\cdot\unitvec{n}\,\dd A$. The left side
is $3 V$, so $V = \tfrac{1}{3}\oint_{\partial V}\vect{r}\cdot
\unitvec{n}\,\dd A$. Regularity: the divergence theorem requires
$\partial V$ to be a piecewise-smooth, closed, orientable surface
with outward-pointing unit normal $\unitvec{n}$ defined almost
everywhere. The position field $\vect{r}$ is smooth everywhere in
$\R^{3}$, so the limiting assumption is on the boundary geometry.
Lipschitz boundaries (which include all piecewise-flat
triangulations) suffice.
\end{exercise}

\begin{exercise}
Derive the formula $V_{\triangle} = \tfrac{1}{6}\vect{p}_{1} \cdot
(\vect{p}_{2} \times \vect{p}_{3})$ for the contribution of a single
oriented triangle with vertices $\vect{p}_{1}, \vect{p}_{2},
\vect{p}_{3}$ to the divergence-theorem volume of a closed
triangulated surface. Start from the divergence-theorem expression
$V = \tfrac{1}{3} \oint \vect{r} \cdot \unitvec{n} \,\dd A$ and
specialise it to a flat triangle.

\paragraph{Solution sketch.} On a flat triangle, the area vector is
$\vect{A}_{\triangle} = \tfrac{1}{2}(\vect{p}_{2} - \vect{p}_{1})
\times(\vect{p}_{3} - \vect{p}_{1})$ with $\unitvec{n}\,\dd A =
\vect{A}_{\triangle}$ when summed over the triangle. The centroid
is $\vect{r}_{c} = (\vect{p}_{1} + \vect{p}_{2} + \vect{p}_{3})/3$.
For a flat triangle, $\int_{\triangle}\vect{r}\,\dd A = \vect{r}_{c}
\,A_{\triangle}$, so the contribution to $\oint\vect{r}\cdot
\unitvec{n}\,\dd A$ is $\vect{r}_{c}\cdot(\unitvec{n} A_{\triangle})
= \vect{r}_{c}\cdot\vect{A}_{\triangle}$. Expand:
$3 V_{\triangle} = \vect{r}_{c}\cdot\vect{A}_{\triangle} =
\tfrac{1}{6}(\vect{p}_{1} + \vect{p}_{2} + \vect{p}_{3})\cdot
((\vect{p}_{2} - \vect{p}_{1})\times(\vect{p}_{3} - \vect{p}_{1}))$.
Expanding the cross product and exploiting the antisymmetry of the
scalar triple product (terms with repeated vectors vanish), the
expression collapses to $\tfrac{1}{2}\vect{p}_{1}\cdot
(\vect{p}_{2}\times\vect{p}_{3})$. Dividing by 3 gives $V_{\triangle}
= \tfrac{1}{6}\vect{p}_{1}\cdot(\vect{p}_{2}\times\vect{p}_{3})$,
which is recognisable as the signed volume of the tetrahedron with
one vertex at the origin and opposite face the triangle
\cite{paper:v2ch08-zhang-chen-2001}.
\end{exercise}

\begin{exercise}
Derive the cylindrical-coordinate expression for $\divg\vect{F}$ by
applying the divergence theorem to a small cylindrical-shell volume
element of dimensions $\dd\rho \times \rho\,\dd\phi \times \dd z$
and computing the net flux through its six faces in the limit as the
volume shrinks to a point.

\paragraph{Solution sketch.} The shell has six bounding faces:
inner ($\rho$), outer ($\rho + \dd\rho$), bottom ($\phi$), top
($\phi + \dd\phi$), back ($z$), front ($z + \dd z$). The volume is
$\dd V = \rho\,\dd\rho\,\dd\phi\,\dd z$. Net radial flux: the outer
face has area $(\rho + \dd\rho)\,\dd\phi\,\dd z$ and flux density
$F_{\rho}(\rho + \dd\rho)$; the inner face has area $\rho\,\dd\phi
\,\dd z$ with flux $-F_{\rho}(\rho)$. Net radial contribution is
$[\partial(\rho F_{\rho})/\partial\rho]\,\dd\rho\,\dd\phi\,\dd z$
to leading order. Net azimuthal flux: top minus bottom face,
each of area $\dd\rho\,\dd z$, contributes $(\partial F_{\phi}/
\partial\phi)\,\dd\phi\,\dd\rho\,\dd z$. Net axial flux: top
minus bottom $z$-face, each of area $\rho\,\dd\rho\,\dd\phi$,
contributes $(\partial F_{z}/\partial z)\,\dd z\,\rho\,\dd\rho
\,\dd\phi$. Sum the three contributions, divide by $\dd V = \rho
\,\dd\rho\,\dd\phi\,\dd z$, and take the limit: $\divg\vect{F} =
\frac{1}{\rho}\partial(\rho F_{\rho})/\partial\rho + \frac{1}{\rho}
\partial F_{\phi}/\partial\phi + \partial F_{z}/\partial z$. The
factor $1/\rho$ in the first two terms is the Jacobian compensating
for the volume element's $\rho$ scaling.
\end{exercise}

\begin{exercise}
Show that the curl of any gradient vanishes: $\curl(\grad\Phi) =
\vect{0}$ for every twice-continuously-differentiable scalar field
$\Phi$. State the assumption on $\Phi$ that the derivation uses.

\paragraph{Solution sketch.} Compute the $x$-component of $\curl
(\grad\Phi)$: $\partial(\grad\Phi)_{z}/\partial y - \partial(\grad
\Phi)_{y}/\partial z = \partial^{2}\Phi/(\partial y\,\partial z) -
\partial^{2}\Phi/(\partial z\,\partial y)$. By Clairaut's theorem on
the equality of mixed partial derivatives, the two are equal
whenever $\Phi \in C^{2}$ (twice continuously differentiable), so
the $x$-component vanishes. The $y$- and $z$-components follow by
the same argument with cyclic permutation. The required assumption
is therefore $\Phi \in C^{2}$; without it the mixed partials need
not commute and the result can fail at points where $\Phi$ has weak
singularities.
\end{exercise}

\begin{exercise}
Show that the divergence of any curl vanishes: $\divg(\curl\vect{F})
= 0$ for every twice-continuously-differentiable vector field
$\vect{F}$. State the assumption on $\vect{F}$ that the derivation
uses.

\paragraph{Solution sketch.} Expand:
$\divg(\curl\vect{F}) = \partial/\partial x(\partial F_{z}/\partial y
- \partial F_{y}/\partial z) + \partial/\partial y(\partial F_{x}/
\partial z - \partial F_{z}/\partial x) + \partial/\partial z(\partial
F_{y}/\partial x - \partial F_{x}/\partial y)$. Each mixed second
partial appears twice with opposite signs (e.g., $\partial^{2}
F_{z}/(\partial x\,\partial y)$ and $-\partial^{2} F_{z}/(\partial y
\,\partial x)$), so all six terms cancel in pairs by Clairaut's
theorem. Required assumption: $\vect{F} \in C^{2}$. The identity
underwrites the conservation laws of incompressible vector
potentials (the magnetic vector potential in electromagnetism, the
stream function in two-dimensional fluid flow) by guaranteeing that
the curl of such a potential is divergence-free.
\end{exercise}

\subsection*{Estimation}\pathtag{core}

\begin{exercise}
Estimate the surface area of a household object of the reader's
choice (a coffee mug, an apple, a small stone) using a vertex-and-
triangle approximation with no more than 12 sampled points. State
the dominant uncertainty source and a factor-of-three confidence
range. Compare to a more careful measurement if one is available.

\paragraph{Solution sketch.} For an apple modelled as a sphere of
radius $4\,\si{\centi\meter}$, the true surface area is $4\pi r^{2}
\approx 0.020\,\si{\meter\squared}$. A 12-point triangulation gives
6 triangles in the upper hemisphere and 6 in the lower; each
triangle has area on the order of $r^{2}\sin(60^{\circ}) \approx
1.4\times 10^{-3}\,\si{\meter\squared}$, totalling roughly $0.017\,
\si{\meter\squared}$. The dominant uncertainty is the
piecewise-flat approximation, which underestimates the area of a
convex object by up to 20\% at this sampling density. A
factor-of-three range, $0.007$ to $0.06\,\si{\meter\squared}$,
covers the variability across object shape (oranges curve more
sharply than cucumbers) and across sampling densities. A careful
measurement (wrap with paper, trace, weigh) typically lands within
the upper half of the range.
\end{exercise}

\begin{exercise}
Estimate the volume of a household object using the divergence-
theorem formula $V = \tfrac{1}{3}\oint \vect{r}\cdot\unitvec{n}
\,\dd A$ on a coarse 8-point sampled approximation. Compare to the
estimate one would obtain from a single overall length-times-width-
times-height bounding box, and identify which estimate is more
defensible and why.

\paragraph{Solution sketch.} An 8-point sample (4 vertices on a
mid-equator and 2 each on top and bottom apices) gives 8 outward
triangles. Summing the signed-tetrahedron contributions yields an
estimate within 30\% of the true volume for a smooth blob, and the
estimate is systematically low because flat triangles cut chords
inside any convex region. The bounding-box estimate is
systematically high by the volume fraction of the object inside its
box: roughly a factor of 2 high for a sphere ($\pi/6 \approx 0.52$
filling fraction), a factor of 1.5 high for a typical rounded
household object. The 8-point divergence-theorem estimate is more
defensible because its error has a known sign (negative) and a
known mechanism (chord-cutting), which the engineer can correct
with a refinement step; the bounding-box estimate's error depends
on the object's shape relative to the box and cannot be corrected
without measuring the object.
\end{exercise}

\begin{exercise}
Estimate the magnitude of the gradient of the temperature field
inside a household oven during the first 60 seconds after the door
is closed, in $\,\si{\kelvin\per\meter}$. State the assumptions
about the heat source location, the oven's interior dimensions, and
the temperature difference between the heater and the air.

\paragraph{Solution sketch.} Assume the oven cavity is roughly
$0.4 \times 0.4 \times 0.4\,\si{\meter}$, the heating element is a
glowing coil at $700\,\degC$ at the top of the cavity, and the
ambient air starts at $20\,\degC$. The temperature drop across the
$0.4\,\si{\meter}$ vertical extent is $\Delta T \approx 680\,
\si{\kelvin}$ in the first few seconds, before convection equalises
the cavity. The gradient magnitude is then $\norm{\grad T}\approx
\Delta T/\ell \approx 680/0.4 \approx 1.7\times 10^{3}\,\si{\kelvin
\per\meter}$, directed upward. By 60 seconds the convection cell
has driven $\Delta T$ across the cavity down to perhaps $100\,
\si{\kelvin}$, giving $\norm{\grad T}\approx 250\,\si{\kelvin\per
\meter}$. A factor-of-three band brackets the strong unsteadiness
of the period; the dominant uncertainty is the convection-onset
time, which depends on cavity shape, element wattage, and door
seal quality.
\end{exercise}

\begin{exercise}
Estimate the divergence of the velocity field of room air during
the first few seconds after a household window is opened on a
breezy day. State whether the divergence is positive (air entering
the room is a volume source), negative (air leaving the room is a
volume sink), or near zero (incompressible-flow approximation), and
state the regime in which each answer is defensible.

\paragraph{Solution sketch.} For room air at habitable temperatures
and Mach numbers well below $0.1$, the incompressible-flow
approximation gives $\divg\vect{v} \approx 0$ at every interior
point. The room's air mass is approximately constant while the
window is open: as air flows in through one part of the window, the
same volumetric rate flows out elsewhere (through the window itself,
through other openings, or by displacing existing room air). A
local pocket can have transient positive or negative divergence on
the order of $\norm{\vect{v}}/\ell \approx 1\,\si{\meter\per\second}
/0.5\,\si{\meter} = 2\,\si{\per\second}$ in the first second after
opening, then decays as the room equilibrates. The
incompressible-flow answer is defensible whenever the Mach number
is small and the temperature is roughly uniform; the
positive-source answer applies in the explicit accounting where the
``room'' is the control volume and the window is a boundary
through which air enters; the negative-sink answer applies on the
other side of the window.
\end{exercise}

\subsection*{Simulation}\pathtag{mastery}

\begin{exercise}
Implement, in a programming environment of choice, the
divergence-theorem volume formula
$V = \tfrac{1}{6}\sum_{\triangle} \vect{p}_{1}\cdot(\vect{p}_{2}
\times \vect{p}_{3})$ for a triangulated closed surface. Test the
implementation against a known-volume shape (a unit cube
triangulated into 12 triangles, a regular tetrahedron, a sphere
triangulated into a geodesic mesh of at least 80 triangles). Report
the relative error of each test and identify the regime in which
the truncation error of the piecewise-flat approximation
dominates the position-sampling uncertainty.

\paragraph{Solution sketch.} A reference implementation is in
\repopath{volumes/02-form/ch08-vectors-and-vector-calculus/code/mesh_volume.py}.
The companion \texttt{triple\_product\_sweep.py} runs the three test
cases. The cube and tetrahedron return the analytical volumes to
floating-point precision; the geodesic icosphere at refinement zero
(20 triangles) has $V_{\text{DT}} \approx 2.54$ against the true
$4\pi/3 \approx 4.19$, a relative error of $0.40$; at refinement
two (320 triangles) the error drops to about $3 \times 10^{-2}$;
at refinement four (5120 triangles) the error drops to about
$2 \times 10^{-3}$. The reference numbers are tabulated in
\repopath{data/sweep_reference.csv}. The truncation error scales as
edge-length-squared (the standard piecewise-flat order), so it
dominates the position-sampling uncertainty whenever the edge
length exceeds the position uncertainty by a factor of $\sqrt{N}$
where $N$ is the number of triangles. For household measurements
with millimetre-scale position uncertainty and centimetre-scale
edges, truncation dominates until the mesh exceeds approximately
1000 triangles.
\end{exercise}

\begin{exercise}
Implement a Monte Carlo estimate of the surface area of a
parametrised surface $\vect{r}(u, v) = (u\cos v, u\sin v, u^{2})$,
$u \in [0, 1]$, $v \in [0, 2\pi]$. Use at least $10\,000$
uniformly distributed sample points in $(u, v)$, evaluate the
local area density $\norm{\vect{r}_{u} \times \vect{r}_{v}}$ at
each, and report the mean and standard deviation of the area
estimate. Compare against the analytical integral.

\paragraph{Solution sketch.} The analytical area is $A = \pi(5
\sqrt{5} - 1)/6 \approx 5.33$ (computed in exercise 8). The Monte
Carlo estimator is $A_{\text{MC}} = \abs{D}\cdot\langle\norm{
\vect{r}_{u}\times\vect{r}_{v}}\rangle$ where $\abs{D} = 2\pi$ is the
area of the parameter domain and the angle brackets denote the
empirical mean over the $N$ samples. With $N = 10^{4}$ uniform
samples in $(u, v) \in [0, 1]\times[0, 2\pi]$, the estimator
converges to $5.33 \pm 0.02$ (sample standard error). The relative
error scales as $1/\sqrt{N}$; achieving $10^{-3}$ relative error
requires roughly $10^{6}$ samples. Monte Carlo is competitive with
deterministic quadrature on smooth two-dimensional integrands only
at low accuracy targets; the divergence-theorem formula uses
deterministic quadrature on the triangulated surface and converges
at second order, faster than Monte Carlo for the same point budget.
\end{exercise}

\subsection*{Design}\pathtag{standard}

\begin{exercise}
Design a sampling strategy for measuring the volume of an irregular
metal casting by the divergence theorem to within $1\,\%$ relative
uncertainty. Specify the number of surface sample points, the
expected position uncertainty per point, and the truncation
uncertainty of the piecewise-flat surface approximation. State
which uncertainty source dominates at the chosen sampling density
and which would dominate at twice or half the density.

\paragraph{Discussion.} For a casting of characteristic size
$0.3\,\si{\meter}$ and volume $\sim 10^{-2}\,\si{\meter\cubed}$, a
$1\,\%$ relative volume uncertainty corresponds to an absolute
budget of $10^{-4}\,\si{\meter\cubed}$. A coordinate-measuring
machine with $0.01\,\si{\milli\meter}$ point uncertainty supports
position-budget contributions far below $1\,\%$ even with sparse
sampling. Truncation, scaling as edge-length-squared, dominates:
the casting needs an edge length below approximately $5\,\si{\milli
\meter}$, equivalent to $\sim 10^{4}$ triangles distributed
preferentially on high-curvature features. At half the sampling
density ($\sim 5\times 10^{3}$ triangles), truncation rises to
about $4\,\%$ and dominates by a wide margin. At twice the density
($\sim 2 \times 10^{4}$ triangles), truncation falls to $0.25\,\%$
and the budget tightens to a different regime: registration error
between scan patches and edge-detection error in the mesh become
the dominant sources. The design choice is roughly $10^{4}$
triangles with CMM-grade position accuracy and a curvature-adaptive
mesh.
\end{exercise}

\begin{exercise}
Design a procedure for verifying the orientation convention of a
triangulated closed surface mesh produced by an unspecified
software tool. State the operational test, the expected output for
a correctly oriented mesh, and the corrective action when the test
indicates inverted normals on some triangles.

\paragraph{Discussion.} The operational test is the divergence-
theorem signed volume itself. Compute $V_{\text{DT}} = \tfrac{1}{6}
\sum_{\triangle}\vect{p}_{1}\cdot(\vect{p}_{2}\times\vect{p}_{3})$.
A correctly oriented mesh returns a positive value; a uniformly
inverted mesh returns the negative of the same magnitude; a mixed
mesh returns a value with magnitude smaller than the bounding-box
volume and possibly of indeterminate sign. The expected output for
a unit cube is exactly $+1.000$. Corrective action depends on the
diagnosis. Uniform inversion: reverse every triangle's vertex order
(swap $\vect{p}_{2}$ and $\vect{p}_{3}$ in each triple). Mixed
orientation: run a flood-fill orientation propagation starting from
one triangle known to be outward (pick the triangle whose centroid
has the largest $\norm{\vect{r}}$; its outward normal is the one
that points away from the centroid of the bounding box) and
propagate the convention to adjacent triangles via shared-edge
matching. Re-run the volume test after the correction; the
absolute volume should remain unchanged but the sign should now be
positive.
\end{exercise}

\begin{exercise}
Design a coordinate frame for a fluid-flow measurement in a
circular pipe of inner radius $50 \,\si{\milli\meter}$. State which
of Cartesian, cylindrical, or spherical coordinates is the working
choice; identify the symmetry that selects the choice; and write
the divergence operator $\divg\vect{v}$ for the velocity field in
the chosen coordinates.

\paragraph{Discussion.} Cylindrical coordinates $(\rho, \phi, z)$
are the working choice. The pipe has axial symmetry: the geometry is
invariant under rotations about the pipe axis. Choose $\unitvec{z}$
along the pipe axis, $\rho$ the radial distance from the axis, and
$\phi$ the azimuthal angle. The velocity field in fully developed
laminar pipe flow is $\vect{v} = v_{z}(\rho)\unitvec{z}$ with no
radial or azimuthal component. The divergence is
\[
\divg\vect{v} = \frac{1}{\rho}\pdv{(\rho v_{\rho})}{\rho} +
\frac{1}{\rho}\pdv{v_{\phi}}{\phi} + \pdv{v_{z}}{z}.
\]
For the fully developed flow above, every term vanishes (the field
is incompressible and uniform along $z$ within the developed
region), consistent with $\divg\vect{v} = 0$. Cartesian coordinates
would force the engineer to carry sines and cosines in the
boundary conditions at every $\rho = 50\,\si{\milli\meter}$
checkpoint; spherical coordinates have a symmetry mismatch (the
pipe has cylindrical, not spherical, symmetry).
\end{exercise}

\begin{exercise}
Design a procedure for computing the area of a planar parcel of
land from a surveyed sequence of $(x, y)$ corner coordinates using
the shoelace formula. State the surveying tolerance required for
the area to be known to within $1 \,\si{\meter\squared}$ on a
parcel of nominal area $1000 \,\si{\meter\squared}$.

\paragraph{Discussion.} Walk the boundary counterclockwise, recording
$(x_{i}, y_{i})$ at each corner. Compute the area as
$A = \tfrac{1}{2}\abs{\sum_{i = 1}^{n}(x_{i} y_{i+1} - x_{i+1}
y_{i})}$ with $x_{n+1} = x_{1}$, $y_{n+1} = y_{1}$. The relative
uncertainty on $A$ from $n$ corners with independent position
uncertainty $\sigma_{p}$ each is approximately $\sigma_{A}/A \approx
2\sigma_{p}/\sqrt{P\cdot A/n}$ where $P$ is the parcel perimeter.
For a square parcel with $A = 1000\,\si{\meter\squared}$, side
length $\ell \approx 31.6\,\si{\meter}$, $P \approx 126\,\si{\meter}$,
$n = 4$ corners. Solving for $\sigma_{p}$ at $\sigma_{A}/A = 10^{-3}$
gives $\sigma_{p} \approx 0.013\,\si{\meter}$, or about $1.3\,
\si{\centi\meter}$ per corner. A handheld GPS unit (uncertainty
$\sim 1\,\si{\meter}$) fails the tolerance by two orders of
magnitude; a survey-grade RTK GPS (uncertainty $\sim 0.01\,\si{
\meter}$, current as of 2026) marginally meets it; a total-station
survey (uncertainty $\sim 1\,\si{\milli\meter}$) comfortably
exceeds it. The shoelace formula is exact; the surveying tolerance
governs the result.
\end{exercise}

\subsection*{Diagnosis and reverse engineering}\pathtag{standard}

\begin{exercise}
A divergence-theorem volume calculation on a triangulated
household-object surface returns a value $V_{\text{DT}}$ that is
$15\,\%$ smaller than the water-displacement volume $V_{\text{WD}}$
measured for the same object. State three plausible mechanisms (the
sampling is too coarse, one or more triangles has an inverted
normal, the water-displacement reading included surface bubbles)
and propose, for each, the test that would confirm or refute it.

\paragraph{Solution sketch.} (1) Coarse sampling: refine the mesh by
adding sample points on high-curvature features and recompute. If
$V_{\text{DT}}$ rises monotonically toward $V_{\text{WD}}$ with
refinement, truncation was the cause. (2) Inverted normals: run the
sign-and-magnitude test against the bounding box. If $V_{\text{DT}}$
is positive but several triangles have outward-pointing normals
that cross paths (visible in a renderer or by checking
$\vect{p}_{i}\cdot\unitvec{n}_{i}$ values that vary in sign for
neighbouring triangles on a convex region), patch the orientation
with a flood-fill from a known-outward triangle. (3) Surface
bubbles: redo the water-displacement measurement with the object
submerged and gently tapped or rotated to dislodge bubbles. If
$V_{\text{WD}}$ drops to within $15\,\%$ of $V_{\text{DT}}$, bubbles
were the cause. The diagnoses are testable independently; the
engineer runs them in order of cost (cheapest first: re-read the
water level after tapping; then check orientation; then refine the
mesh).
\end{exercise}

\begin{exercise}
A finite-element heat-conduction code reports a positive divergence
of the heat flux at every interior point of a steel block held at
constant temperature. State whether the report is physically
plausible and identify the most likely software-side mechanism if it
is not.

\paragraph{Solution sketch.} The report is not physically plausible.
At steady state, $\divg\vect{q} = 0$ at every interior point of a
source-free conductor; a positive divergence everywhere would
require a uniform volumetric heat source the problem statement does
not include, and energy conservation would force the temperature to
rise (contradicting the constant-temperature constraint). Most
likely software-side mechanism: a sign error in the constitutive
relation, with $\vect{q} = +k\grad T$ entered instead of $\vect{q}
= -k\grad T$. The substitution flips the sign of every flux vector
and changes the boundary balance from outward-to-inward, which a
naive divergence calculation reports as a positive interior
source. A secondary suspect is a unit error in $k$ that crosses a
sign change (rare but seen when imperial conductivity tables are
mixed with SI grids). The diagnostic is to print $\vect{q}$ at a
known boundary node and check whether the vector points from hot to
cold (correct) or cold to hot (sign error).
\end{exercise}

\begin{exercise}
A torque computation in an inertial-navigation simulation returns a
torque vector pointing along $-\unitvec{z}$ when the corresponding
physical experiment shows the platform rotating about $+\unitvec{z}$.
State the most likely cause and the operational test that would
distinguish a hand-rule swap from a passive-versus-active
transformation error.

\paragraph{Solution sketch.} The two leading candidates are (a) a
left-handed coordinate frame applied with the right-hand rule (or
vice versa) at the cross product, and (b) a passive-active confusion
in the rotation matrix that converts body-frame torque to
inertial-frame torque. The operational test: compute the scalar
triple product $\unitvec{x}\cdot(\unitvec{y}\times\unitvec{z})$ in
the simulation's frame. If the result is $+1$, the frame is
right-handed (hand-rule swap is not the cause; suspect (b)). If the
result is $-1$, the frame is left-handed (hand-rule swap is the
cause; either change the rule or swap two basis vectors to make the
frame right-handed). The passive-active error is then distinguished
by feeding a known input rotation (say, $+10^{\circ}$ about
$+\unitvec{z}$) and checking whether the simulation reports the body
spinning at $+10^{\circ}$ (active correct) or the frame spinning at
$-10^{\circ}$ as observed from the rotated body (passive
misinterpreted as active).
\end{exercise}

\subsection*{Failure analysis}\pathtag{core}

\begin{exercise}
Construct a worked example, parallel to the cube example in
section~8.8, in which a piecewise parametrisation of a closed
surface produces inverted normals on at least one face, and in
which the resulting flux integral differs from the divergence-
theorem volume integral by a known amount. State the operational
test that would catch the error before the calculation is
completed.

\paragraph{Solution sketch.} Consider a cylindrical can: side
surface $\rho = 1$ for $z \in [0, 1]$, top cap $z = 1$, bottom cap
$z = 0$. Parametrise the side as $\vect{r}(\phi, z) = (\cos\phi,
\sin\phi, z)$ for $\phi \in [0, 2\pi]$, $z \in [0, 1]$; the induced
normal $\vect{r}_{\phi}\times\vect{r}_{z} = (\cos\phi, \sin\phi,
0) = +\unitvec{\rho}$ points outward as required. Parametrise both
caps using the same Cartesian rule $\vect{r}(x, y) = (x, y, z_{0})$
on the unit disc. The top cap's induced normal is $+\unitvec{z}$
(outward, correct); the bottom cap's induced normal is also
$+\unitvec{z}$ (inward, wrong). For the test field $\vect{F} = (0,
0, 1)$, the flux through the top is $+\pi$, through the bottom in
the wrong orientation is $+\pi$ (added to the total instead of
subtracted); side contributions are zero. The reader's result is
$+2\pi$ instead of the correct $0$ (divergence is zero). The
operational test: at parametrisation time, evaluate the induced
normal on a sample point of each cap and check by inspection
whether it points away from the can's interior; the bottom cap fails
the test before any integration begins.
\end{exercise}

\begin{exercise}
A team computes the volume of a triangulated closed surface mesh
using $V = \tfrac{1}{6}\sum_{\triangle} \vect{p}_{1} \cdot
(\vect{p}_{2} \times \vect{p}_{3})$ and obtains a value smaller in
magnitude than the bounding-box volume by a factor of two and of
the wrong sign. State the most likely failure (an inverted
orientation convention applied uniformly to every triangle, so the
sum reports $-V$ instead of $+V$) and explain why the magnitude
agrees with the bounding-box estimate scaled by the
volume-fraction the bounding-box overestimates.

\paragraph{Solution sketch.} The most likely failure is exactly what
the prompt names: a uniformly inverted orientation. The
divergence-theorem sum is signed, so swapping every triangle's
winding order flips the sign of every contribution; the sum becomes
$-V$. The magnitude of $\abs{-V}$ is still $V$, the correct volume,
which for a typical rounded object fills approximately half the
bounding box (factor $\sim 0.5$). The team's result therefore
matches $-(\text{bounding box}/2)$ in magnitude and sign, which is
the observed disagreement. The remedy is to multiply every
triangle's vertex order by a $(\vect{p}_{1}, \vect{p}_{2},
\vect{p}_{3}) \to (\vect{p}_{1}, \vect{p}_{3}, \vect{p}_{2})$ swap
and re-run the computation; the result should be the absolute value
of the original sum with positive sign. The reference code in
\repopath{code/mesh_volume.py} prints a warning when the sum is
negative, which catches the failure at the moment the script runs
rather than at the moment the team presents the result.
\end{exercise}

\begin{exercise}
A student computes the curl of a velocity field for an
incompressible fluid flow using a left-handed coordinate frame and
the right-hand rule, and is surprised when the predicted vorticity
direction disagrees with a smoke-visualisation experiment. State
the corrective action and identify the engineering domain (computer
graphics, robotics, geodesy) in which left-handed frames are common
enough that the cross-domain reader should check the convention
before trusting any cross product.

\paragraph{Solution sketch.} The corrective action is to apply the
left-hand rule in a left-handed frame (and the right-hand rule in a
right-handed frame), or, equivalently, to swap two basis vectors to
convert the left-handed frame into a right-handed one and continue
with the right-hand rule. The vorticity direction will then agree
with the smoke experiment. Engineering domains where left-handed
frames are standard: computer graphics (Direct3D, screen
coordinates with $+y$ down and $+z$ into the screen); some aerospace
and naval engineering body-fixed frames ($+x$ forward, $+y$
starboard, $+z$ down). The reader crossing from fluid mechanics to
graphics-pipeline integration or to flight-control software checks
the frame's handedness as the first step before any cross product
or curl computation.
\end{exercise}

\subsection*{Judgment}\pathtag{standard}

\begin{exercise}
Argue, in 200 to 300 words, whether the divergence theorem or the
empirical water-displacement method is the more defensible volume
measurement for an irregular household object. Take a position;
defend it with a specific scenario in which your preferred answer
fails and the other method is preferred.

\paragraph{Discussion.} A defensible answer for the household
setting favours water displacement: it is direct, requires no
orientation discipline, has a single dominant uncertainty source
(meniscus reading) the reader can quantify, and produces the volume
in seconds. The divergence theorem requires a triangulation, an
orientation check, and a sum over many signed contributions, each
introducing its own error mode. The scenario in which water
displacement fails: a porous or water-soluble object (sponge,
cardboard, bath salt), where the wetting changes the geometry being
measured; or an object too large for available beakers. In those
cases the divergence-theorem method on a 3D-scanned mesh is
preferred. The deeper point is that the methods are complementary:
the divergence theorem scales to objects of any size and any
material once the surface is sampled, while water displacement
scales to whatever fits in a beaker. The engineering choice is
governed by which uncertainty source the engineer prefers to
defend.
\end{exercise}

\begin{exercise}
A senior colleague tells the reader that ``vector calculus is for
physicists; in industry we use software, and the software handles
the orientations.'' Write a one-paragraph reply that takes the
colleague seriously, identifies the regime in which the claim is
defensible, and identifies the regime in which it is not.

\paragraph{Discussion.} The colleague is right in the regime of
mature, single-vendor, well-validated software applied to standard
problems: a commercial CFD code computing flow over a wing at
design-point Reynolds number does carry an orientation convention
in its mesh format and applies it consistently. The claim fails
when (a) the engineer integrates two pieces of software whose
orientation conventions differ (a body-fixed flight-control frame
imported into a graphics renderer; a CAD mesh imported into a
finite-element solver with inverted face normals); (b) the engineer
runs the software outside its validated regime, where the user
becomes responsible for verifying that the result makes physical
sense; or (c) the result is anomalous and the engineer needs to
diagnose whether the cause is the software or the input. In each
case the engineer who cannot reason about the orientation by hand
cannot diagnose the failure. The vector-calculus literacy is
insurance against the day the software does not handle the
orientation, which is the day it matters most.
\end{exercise}

\begin{exercise}
A working engineer computes a quantity by both a vector-calculus
derivation and a direct numerical simulation, and the two disagree
by $5\,\%$. Argue, in one paragraph, which of the following actions
is more defensible: trust the analytical derivation and search for
a software bug; trust the simulation and search for a missed term in
the derivation; treat both as candidate answers and seek a third
independent method to break the tie. Reference the cross-checking
discipline of Volume~I Chapter~7 in defending the position.

\paragraph{Discussion.} The third action is the most defensible.
Volume~I Chapter~7's cross-checking discipline insists that a
single calculation, no matter how carefully done, cannot defeat a
second independent measurement; only a third independent estimate
can break a tie, and even then the three should agree within their
combined uncertainties before the engineer accepts the result.
Trusting the analytical derivation assumes the derivation has no
missed term, no algebra slip, no boundary-condition oversight;
trusting the simulation assumes no software bug, no mesh artefact,
no input-data error. Both assumptions are routinely violated in
practice at the $5\,\%$ level. A third method (an experimental
measurement, a published reference value, a back-of-envelope
estimate, a different numerical scheme) raises the engineer's
defensible confidence from ``my derivation says X'' to ``three
independent methods agree on X within uncertainty.'' The discipline
is more expensive in time and is the difference between a
defensible answer and a result that survives until it does not.
\end{exercise}
